% Options for packages loaded elsewhere
\PassOptionsToPackage{unicode}{hyperref}
\PassOptionsToPackage{hyphens}{url}
\PassOptionsToPackage{dvipsnames,svgnames,x11names}{xcolor}
\documentclass[
  10pt,
  letterpaper,
]{article}
\usepackage{xcolor}
\usepackage[margin=0.60in]{geometry}
\usepackage{amsmath,amssymb}
\setcounter{secnumdepth}{-\maxdimen} % remove section numbering
\usepackage{iftex}
\ifPDFTeX
  \usepackage[T1]{fontenc}
  \usepackage[utf8]{inputenc}
  \usepackage{textcomp} % provide euro and other symbols
\else % if luatex or xetex
  \usepackage{unicode-math} % this also loads fontspec
  \defaultfontfeatures{Scale=MatchLowercase}
  \defaultfontfeatures[\rmfamily]{Ligatures=TeX,Scale=1}
\fi
\usepackage{lmodern}
\ifPDFTeX\else
  % xetex/luatex font selection
\fi
% Use upquote if available, for straight quotes in verbatim environments
\IfFileExists{upquote.sty}{\usepackage{upquote}}{}
\IfFileExists{microtype.sty}{% use microtype if available
  \usepackage[]{microtype}
  \UseMicrotypeSet[protrusion]{basicmath} % disable protrusion for tt fonts
}{}
\makeatletter
\@ifundefined{KOMAClassName}{% if non-KOMA class
  \IfFileExists{parskip.sty}{%
    \usepackage{parskip}
  }{% else
    \setlength{\parindent}{0pt}
    \setlength{\parskip}{6pt plus 2pt minus 1pt}}
}{% if KOMA class
  \KOMAoptions{parskip=half}}
\makeatother
\usepackage{listings}
\newcommand{\passthrough}[1]{#1}
\lstset{defaultdialect=[5.3]Lua}
\lstset{defaultdialect=[x86masm]Assembler}
\usepackage{longtable,booktabs,array}
\usepackage{caption}
\captionsetup[table]{skip=6pt}
\newcounter{none} % for unnumbered tables
\usepackage{calc} % for calculating minipage widths
% Correct order of tables after \paragraph or \subparagraph
\usepackage{etoolbox}
\makeatletter
\patchcmd\longtable{\par}{\if@noskipsec\mbox{}\fi\par}{}{}
\makeatother
% Allow footnotes in longtable head/foot
\IfFileExists{footnotehyper.sty}{\usepackage{footnotehyper}}{\usepackage{footnote}}
\makesavenoteenv{longtable}
\setlength{\emergencystretch}{3em} % prevent overfull lines
\providecommand{\tightlist}{%
  \setlength{\itemsep}{0pt}\setlength{\parskip}{0pt}}
\setlength{\emergencystretch}{3em}
\AtBeginDocument{%
  \lstset{%
    basicstyle=\ttfamily\footnotesize,%
    breaklines=true,%
    breakatwhitespace=false,%
    columns=fullflexible,%
    keepspaces=true,%
    showstringspaces=false%
  }%
}
\usepackage{bookmark}
\IfFileExists{xurl.sty}{\usepackage{xurl}}{} % add URL line breaks if available
\urlstyle{same}
% fallback for those not using the hyperref driver hyperxmp:
\makeatletter
\@ifundefined{xmpquote}{\newcommand{\xmpquote}[1]{#1}}{}
\makeatother
\hypersetup{
  colorlinks=true,
  linkcolor={blue},
  filecolor={Maroon},
  citecolor={Blue},
  urlcolor={blue},
  pdfcreator={LaTeX via pandoc}}

\author{}
\date{}

\begin{document}

\section{First Gear Bank v1 Foundation Implementation
Plan}\label{first-gear-bank-v1-foundation-implementation-plan}

\begin{quote}
\textbf{For agentic workers:} REQUIRED SUB-SKILL: Use
superpowers:subagent-driven-development (recommended) or
superpowers:executing-plans to implement this plan task-by-task. Steps
use checkbox (\passthrough{\lstinline!- [ ]!}) syntax for tracking.
\end{quote}

\textbf{Goal:} Build the verified, test-driven foundation for First Gear
Bank: a safe package pipeline, engine-independent banking core,
fixed-point domain model, balanced journal, atomic coordinator,
sectioned persistence, and a thin Vintage Story 1.22.7 world-save host.

\textbf{Architecture:} Keep all financial rules and deterministic state
transitions in \passthrough{\lstinline!FirstGearBank.Core!}, which
references no Vintage Story assembly. The existing
\passthrough{\lstinline!First Gear Bank!} project remains the packaged
engine adapter and translates verified 1.22.7 lifecycle calls into
narrow core ports. State is published as immutable revisions and saved
as independently checksummed sections so corruption can isolate a
subsystem without inventing or silently resetting money.

\textbf{Tech Stack:} .NET SDK 10.0.400, C\# 14, Vintage Story API
1.22.7, xUnit.net v3 4.0.0 with Microsoft Testing Platform v2, Cake
Frosting 6.2.0, \passthrough{\lstinline!System.Text.Json!},
\passthrough{\lstinline!System.IO.Compression!}, and
SHA-256/HMAC-SHA-256 from the base class library.

\textbf{Spec:}
\texttt{docs/}\allowbreak{}\texttt{superpowers/}\allowbreak{}\texttt{specs/}\allowbreak{}\texttt{2026-}\allowbreak{}\texttt{09-}\allowbreak{}\texttt{08-}\allowbreak{}\texttt{first-}\allowbreak{}\texttt{gear-}\allowbreak{}\texttt{bank-}\allowbreak{}\texttt{v1-}\allowbreak{}\texttt{design.}\allowbreak{}\texttt{md}

\subsection{Global Constraints}\label{global-constraints}

\begin{itemize}
\tightlist
\item
  Target Vintage Story exactly 1.22.7 and use mod ID exactly
  \passthrough{\lstinline!firstgearbank!}.
\item
  If execution uses a Git worktree, create it as a direct sibling of
  this repository under the \passthrough{\lstinline!mods!} directory so
  the required default reference path
  \texttt{.}\allowbreak{}\texttt{.}\allowbreak{}\texttt{/}\allowbreak{}\texttt{Vintage\ }\allowbreak{}\texttt{Story\ }\allowbreak{}\texttt{Reference}
  remains valid.
\item
  Ship one mod package containing a pure .NET 10 core plus a thin
  Vintage Story adapter.
\item
  Keep all financial authority on the server; the client never supplies
  a trusted UID, balance, timestamp, rate, or result.
\item
  Use six-decimal signed fixed-point money: one gear is
  \passthrough{\lstinline!1,000,000!} bank units.
\item
  Round every floating-rate result to the nearest bank unit with
  midpoint ties to even.
\item
  Cap customer rusty-savings and temporal-vault balances at
  \passthrough{\lstinline!2,147,483,647,000,000!} units.
\item
  Never allow negative savings or temporal-vault balances; signed system
  and future contract accounts remain valid.
\item
  Treat the append-only, independently balanced journal as the monetary
  source of truth.
\item
  Publish each command as one validated immutable revision containing
  zero or more complete journal records.
\item
  Never treat default record equality over
  \passthrough{\lstinline!ImmutableArray!},
  \passthrough{\lstinline!ImmutableDictionary!}, or
  \passthrough{\lstinline!ReadOnlyMemory!} as semantic equality; use
  explicit structural comparison, canonical bytes, or a domain
  validator.
\item
  Preserve corrupt authoritative bytes and quarantine the affected
  domain; never synthesize a repair to money.
\item
  Never display or normally log Vintage Story PlayerUIDs.
\item
  Add no Harmony patch unless a public-API gap is documented and the
  user separately approves the exact patch.
\item
  Keep Better Loot, Better Loot Plus, Config Lib, and Integrated Mod
  Manager optional.
\item
  Do not implement any section 2.1 non-goal in this phase.
\item
  Follow all checked-in common, C\#, hygiene, and Vintage Story
  standards, including file headers, callable comments, exact
  three-blank-line spacing, 120-character documentation lines, and
  \passthrough{\lstinline!INFO!}/\passthrough{\lstinline!DEBG!}/\passthrough{\lstinline!WARN!}/\passthrough{\lstinline!CRIT!}
  log labels.
\item
  Use test-driven development for every behavior change: observe the
  intended failure, make the smallest passing change, then refactor
  while green.
\item
  Run formatter verification, build, target-project tests, the full test
  suite, and package verification before a task is complete.
\end{itemize}

\begin{center}\rule{0.5\linewidth}{0.5pt}\end{center}

\subsection{Phase Boundary}\label{phase-boundary}

This plan intentionally stops before playable deposits, interest, CDs,
transfers, GUI, or branches. Its output is a loadable foundation-only
mod whose state and package contracts are safe enough for the remaining
plans in
\texttt{2026-}\allowbreak{}\texttt{09-}\allowbreak{}\texttt{08-}\allowbreak{}\texttt{first-}\allowbreak{}\texttt{gear-}\allowbreak{}\texttt{bank-}\allowbreak{}\texttt{v1-}\allowbreak{}\texttt{roadmap.}\allowbreak{}\texttt{md}.

Two proof gates occur here because they affect persisted contracts:

\begin{enumerate}
\def\labelenumi{\arabic{enumi}.}
\tightlist
\item
  The physical-settlement spike records whether player mod data is
  provably co-serialized with inventory. If it is not, the supported
  recovery policy is authenticated evidence where available and
  quarantine for every ambiguous state.
\item
  The accrued-liability-index spike must prove a sublinear exact result
  for the specification\textquotesingle s per-account, month-by-month
  ties-to-even accrual. A scalar aggregate is known to be insufficient.
  If the proof fails, stop before the CD/liquidity projection schema is
  frozen and present the user with the measured conflict.
\end{enumerate}

\subsection{File and Responsibility
Map}\label{file-and-responsibility-map}

\begin{lstlisting}
AGENTS.md                                      repository-specific agent entry point
README.md                                      setup, architecture, build, test, and package commands
Directory.Build.props                         common C# compiler and reproducibility settings
Directory.Packages.props                      centrally pinned third-party packages
global.json                                   .NET 10.0.400 and Microsoft Testing Platform selection

FirstGearBank.Core/
  FirstGearBank.Core.csproj                    engine-independent library
  CoreAssemblyMarker.cs                        dependency-boundary test anchor
  Domain/                                      currencies, identifiers, results, revisions
  Money/                                       fixed-point arithmetic, parsing, formatting, caps
  Time/                                        normalized financial and audit timestamps
  Configuration/                               JSONC parsing, defaults, validation, activation candidates
  Ledger/                                      typed immutable records, postings, hash chain, validation
  Projection/                                  replay-built account and system views
  Requests/                                    scope sequencing, payload digests, cached terminal results
  Coordination/                                serialized command admission and immutable publication
  Security/                                    production entropy port/implementation
  Settlement/                                  recovery evidence contracts and authentication
  Persistence/                                 section framing, codecs, migrations, limits, quarantine

First Gear Bank/
  FirstGearBankModSystem.cs                    Vintage Story registration and lifecycle wiring only
  Runtime/FirstGearBankServerHost.cs           testable startup/load/save/dispose orchestration
  Persistence/VintageStoryWorldStateStore.cs   adapter for the namespaced install marker and state envelope
  Configuration/VintageStoryConfigurationStore.cs
                                                comment-preserving config file ownership
  Diagnostics/FgbFileLogger.cs                 bounded file sink with exact severity labels
  Diagnostics/VintageStoryBankLogger.cs        file plus native Vintage Story log fan-out

CakeBuild/
  BuildPaths.cs                                distinct repository/project/stage/release/deploy paths
  ModManifest.cs                               build-only manifest parsing with no game dependency
  Packaging/                                   deterministic ZIP construction, inspection, safe selection
  Tasks/                                       unique, nondestructive build/test/package/deploy graph

tests/
  FirstGearBank.Core.Tests/                    deterministic domain and state-machine tests
  FirstGearBank.Adapter.Tests/                 host/config/log tests behind project-owned fakes
  CakeBuild.Tests/                             packaging and deployment safety tests

docs/research/                                 source-linked Vintage Story evidence
docs/architecture/                             recovery, identity/time, persistence, and index decisions
docs/testing/                                  controlled 1.22.7 engine smoke procedures and results
\end{lstlisting}

\subsubsection{Task 1: Record the 1.22.7 Evidence and Foundation
Decisions}\label{task-1-record-the-1227-evidence-and-foundation-decisions}

\textbf{Files:}

\begin{itemize}
\tightlist
\item
  Create:
  \texttt{docs/}\allowbreak{}\texttt{research/}\allowbreak{}\texttt{vintage-}\allowbreak{}\texttt{story-}\allowbreak{}\texttt{1.}\allowbreak{}\texttt{22.}\allowbreak{}\texttt{7-}\allowbreak{}\texttt{api-}\allowbreak{}\texttt{evidence.}\allowbreak{}\texttt{md}
\item
  Create:
  \texttt{docs/}\allowbreak{}\texttt{architecture/}\allowbreak{}\texttt{0001-}\allowbreak{}\texttt{physical-}\allowbreak{}\texttt{settlement-}\allowbreak{}\texttt{recovery-}\allowbreak{}\texttt{threat-}\allowbreak{}\texttt{model.}\allowbreak{}\texttt{md}
\item
  Create:
  \texttt{docs/}\allowbreak{}\texttt{architecture/}\allowbreak{}\texttt{0002-}\allowbreak{}\texttt{economically-}\allowbreak{}\texttt{accrued-}\allowbreak{}\texttt{liability-}\allowbreak{}\texttt{index.}\allowbreak{}\texttt{md}
\item
  Create:
  \texttt{docs/}\allowbreak{}\texttt{architecture/}\allowbreak{}\texttt{0003-}\allowbreak{}\texttt{domain-}\allowbreak{}\texttt{identifiers-}\allowbreak{}\texttt{and-}\allowbreak{}\texttt{financial-}\allowbreak{}\texttt{time.}\allowbreak{}\texttt{md}
\item
  Create:
  \texttt{docs/}\allowbreak{}\texttt{architecture/}\allowbreak{}\texttt{0004-}\allowbreak{}\texttt{persistence-}\allowbreak{}\texttt{sections-}\allowbreak{}\texttt{and-}\allowbreak{}\texttt{revisions.}\allowbreak{}\texttt{md}
\item
  Create:
  \texttt{docs/}\allowbreak{}\texttt{testing/}\allowbreak{}\texttt{foundation-}\allowbreak{}\texttt{engine-}\allowbreak{}\texttt{smoke.}\allowbreak{}\texttt{md}
\item
  Create: same-basename \passthrough{\lstinline!.tex!} files for all six
  Markdown documents through Pandoc
\item
  Generate but do not commit: same-basename
  \passthrough{\lstinline!.pdf!} files for rendered review
\end{itemize}

\textbf{Interfaces:}

\begin{itemize}
\tightlist
\item
  Consumes: the locked v03 specification and the indexed 1.22.7
  reference at
  \texttt{.}\allowbreak{}\texttt{.}\allowbreak{}\texttt{/}\allowbreak{}\texttt{Vintage\ }\allowbreak{}\texttt{Story\ }\allowbreak{}\texttt{Reference}.
\item
  Produces: source-linked go/no-go evidence and the serialized
  conventions consumed by every following task.
\end{itemize}

\begin{itemize}
\tightlist
\item[$\square$]
  \textbf{Step 1: Re-run narrow reference queries and capture the exact
  source paths}
\end{itemize}

\begin{lstlisting}
$db = "..\Vintage Story Reference\versions\1.22.7\Index\code.sqlite"
python -B "..\Vintage Story Reference\tools\code-index\query_code_index.py" --db $db --type ModSystem --limit 10
python -B "..\Vintage Story Reference\tools\code-index\query_code_index.py" --db $db --type ISaveGame --limit 10
python -B "..\Vintage Story Reference\tools\code-index\query_code_index.py" --db $db --type IServerEventAPI --limit 10
python -B "..\Vintage Story Reference\tools\code-index\query_code_index.py" --db $db --type IServerPlayer --limit 10
python -B "..\Vintage Story Reference\tools\code-index\query_code_index.py" --db $db --type INetworkAPI --limit 10
python -B "..\Vintage Story Reference\tools\code-index\query_code_index.py" --db $db --type ILogger --limit 10
\end{lstlisting}

Expected: every query resolves only 1.22.7 symbols; no 1.21 or earlier
type is used as implementation authority.

\begin{itemize}
\tightlist
\item[$\square$]
  \textbf{Step 2: Write the API evidence matrix}
\end{itemize}

Use this exact column contract and fill each row with a concrete source
file and line span:

\begin{lstlisting}
| Capability | Verified symbol or asset | 1.22.7 source | Conclusion | Dependent phase |
|---|---|---|---|---|
| Common/server lifecycle | `ModSystem.Start`, `StartServerSide`, `Dispose` | `API/VintagestoryAPI_decompiled_1.22.7/Vintagestory/API/Common/ModSystem.cs:74,112,119` | Verified public API | Foundation |
| World state | `ISaveGame.GetData`, `StoreData`, `SavegameIdentifier` | `API/VintagestoryAPI_decompiled_1.22.7/Vintagestory/API/Server/ISaveGame.cs:30,54,61` | Verified blob API; no transaction/fsync guarantee | Foundation |
| Save events | `SaveGameCreated`, `SaveGameLoaded`, `GameWorldSave` | `API/VintagestoryAPI_decompiled_1.22.7/Vintagestory/API/Server/IServerEventAPI.cs:124,129,142-144` | Verified pre-write staging event; no post-save acknowledgment | Foundation |
| Player capsule | `IServerPlayer.SetModData`, `GetModData` | `API/VintagestoryAPI_decompiled_1.22.7/Vintagestory/API/Server/IServerPlayer.cs:114,116` | API exists; co-serialization with inventory unproved | Settlement spike |
| Reusable recipe input | `CraftingRecipeIngredient.Consume`; vanilla glider/copy recipes | `API/VintagestoryAPI_decompiled_1.22.7/Vintagestory/API/Common/CraftingRecipeIngredient.cs:136`; `Assets/survival/recipes/grid/glider.json:8`; `schematiccopy.json:5` | Use `consume: false`, never `isTool` | Charter |
\end{lstlisting}

Also include verified rows for current \passthrough{\lstinline!Playing!}
connections, typed networking, registrations, entity spawning, GUI,
commands/permissions, the data path, rusty/temporal item codes, and
asset layout. Mark player/inventory durability, Charter action
interception, Banker base/dialogue, generated-trader identity, and
absent optional mods as explicit gates.

\begin{itemize}
\tightlist
\item[$\square$]
  \textbf{Step 3: Record the physical-settlement threat model}
\end{itemize}

The decision record must contain this trust statement verbatim:

\begin{lstlisting}
The authenticated server and its world save are trusted.  Clients, packet contents, item attributes supplied by clients,
stale player snapshots, filenames, and unauthenticated recovery bytes are not trusted.  A valid capsule proves only the
canonical evidence it authenticates; it does not prove an unstated save-order guarantee.  If the Bank record, recovery key,
world binding, player binding, manifest, fingerprint, phase, or authentication tag cannot be reconciled, First Gear Bank
quarantines the affected physical settlement and neither creates nor destroys compensating value by inference.
\end{lstlisting}

Specify HMAC-SHA-256, a random 256-bit world recovery key, binding to
the hash of \texttt{ISaveGame.}\allowbreak{}\texttt{SavegameIdentifier},
canonical length-prefixed UTF-8/binary encoding, constant-time tag
comparison, key nonrotation while evidence exists, and raw-evidence
retention.

\begin{itemize}
\tightlist
\item[$\square$]
  \textbf{Step 4: Run and document the accrued-liability counterexample
  and index gate}
\end{itemize}

Use this minimal counterexample in the decision record:

\begin{lstlisting}
balances = [1, 1] bank unit
factor   = 1.5
sum(roundToEven(factor * each balance)) = 2 + 2 = 4
roundToEven(factor * sum(balances))      = roundToEven(3) = 3
\end{lstlisting}

The record must require an indexed algorithm to match the per-account
month-by-month function for adversarial balances, partial months, caps,
and account mutations. If no exact sublinear algorithm is proved, record
\passthrough{\lstinline!Gate result: failed!} and stop before Phase 4 of
the roadmap; do not substitute aggregate rounding.

\begin{itemize}
\tightlist
\item[$\square$]
  \textbf{Step 5: Record identifier, time, section, and revision
  encodings}
\end{itemize}

\begin{lstlisting}
GUID-backed IDs: 16 canonical bytes in binary; lowercase 32-hex characters in diagnostic text.
Player key: exact authenticated engine UID internally; bounded UTF-8 string; never normalized or player-facing.
Financial instant: nonnegative completed financial months plus decimal fractional month in [0, 1).
Journal sequence: signed Int64 beginning at 1 and increasing contiguously.
Global revision: signed Int64 beginning at 0; one increment for every published state revision.
Subsystem revision: one signed Int64 per independently persisted domain; increment only when that domain changes.
Hash domains: explicit ASCII prefix and version followed by canonical length-prefixed bytes; never `GetHashCode`.
\end{lstlisting}

\begin{itemize}
\tightlist
\item[$\square$]
  \textbf{Step 6: Write the controlled engine smoke matrix}
\end{itemize}

The procedure must use a disposable 1.22.7 test world and record, for
each run, game version, mod ZIP hash, world backup, action, expected
state, observed state, FGB log excerpt, and pass/fail. Include fresh
load, save/restart, missing blob, corrupt copied blob, player-mod-data
persistence, and absence of all optional mods. List
inventory-plus-capsule save ordering and server shutdown during every
physical-settlement phase as explicitly deferred, nonblocking Phase 3
rows: Phase 1 has no inventory mutation capable of producing those
phases and must retain
\passthrough{\lstinline!UnprovedCoSerialization!}.

\begin{itemize}
\tightlist
\item[$\square$]
  \textbf{Step 7: Render, validate, and commit the documentation
  sources}
\end{itemize}

For each substantial Markdown file created by this task, generate a
same-basename standalone LaTeX source with Pandoc, compile a
same-basename PDF with XeLaTeX, render the PDF pages through Poppler,
and inspect them for clipping, missing text, bad glyphs, and blank
pages. Regenerate from Markdown after every correction; never edit
generated LaTeX or PDF directly. Commit Markdown and LaTeX, but leave
PDFs uncommitted as required by project hygiene.

\begin{lstlisting}
$ErrorActionPreference = "Stop"
function Invoke-NativeChecked([string] $tool, [string[]] $arguments) {
    & $tool @arguments
    if ($LASTEXITCODE -ne 0) {
        throw "$tool failed with exit code $LASTEXITCODE"
    }
}

$sourceDocs = @(
    "docs/research/vintage-story-1.22.7-api-evidence.md",
    "docs/architecture/0001-physical-settlement-recovery-threat-model.md",
    "docs/architecture/0002-economically-accrued-liability-index.md",
    "docs/architecture/0003-domain-identifiers-and-financial-time.md",
    "docs/architecture/0004-persistence-sections-and-revisions.md",
    "docs/testing/foundation-engine-smoke.md"
)
$systemTemp = [IO.Path]::GetFullPath([IO.Path]::GetTempPath())
$renderRoot = Join-Path $systemTemp ("fgb-docs-" + [Guid]::NewGuid().ToString("N"))
[IO.Directory]::CreateDirectory($renderRoot) | Out-Null
foreach ($sourceDoc in $sourceDocs) {
    $tex = [IO.Path]::ChangeExtension($sourceDoc, ".tex")
    $pdf = [IO.Path]::ChangeExtension($sourceDoc, ".pdf")
    $stem = [IO.Path]::GetFileNameWithoutExtension($sourceDoc)
    $documentOutput = Join-Path $renderRoot $stem
    $pageOutput = Join-Path $documentOutput "pages"
    [IO.Directory]::CreateDirectory($pageOutput) | Out-Null

    Invoke-NativeChecked "pandoc" @("--standalone", "--from=gfm", "--to=latex", "--output=$tex", $sourceDoc)
    Invoke-NativeChecked "xelatex" @("-interaction=nonstopmode", "-halt-on-error", "-output-directory=$documentOutput", $tex)
    Invoke-NativeChecked "xelatex" @("-interaction=nonstopmode", "-halt-on-error", "-output-directory=$documentOutput", $tex)
    Copy-Item -Force -LiteralPath (Join-Path $documentOutput "$stem.pdf") -Destination $pdf
    Invoke-NativeChecked "pdftoppm" @("-png", "-r", "144", $pdf, (Join-Path $pageOutput "page"))
}
Write-Output "FGB_RENDER_ROOT=$renderRoot"
\end{lstlisting}

Stop here. Retain the printed render root and inspect every generated
\texttt{pages/}\allowbreak{}\texttt{page-}\allowbreak{}\texttt{*.}\allowbreak{}\texttt{png}
with the image viewer. Record the page count and pass/fail for each
document; correct Markdown and repeat the complete render if any page
fails. Only after all pages pass, reassign
\passthrough{\lstinline!$renderRoot!} to the exact
\passthrough{\lstinline!FGB\_RENDER\_ROOT!} path if the shell changed
and run this guarded cleanup and source check:

\begin{lstlisting}
$systemTemp = [IO.Path]::GetFullPath([IO.Path]::GetTempPath())
if ([string]::IsNullOrWhiteSpace($renderRoot)) { throw "The inspected render root is required." }
$resolvedRenderRoot = [IO.Path]::GetFullPath($renderRoot)
$requiredPrefix = $systemTemp.TrimEnd([IO.Path]::DirectorySeparatorChar) + [IO.Path]::DirectorySeparatorChar
if (-not $resolvedRenderRoot.StartsWith($requiredPrefix, [StringComparison]::OrdinalIgnoreCase) -or
    [IO.Path]::GetFileName($resolvedRenderRoot) -notmatch '^fgb-docs-[0-9a-f]{32}$') {
    throw "Refusing to remove an unexpected render directory: $resolvedRenderRoot"
}
[IO.Directory]::Delete($resolvedRenderRoot, $true)
rg -n "Verified public|unproved|Gate result|quarantine|1\.22\.7" docs/research docs/architecture docs/testing
git diff --check
git add docs/research/vintage-story-1.22.7-api-evidence.md `
  docs/research/vintage-story-1.22.7-api-evidence.tex `
  docs/architecture/0001-physical-settlement-recovery-threat-model.md `
  docs/architecture/0001-physical-settlement-recovery-threat-model.tex `
  docs/architecture/0002-economically-accrued-liability-index.md `
  docs/architecture/0002-economically-accrued-liability-index.tex `
  docs/architecture/0003-domain-identifiers-and-financial-time.md `
  docs/architecture/0003-domain-identifiers-and-financial-time.tex `
  docs/architecture/0004-persistence-sections-and-revisions.md `
  docs/architecture/0004-persistence-sections-and-revisions.tex `
  docs/testing/foundation-engine-smoke.md docs/testing/foundation-engine-smoke.tex
git commit -m "docs: record Vintage Story foundation evidence"
\end{lstlisting}

Expected: the evidence distinguishes verified API facts from runtime
observations and unresolved gates.

\subsubsection{Task 2: Establish the Project Graph, Test Runner, and
Repository
Hygiene}\label{task-2-establish-the-project-graph-test-runner-and-repository-hygiene}

\textbf{Files:}

\begin{itemize}
\tightlist
\item
  Create: \passthrough{\lstinline!AGENTS.md!}
\item
  Create: \passthrough{\lstinline!README.md!}
\item
  Create: \passthrough{\lstinline!.editorconfig!}
\item
  Create: \passthrough{\lstinline!.gitattributes!}
\item
  Create: \passthrough{\lstinline!global.json!}
\item
  Create: \passthrough{\lstinline!Directory.Build.props!}
\item
  Create:
  \texttt{FirstGearBank.}\allowbreak{}\texttt{Core/}\allowbreak{}\texttt{FirstGearBank.}\allowbreak{}\texttt{Core.}\allowbreak{}\texttt{csproj}
\item
  Create:
  \texttt{FirstGearBank.}\allowbreak{}\texttt{Core/}\allowbreak{}\texttt{CoreAssemblyMarker.}\allowbreak{}\texttt{cs}
\item
  Create:
  \texttt{tests/}\allowbreak{}\texttt{FirstGearBank.}\allowbreak{}\texttt{Core.}\allowbreak{}\texttt{Tests/}\allowbreak{}\texttt{FirstGearBank.}\allowbreak{}\texttt{Core.}\allowbreak{}\texttt{Tests.}\allowbreak{}\texttt{csproj}
\item
  Create:
  \texttt{tests/}\allowbreak{}\texttt{FirstGearBank.}\allowbreak{}\texttt{Core.}\allowbreak{}\texttt{Tests/}\allowbreak{}\texttt{Architecture/}\allowbreak{}\texttt{CoreAssemblyBoundaryTests.}\allowbreak{}\texttt{cs}
\item
  Create:
  \texttt{tests/}\allowbreak{}\texttt{FirstGearBank.}\allowbreak{}\texttt{Adapter.}\allowbreak{}\texttt{Tests/}\allowbreak{}\texttt{FirstGearBank.}\allowbreak{}\texttt{Adapter.}\allowbreak{}\texttt{Tests.}\allowbreak{}\texttt{csproj}
\item
  Create:
  \texttt{tests/}\allowbreak{}\texttt{CakeBuild.}\allowbreak{}\texttt{Tests/}\allowbreak{}\texttt{CakeBuild.}\allowbreak{}\texttt{Tests.}\allowbreak{}\texttt{csproj}
\item
  Create:
  \texttt{FirstGearBank.}\allowbreak{}\texttt{Core/}\allowbreak{}\texttt{packages.}\allowbreak{}\texttt{lock.}\allowbreak{}\texttt{json}
  through \passthrough{\lstinline!dotnet restore!}
\item
  Create:
  \texttt{First\ }\allowbreak{}\texttt{Gear\ }\allowbreak{}\texttt{Bank/}\allowbreak{}\texttt{packages.}\allowbreak{}\texttt{lock.}\allowbreak{}\texttt{json}
  through \passthrough{\lstinline!dotnet restore!}
\item
  Create:
  \texttt{CakeBuild/}\allowbreak{}\texttt{packages.}\allowbreak{}\texttt{lock.}\allowbreak{}\texttt{json}
  through \passthrough{\lstinline!dotnet restore!}
\item
  Create:
  \texttt{tests/}\allowbreak{}\texttt{FirstGearBank.}\allowbreak{}\texttt{Core.}\allowbreak{}\texttt{Tests/}\allowbreak{}\texttt{packages.}\allowbreak{}\texttt{lock.}\allowbreak{}\texttt{json}
  through \passthrough{\lstinline!dotnet restore!}
\item
  Create:
  \texttt{tests/}\allowbreak{}\texttt{FirstGearBank.}\allowbreak{}\texttt{Adapter.}\allowbreak{}\texttt{Tests/}\allowbreak{}\texttt{packages.}\allowbreak{}\texttt{lock.}\allowbreak{}\texttt{json}
  through \passthrough{\lstinline!dotnet restore!}
\item
  Create:
  \texttt{tests/}\allowbreak{}\texttt{CakeBuild.}\allowbreak{}\texttt{Tests/}\allowbreak{}\texttt{packages.}\allowbreak{}\texttt{lock.}\allowbreak{}\texttt{json}
  through \passthrough{\lstinline!dotnet restore!}
\item
  Modify: \passthrough{\lstinline!.gitignore:1!}
\item
  Modify: \passthrough{\lstinline!First Gear Bank.sln:1!}
\item
  Modify:
  \texttt{First\ }\allowbreak{}\texttt{Gear\ }\allowbreak{}\texttt{Bank/}\allowbreak{}\texttt{First\ }\allowbreak{}\texttt{Gear\ }\allowbreak{}\texttt{Bank.}\allowbreak{}\texttt{csproj:}\allowbreak{}\texttt{1}
\end{itemize}

\textbf{Interfaces:}

\begin{itemize}
\tightlist
\item
  Consumes: SDK 10.0.400 and
  \texttt{VINTAGE\_}\allowbreak{}\texttt{STORY/}\allowbreak{}\texttt{VintagestoryAPI.}\allowbreak{}\texttt{dll}
  version 1.22.7.0.
\item
  Produces: \passthrough{\lstinline!FirstGearBank.Core!},
  \passthrough{\lstinline!FirstGearBank!}, and three test projects under
  one reproducible solution.
\end{itemize}

\begin{itemize}
\tightlist
\item[$\square$]
  \textbf{Step 1: Add the SDK and compiler contracts}
\end{itemize}

Create \passthrough{\lstinline!global.json!} with:

\begin{lstlisting}
{
  "sdk": {
    "version": "10.0.400",
    "rollForward": "latestPatch"
  },
  "test": {
    "runner": "Microsoft.Testing.Platform"
  }
}
\end{lstlisting}

Create \passthrough{\lstinline!Directory.Build.props!} with
\passthrough{\lstinline!net10.0!}, nullable annotations, implicit
usings, deterministic builds,
\passthrough{\lstinline!LangVersion=14.0!}, generated-file exclusions,
and \texttt{RestorePackagesWithLockFile=}\allowbreak{}\texttt{true}. Do
not globally suppress NuGet audit, nullable, or analyzer warnings.

\begin{itemize}
\tightlist
\item[$\square$]
  \textbf{Step 2: Create the core and test project files before the
  marker type exists}
\end{itemize}

Use this test-project shape for each xUnit project, changing project
references per target:

\begin{lstlisting}[language=XML]
<Project Sdk="Microsoft.NET.Sdk">
  <PropertyGroup>
    <OutputType>Exe</OutputType>
    <IsPackable>false</IsPackable>
  </PropertyGroup>
  <ItemGroup>
    <Using Include="Xunit" />
    <PackageReference Include="xunit.v3.mtp-v2" Version="4.0.0" />
  </ItemGroup>
  <ItemGroup>
    <ProjectReference Include="..\..\FirstGearBank.Core\FirstGearBank.Core.csproj" />
  </ItemGroup>
</Project>
\end{lstlisting}

\texttt{FirstGearBank.}\allowbreak{}\texttt{Adapter.}\allowbreak{}\texttt{Tests}
references both core and
\passthrough{\lstinline!First Gear Bank.csproj!}; its test-only Vintage
Story API reference may copy locally so the runner can load adapter
types. The packaged adapter keeps
\passthrough{\lstinline!Private="false"!}.

\begin{itemize}
\tightlist
\item[$\square$]
  \textbf{Step 3: Write the failing dependency-boundary test}
\end{itemize}

\begin{lstlisting}
/*
 * Protects the engine-independent assembly boundary required by the locked architecture.
 */

using FirstGearBank.Core;

namespace FirstGearBank.Core.Tests.Architecture;

public sealed class CoreAssemblyBoundaryTests
{



    //// Verifies that the banking core cannot acquire an engine reference unnoticed.
    ////
    [Fact]
    public void Core_assembly_references_no_vintage_story_assembly()
    {
        var references = typeof(CoreAssemblyMarker).Assembly.GetReferencedAssemblies();

        Assert.DoesNotContain(
            references,
            reference => reference.Name?.StartsWith("Vintagestory", StringComparison.OrdinalIgnoreCase) == true);
    }



}
\end{lstlisting}

\begin{itemize}
\tightlist
\item[$\square$]
  \textbf{Step 4: Run the test and observe the intended compile failure}
\end{itemize}

\begin{lstlisting}
dotnet test tests/FirstGearBank.Core.Tests/FirstGearBank.Core.Tests.csproj -c Release
\end{lstlisting}

Expected: FAIL because \passthrough{\lstinline!CoreAssemblyMarker!} does
not exist.

\begin{itemize}
\tightlist
\item[$\square$]
  \textbf{Step 5: Add the minimal core marker and project reference}
\end{itemize}

\begin{lstlisting}
/*
 * Provides a stable assembly anchor for dependency and packaging verification.
 */

namespace FirstGearBank.Core;

public static class CoreAssemblyMarker
{
}
\end{lstlisting}

Set the adapter project to
\texttt{AssemblyName=}\allowbreak{}\texttt{FirstGearBank},
\texttt{RootNamespace=}\allowbreak{}\texttt{FirstGearBank}, and add a
project reference to core. Keep only
\passthrough{\lstinline!VintagestoryAPI.dll!} as a non-copying engine
reference. Remove the custom publish output path and unused folder item;
packaging owns final layout. Expose core internals only to
\texttt{FirstGearBank.}\allowbreak{}\texttt{Core.}\allowbreak{}\texttt{Tests}
and the named \passthrough{\lstinline!FirstGearBank!} adapter assembly
through SDK \passthrough{\lstinline!InternalsVisibleTo!} items: tests
need deliberately corrupt fixtures, while the trusted server adapter
needs to construct permission grants after engine authorization. No
client assembly is a friend, and invariant-breaking constructors remain
nonpublic.

\begin{itemize}
\tightlist
\item[$\square$]
  \textbf{Step 6: Add repository instructions and reproducible setup
  documentation}
\end{itemize}

\passthrough{\lstinline!AGENTS.md!} must direct contributors to the four
checked-in standards, locked v03 spec, this plan/roadmap, and the
indexed 1.22.7 reference instructions.
\passthrough{\lstinline!README.md!} must document .NET 10.0.400,
\passthrough{\lstinline!VINTAGE\_STORY!}, the exact solution commands,
explicit packaging/deployment, the pure-core boundary, and the fact that
Phase 1 has no playable banking UI.

Configure UTF-8, final newlines, four-space C\# indentation, LF for
text, and binary treatment for
DLL/\allowbreak{}PNG/\allowbreak{}DOCX/\allowbreak{}PDF/\allowbreak{}ZIP.
Set
\texttt{dotnet\_}\allowbreak{}\texttt{style\_}\allowbreak{}\texttt{allow\_}\allowbreak{}\texttt{multiple\_}\allowbreak{}\texttt{blank\_}\allowbreak{}\texttt{lines\_}\allowbreak{}\texttt{experimental\ }\allowbreak{}\texttt{=}\allowbreak{}\texttt{\ }\allowbreak{}\texttt{true}
for C\# so Roslyn does not collapse the profile\textquotesingle s
mandatory three-blank-line callable spacing; the source-standard gate
still checks the exact project rule. Extend
\passthrough{\lstinline!.gitignore!} with
\texttt{Artifacts/}\allowbreak{}, \texttt{Releases/}\allowbreak{},
\texttt{TestResults/}\allowbreak{}, \texttt{logs/}\allowbreak{},
\texttt{.}\allowbreak{}\texttt{vs/}\allowbreak{}, and
\passthrough{\lstinline!*.zip!} without removing existing entries.
Ignore generated PDFs only in
\texttt{docs/}\allowbreak{}\texttt{research/}\allowbreak{},
\texttt{docs/}\allowbreak{}\texttt{architecture/}\allowbreak{},
\texttt{docs/}\allowbreak{}\texttt{testing/}\allowbreak{}, and
\texttt{docs/}\allowbreak{}\texttt{superpowers/}\allowbreak{}\texttt{plans/}\allowbreak{};
keep their canonical Markdown and generated LaTeX tracked.

\begin{itemize}
\tightlist
\item[$\square$]
  \textbf{Step 7: Add all projects to the solution and prove the
  boundary test passes}
\end{itemize}

\begin{lstlisting}
dotnet sln "First Gear Bank.sln" add "FirstGearBank.Core/FirstGearBank.Core.csproj"
dotnet sln "First Gear Bank.sln" add "tests/FirstGearBank.Core.Tests/FirstGearBank.Core.Tests.csproj"
dotnet sln "First Gear Bank.sln" add "tests/FirstGearBank.Adapter.Tests/FirstGearBank.Adapter.Tests.csproj"
dotnet sln "First Gear Bank.sln" add "tests/CakeBuild.Tests/CakeBuild.Tests.csproj"
dotnet test tests/FirstGearBank.Core.Tests/FirstGearBank.Core.Tests.csproj -c Release
\end{lstlisting}

Expected: PASS, and
\texttt{FirstGearBank.}\allowbreak{}\texttt{Core.}\allowbreak{}\texttt{csproj}
contains no assembly or package reference to Vintage Story.

\begin{itemize}
\tightlist
\item[$\square$]
  \textbf{Step 8: Build, inspect, and commit the project foundation}
\end{itemize}

\begin{lstlisting}
dotnet restore "First Gear Bank.sln"
dotnet build "First Gear Bank.sln" -c Release --no-restore
dotnet list FirstGearBank.Core/FirstGearBank.Core.csproj reference
git diff --check
git add AGENTS.md README.md .editorconfig .gitattributes .gitignore global.json Directory.Build.props `
  "First Gear Bank.sln" "First Gear Bank/First Gear Bank.csproj" `
  "First Gear Bank/packages.lock.json" CakeBuild/packages.lock.json FirstGearBank.Core tests
git commit -m "build: establish core and test project foundation"
\end{lstlisting}

Expected: zero build errors and no
\passthrough{\lstinline!Vintagestory.*!} core reference.

\subsubsection{Task 3: Build Deterministic Package and Deployment-Safety
Primitives}\label{task-3-build-deterministic-package-and-deployment-safety-primitives}

\textbf{Files:}

\begin{itemize}
\tightlist
\item
  Create:
  \texttt{CakeBuild/}\allowbreak{}\texttt{BuildPaths.}\allowbreak{}\texttt{cs}
\item
  Create:
  \texttt{CakeBuild/}\allowbreak{}\texttt{RepositoryLocator.}\allowbreak{}\texttt{cs}
\item
  Create:
  \texttt{CakeBuild/}\allowbreak{}\texttt{ModManifest.}\allowbreak{}\texttt{cs}
\item
  Create:
  \texttt{CakeBuild/}\allowbreak{}\texttt{Packaging/}\allowbreak{}\texttt{SemanticVersion.}\allowbreak{}\texttt{cs}
\item
  Create:
  \texttt{CakeBuild/}\allowbreak{}\texttt{Packaging/}\allowbreak{}\texttt{PackageArchiveLimits.}\allowbreak{}\texttt{cs}
\item
  Create:
  \texttt{CakeBuild/}\allowbreak{}\texttt{Packaging/}\allowbreak{}\texttt{PackageArchiveWriter.}\allowbreak{}\texttt{cs}
\item
  Create:
  \texttt{CakeBuild/}\allowbreak{}\texttt{Packaging/}\allowbreak{}\texttt{PackageArchiveInspector.}\allowbreak{}\texttt{cs}
\item
  Create:
  \texttt{CakeBuild/}\allowbreak{}\texttt{Packaging/}\allowbreak{}\texttt{PackageInspection.}\allowbreak{}\texttt{cs}
\item
  Create:
  \texttt{CakeBuild/}\allowbreak{}\texttt{Packaging/}\allowbreak{}\texttt{DeploymentPackageSelector.}\allowbreak{}\texttt{cs}
\item
  Create:
  \texttt{CakeBuild/}\allowbreak{}\texttt{Packaging/}\allowbreak{}\texttt{DeploymentPathValidator.}\allowbreak{}\texttt{cs}
\item
  Create:
  \texttt{CakeBuild/}\allowbreak{}\texttt{Packaging/}\allowbreak{}\texttt{DeploymentPathValidation.}\allowbreak{}\texttt{cs}
\item
  Create:
  \texttt{CakeBuild/}\allowbreak{}\texttt{Packaging/}\allowbreak{}\texttt{ValidatedDeploymentDirectory.}\allowbreak{}\texttt{cs}
\item
  Create:
  \texttt{CakeBuild/}\allowbreak{}\texttt{Packaging/}\allowbreak{}\texttt{IDeploymentFileOperations.}\allowbreak{}\texttt{cs}
\item
  Create:
  \texttt{CakeBuild/}\allowbreak{}\texttt{Packaging/}\allowbreak{}\texttt{PhysicalDeploymentFileOperations.}\allowbreak{}\texttt{cs}
\item
  Create:
  \texttt{CakeBuild/}\allowbreak{}\texttt{Packaging/}\allowbreak{}\texttt{PackageDeploymentInstaller.}\allowbreak{}\texttt{cs}
\item
  Create:
  \texttt{CakeBuild/}\allowbreak{}\texttt{Packaging/}\allowbreak{}\texttt{PackageDeploymentResult.}\allowbreak{}\texttt{cs}
\item
  Create:
  \texttt{tests/}\allowbreak{}\texttt{CakeBuild.}\allowbreak{}\texttt{Tests/}\allowbreak{}\texttt{BuildPathsTests.}\allowbreak{}\texttt{cs}
\item
  Create:
  \texttt{tests/}\allowbreak{}\texttt{CakeBuild.}\allowbreak{}\texttt{Tests/}\allowbreak{}\texttt{RepositoryLocatorTests.}\allowbreak{}\texttt{cs}
\item
  Create:
  \texttt{tests/}\allowbreak{}\texttt{CakeBuild.}\allowbreak{}\texttt{Tests/}\allowbreak{}\texttt{SemanticVersionTests.}\allowbreak{}\texttt{cs}
\item
  Create:
  \texttt{tests/}\allowbreak{}\texttt{CakeBuild.}\allowbreak{}\texttt{Tests/}\allowbreak{}\texttt{ModManifestTests.}\allowbreak{}\texttt{cs}
\item
  Create:
  \texttt{tests/}\allowbreak{}\texttt{CakeBuild.}\allowbreak{}\texttt{Tests/}\allowbreak{}\texttt{PackageArchiveWriterTests.}\allowbreak{}\texttt{cs}
\item
  Create:
  \texttt{tests/}\allowbreak{}\texttt{CakeBuild.}\allowbreak{}\texttt{Tests/}\allowbreak{}\texttt{PackageArchiveInspectorTests.}\allowbreak{}\texttt{cs}
\item
  Create:
  \texttt{tests/}\allowbreak{}\texttt{CakeBuild.}\allowbreak{}\texttt{Tests/}\allowbreak{}\texttt{DeploymentPackageSelectorTests.}\allowbreak{}\texttt{cs}
\item
  Create:
  \texttt{tests/}\allowbreak{}\texttt{CakeBuild.}\allowbreak{}\texttt{Tests/}\allowbreak{}\texttt{DeploymentPathValidatorTests.}\allowbreak{}\texttt{cs}
\item
  Create:
  \texttt{tests/}\allowbreak{}\texttt{CakeBuild.}\allowbreak{}\texttt{Tests/}\allowbreak{}\texttt{PackageDeploymentInstallerTests.}\allowbreak{}\texttt{cs}
\item
  Create:
  \texttt{tests/}\allowbreak{}\texttt{CakeBuild.}\allowbreak{}\texttt{Tests/}\allowbreak{}\texttt{TestSupport/}\allowbreak{}\texttt{TestArchive.}\allowbreak{}\texttt{cs}
\item
  Create:
  \texttt{tests/}\allowbreak{}\texttt{CakeBuild.}\allowbreak{}\texttt{Tests/}\allowbreak{}\texttt{TestSupport/}\allowbreak{}\texttt{TestDeploymentDirectory.}\allowbreak{}\texttt{cs}
\item
  Create:
  \texttt{tests/}\allowbreak{}\texttt{CakeBuild.}\allowbreak{}\texttt{Tests/}\allowbreak{}\texttt{TestSupport/}\allowbreak{}\texttt{TestManifest.}\allowbreak{}\texttt{cs}
\item
  Create:
  \texttt{tests/}\allowbreak{}\texttt{CakeBuild.}\allowbreak{}\texttt{Tests/}\allowbreak{}\texttt{TestSupport/}\allowbreak{}\texttt{FaultInjectingDeploymentFileOperations.}\allowbreak{}\texttt{cs}
\item
  Modify:
  \texttt{tests/}\allowbreak{}\texttt{CakeBuild.}\allowbreak{}\texttt{Tests/}\allowbreak{}\texttt{CakeBuild.}\allowbreak{}\texttt{Tests.}\allowbreak{}\texttt{csproj:}\allowbreak{}\texttt{1}
\end{itemize}

\textbf{Interfaces:}

\begin{itemize}
\tightlist
\item
  Consumes: the repository project graph from Task 2 and build-only BCL
  APIs.
\item
  Produces: tested, Cake-independent path, version, manifest, archive,
  and deployment-selection logic used by Task 4.
\end{itemize}

\begin{itemize}
\tightlist
\item[$\square$]
  \textbf{Step 1: Point the build tests at
  \passthrough{\lstinline!CakeBuild!} and write the failing path/version
  tests}
\end{itemize}

Add a project reference to
\texttt{.}\allowbreak{}\texttt{.}\allowbreak{}\texttt{/}\allowbreak{}\texttt{.}\allowbreak{}\texttt{.}\allowbreak{}\texttt{/}\allowbreak{}\texttt{CakeBuild/}\allowbreak{}\texttt{CakeBuild.}\allowbreak{}\texttt{csproj},
then add these cases:

\begin{lstlisting}
[Fact]
public void Repository_paths_preserve_the_adapter_directory_with_spaces()
{
    var paths = BuildPaths.FromRepositoryRoot(new DirectoryInfo(@"C:\repo\First Gear Bank"));

    Assert.Equal(
        @"C:\repo\First Gear Bank\First Gear Bank\First Gear Bank.csproj",
        paths.AdapterProject.FullName);
}

[Theory]
[InlineData("1.2.3", VersionComponent.Major, "2.0.0")]
[InlineData("1.2.3", VersionComponent.Minor, "1.3.0")]
[InlineData("1.2.3", VersionComponent.Patch, "1.2.4")]
public void Bump_changes_exactly_one_semantic_component(
    string source,
    VersionComponent component,
    string expected)
{
    Assert.Equal(expected, SemanticVersion.Parse(source).Bump(component).ToString());
}
\end{lstlisting}

\begin{itemize}
\tightlist
\item[$\square$]
  \textbf{Step 2: Run the focused tests and observe missing-type
  failures}
\end{itemize}

\begin{lstlisting}
dotnet test tests/CakeBuild.Tests/CakeBuild.Tests.csproj -c Release
\end{lstlisting}

Expected: FAIL because \passthrough{\lstinline!BuildPaths!},
\passthrough{\lstinline!SemanticVersion!}, and
\passthrough{\lstinline!VersionComponent!} do not exist.

\begin{itemize}
\tightlist
\item[$\square$]
  \textbf{Step 3: Implement path separation and strict semantic
  versions}
\end{itemize}

Use these public contracts:

\begin{lstlisting}
public sealed record BuildPaths(
    DirectoryInfo RepositoryRoot,
    FileInfo Solution,
    DirectoryInfo AdapterDirectory,
    FileInfo AdapterProject,
    FileInfo Manifest,
    DirectoryInfo PublishDirectory,
    DirectoryInfo StagingDirectory,
    DirectoryInfo ReleasesDirectory)
{
    public static BuildPaths FromRepositoryRoot(DirectoryInfo repositoryRoot);
}

public enum VersionComponent
{
    Major = 1,
    Minor = 2,
    Patch = 3
}

public readonly record struct SemanticVersion(int Major, int Minor, int Patch)
{
    public static SemanticVersion Parse(string text);
    public SemanticVersion Bump(VersionComponent component);
    public override string ToString();
}
\end{lstlisting}

\passthrough{\lstinline!Parse!} accepts exactly three nonnegative
invariant integer components and rejects whitespace, signs,
leading/trailing text, overflow, and extra components.
\texttt{BuildPaths.}\allowbreak{}\texttt{FromRepositoryRoot} resolves
absolute descendants and verifies that each build-owned path remains
under the supplied repository root.

\texttt{RepositoryLocator.}\allowbreak{}\texttt{Locate(}\allowbreak{}\texttt{DirectoryInfo\ }\allowbreak{}\texttt{start)}\allowbreak{}
checks only \passthrough{\lstinline!start!} and its ancestors for
\passthrough{\lstinline!First Gear Bank.sln!}; it never recursively
scans drives. Tests cover invocation from the repository root,
\passthrough{\lstinline!CakeBuild!}, a compiled
\passthrough{\lstinline!bin!} descendant, and a directory with no
matching ancestor.

\begin{itemize}
\tightlist
\item[$\square$]
  \textbf{Step 4: Prove path/version behavior is green}
\end{itemize}

\begin{lstlisting}
dotnet test tests/CakeBuild.Tests/CakeBuild.Tests.csproj -c Release
\end{lstlisting}

Expected: PASS for valid bumps and rejection tests for
\passthrough{\lstinline!1!}, \passthrough{\lstinline!1.2!},
\passthrough{\lstinline!1.2.3.4!}, \passthrough{\lstinline!-1.2.3!}, and
numeric overflow.

\begin{itemize}
\tightlist
\item[$\square$]
  \textbf{Step 5: Write failing manifest and archive-inspection tests}
\end{itemize}

Construct temporary ZIPs in the test process and assert all of these
rules:

\begin{lstlisting}
[Fact]
public void Valid_release_requires_both_mod_assemblies_at_the_root()
{
    using var archive = TestArchive.Create(
        ("modinfo.json", TestManifest.ValidBytes),
        ("FirstGearBank.dll", [1, 2, 3]));

    var result = PackageArchiveInspector.Inspect(archive.Path);

    Assert.Contains(result.Errors, error => error.Code == "MissingCoreAssembly");
}

[Theory]
[InlineData("../outside.txt")]
[InlineData("/rooted.txt")]
[InlineData("C:/rooted.txt")]
[InlineData("assets\\firstgearbank\\bad.json")]
public void Unsafe_entry_names_are_rejected(string entryName)
{
    using var archive = TestArchive.Create(
        ("modinfo.json", TestManifest.ValidBytes),
        ("FirstGearBank.dll", [1]),
        ("FirstGearBank.Core.dll", [2]),
        (entryName, [3]));

    Assert.False(PackageArchiveInspector.Inspect(archive.Path).IsValid);
}
\end{lstlisting}

Add explicit tests for wrong \passthrough{\lstinline!modid!}, wrong game
dependency, package/manifest version mismatch, duplicate names under
\passthrough{\lstinline!OrdinalIgnoreCase!}, malformed JSON assets, and
prohibited \passthrough{\lstinline!VintagestoryAPI.dll!}, Cake DLLs,
PDB, \passthrough{\lstinline!.deps.json!}, and
\passthrough{\lstinline!.runtimeconfig.json!} entries. Exercise each
archive bound at exactly the limit and at limit plus one: 10,000/10,001
entries, 64 KiB/64 KiB plus one manifest bytes, 64 MiB/64 MiB plus one
entry bytes, and 256 MiB/256 MiB plus one total declared uncompressed
bytes. Also exercise 256 MiB/plus one physical archive bytes, 8 MiB/plus
one central-directory bytes, 64 MiB/plus one declared compressed entry
bytes, and 256 MiB/plus one total declared compressed bytes. Include a
sparse oversized foreign ZIP and prove it is rejected from file length
before hashing or central-directory traversal. The fixtures may write
central-directory lengths or a sparse file length without allocating the
declared payload.

Keep \passthrough{\lstinline!TestArchive!},
\passthrough{\lstinline!TestDeploymentDirectory!}, and
\passthrough{\lstinline!TestManifest!} under the test project. They
create uniquely named temporary directories, expose exact fixture bytes,
and delete only their own resolved directory during
\passthrough{\lstinline!Dispose!}.

\begin{itemize}
\tightlist
\item[$\square$]
  \textbf{Step 6: Run the archive tests and observe missing behavior}
\end{itemize}

\begin{lstlisting}
dotnet test tests/CakeBuild.Tests/CakeBuild.Tests.csproj -c Release
\end{lstlisting}

Expected: FAIL because the manifest codec, archive writer, and inspector
are absent.

\begin{itemize}
\tightlist
\item[$\square$]
  \textbf{Step 7: Implement strict manifest parsing and deterministic
  ZIP construction}
\end{itemize}

Use these contracts:

\begin{lstlisting}
public sealed record ModManifest(
    string Type,
    string ModId,
    string Name,
    ImmutableArray<string> Authors,
    string Description,
    SemanticVersion Version,
    ImmutableDictionary<string, string> Dependencies);

public sealed record PackageInspection(
    bool IsValid,
    ModManifest? Manifest,
    ImmutableArray<PackageIssue> Errors,
    string? Sha256Hex);

public static class PackageArchiveLimits
{
    public const long MaximumArchiveBytes = 256L * 1024 * 1024;
    public const int MaximumCentralDirectoryBytes = 8 * 1024 * 1024;
    public const long MaximumCompressedEntryBytes = 64L * 1024 * 1024;
    public const long MaximumTotalCompressedBytes = 256L * 1024 * 1024;
}

public static class PackageArchiveWriter
{
    public static void Write(
        DirectoryInfo stagingDirectory,
        FileInfo destination,
        SemanticVersion expectedVersion);
}

public static class PackageArchiveInspector
{
    public static PackageInspection Inspect(
        FileInfo archive,
        SemanticVersion? expectedVersion = null);
}
\end{lstlisting}

Parse with \passthrough{\lstinline!System.Text.Json!}, require
\passthrough{\lstinline!type=code!},
\passthrough{\lstinline!modid=firstgearbank!}, and
\texttt{dependencies.}\allowbreak{}\texttt{game=}\allowbreak{}\texttt{1.}\allowbreak{}\texttt{22.}\allowbreak{}\texttt{7},
and reject missing/extra root placement for required files. Enumerate
staging files with ordinal normalized relative paths, reject unsafe
names and filesystem reparse points before opening the destination, and
assign every ZIP entry the fixed timestamp
\passthrough{\lstinline!2000-01-01T00:00:00Z!}. Bound inspection before
decompression to 10,000 entries, a 64 KiB root manifest, 64 MiB per
entry, 256 MiB total uncompressed content, and the
physical/\allowbreak{}compressed/\allowbreak{}central-directory limits
above. Read file length before opening
\passthrough{\lstinline!ZipArchive!}; validate central-directory range
and every declared compressed range against that length before entry
content is opened. Hash a file only after all structural bounds pass,
and do not hash unrelated foreign archives during deployment selection.
Write to a uniquely named same-directory temporary ZIP, inspect and hash
it, and only then atomically replace the destination. Before replacing
an existing release path, inspect it and refuse replacement unless its
root manifest also identifies \passthrough{\lstinline!firstgearbank!};
preserve both files if replacement cannot be completed safely.

\begin{itemize}
\tightlist
\item[$\square$]
  \textbf{Step 8: Verify deterministic archives and all rejection cases}
\end{itemize}

\begin{lstlisting}
dotnet test tests/CakeBuild.Tests/CakeBuild.Tests.csproj -c Release
\end{lstlisting}

Expected: PASS, including byte-identical SHA-256 hashes for two archives
built from identical staged bytes.

\begin{itemize}
\tightlist
\item[$\square$]
  \textbf{Step 9: Write failing deployment-selection, path, and
  transaction tests}
\end{itemize}

\begin{lstlisting}
[Fact]
public void Selection_uses_the_internal_mod_id_and_preserves_unreadable_archives()
{
    using var directory = TestDeploymentDirectory.Create(
        ("misleading-firstgearbank.zip", TestArchive.Foreign),
        ("unrelated-name.zip", TestArchive.FirstGearBank),
        ("damaged.zip", TestArchive.Unreadable));

    var result = DeploymentPackageSelector.FindExistingMatches(directory.Path, "firstgearbank");

    Assert.Equal(["unrelated-name.zip"], result.MatchingArchives.Select(file => file.Name));
    Assert.Contains(result.Warnings, warning => warning.Path.Name == "damaged.zip");
}
\end{lstlisting}

Add table-driven \texttt{DeploymentPathValidatorTests} for a relative
path, filesystem root, repository root/descendant,
publish/\allowbreak{}stage/\allowbreak{}release directory, a normal
external Mods directory, and every ancestor reparse point whose resolved
target enters a forbidden tree. Every invalid result has no
validated-directory token. Use the file-operations seam rather than
requiring OS symlink privileges.

\texttt{PackageDeploymentInstallerTests} use uniquely scoped temporary
directories and a fault-injecting file-operations adapter. For each
copy, source-hash, temporary-package inspection, existing-package backup
move, final rename, destination-hash, final inspection, and
backup-cleanup boundary, inject one failure and assert that no
unconfirmed archive is deleted, every old confirmed package is restored
when installation did not complete, the source is unchanged, and any
retained recovery file is reported. Also test hash mismatch, an existing
exact-match package, misleading filenames, foreign/unreadable ZIPs, and
successful removal of confirmed backups only after the new ZIP verifies.
Prove the install API has no raw-\passthrough{\lstinline!DirectoryInfo!}
overload, that a reparse target changed after initial validation fails
revalidation, and that every invalid/stale token path causes zero
mutating file-operation calls.

\begin{lstlisting}
dotnet test tests/CakeBuild.Tests/CakeBuild.Tests.csproj -c Release
\end{lstlisting}

Expected: FAIL because deployment selection by bounded internal manifest
identity is not implemented.

\begin{itemize}
\tightlist
\item[$\square$]
  \textbf{Step 10: Implement and verify safe selection}
\end{itemize}

\begin{lstlisting}
public sealed record DeploymentSelection(
    ImmutableArray<FileInfo> MatchingArchives,
    ImmutableArray<DeploymentWarning> Warnings);

public static class DeploymentPackageSelector
{
    public static DeploymentSelection FindExistingMatches(
        DirectoryInfo modsDirectory,
        string expectedModId);
}

public static class DeploymentPathValidator
{
    public static DeploymentPathValidation Validate(
        BuildPaths buildPaths,
        DirectoryInfo candidate,
        IDeploymentFileOperations fileOperations);
}

public sealed record DeploymentPathValidation(
    bool IsValid,
    ValidatedDeploymentDirectory? Directory,
    ImmutableArray<DeploymentIssue> Issues);

public sealed class ValidatedDeploymentDirectory
{
    public DirectoryInfo Path { get; }
    public string ResolvedPath { get; }

    internal ValidatedDeploymentDirectory(
        DirectoryInfo path,
        string resolvedPath,
        ImmutableArray<string> forbiddenRoots,
        ImmutableArray<string> resolvedAncestors);
    internal ImmutableArray<DeploymentIssue> Revalidate(IDeploymentFileOperations fileOperations);
}

public static class PackageDeploymentInstaller
{
    public static PackageDeploymentResult Install(
        FileInfo sourceArchive,
        ValidatedDeploymentDirectory modsDirectory,
        string expectedModId,
        IDeploymentFileOperations fileOperations);
}
\end{lstlisting}

Production supplies \texttt{PhysicalDeploymentFileOperations}; tests
supply the fault-injecting implementation. Both lexical absolute paths
and every existing ancestor\textquotesingle s resolved reparse target
must remain outside the repository and all build-owned trees. Only the
validator\textquotesingle s internal constructor can create the token,
and the token captures the normalized destination, forbidden roots, and
resolved-ancestor fingerprint. The installer revalidates that token
immediately before its first mutation and again before final rename; it
has no raw-path overload. It uses uniquely named same-directory
temporary/backup files, restores in reverse order on failure, and
returns every retained recovery path in
\passthrough{\lstinline!PackageDeploymentResult!}. It never accepts a
broad directory based only on string prefix; path comparison is
separator-aware and case-insensitive on Windows.

\begin{lstlisting}
dotnet test tests/CakeBuild.Tests/CakeBuild.Tests.csproj -c Release
\end{lstlisting}

Expected: PASS; malformed, unreadable, missing-manifest, and foreign
archives are never returned for deletion.

\begin{itemize}
\tightlist
\item[$\square$]
  \textbf{Step 11: Run all build-logic tests and commit}
\end{itemize}

\begin{lstlisting}
dotnet test tests/CakeBuild.Tests/CakeBuild.Tests.csproj -c Release
git diff --check
git add CakeBuild/BuildPaths.cs CakeBuild/RepositoryLocator.cs CakeBuild/ModManifest.cs CakeBuild/Packaging `
  tests/CakeBuild.Tests
git commit -m "build: add deterministic package safety primitives"
\end{lstlisting}

\subsubsection{Task 4: Replace the Unsafe Cake Task Graph and Correct
Mod
Metadata}\label{task-4-replace-the-unsafe-cake-task-graph-and-correct-mod-metadata}

\textbf{Files:}

\begin{itemize}
\tightlist
\item
  Create:
  \texttt{Directory.}\allowbreak{}\texttt{Packages.}\allowbreak{}\texttt{props}
\item
  Create:
  \texttt{CakeBuild/}\allowbreak{}\texttt{Tasks/}\allowbreak{}\texttt{ValidateEnvironmentTask.}\allowbreak{}\texttt{cs}
\item
  Create:
  \texttt{CakeBuild/}\allowbreak{}\texttt{Tasks/}\allowbreak{}\texttt{ValidateMetadataTask.}\allowbreak{}\texttt{cs}
\item
  Create:
  \texttt{CakeBuild/}\allowbreak{}\texttt{Tasks/}\allowbreak{}\texttt{RestoreTask.}\allowbreak{}\texttt{cs}
\item
  Create:
  \texttt{CakeBuild/}\allowbreak{}\texttt{Tasks/}\allowbreak{}\texttt{TestTask.}\allowbreak{}\texttt{cs}
\item
  Create:
  \texttt{CakeBuild/}\allowbreak{}\texttt{Tasks/}\allowbreak{}\texttt{StageTask.}\allowbreak{}\texttt{cs}
\item
  Create:
  \texttt{CakeBuild/}\allowbreak{}\texttt{Tasks/}\allowbreak{}\texttt{PackageTask.}\allowbreak{}\texttt{cs}
\item
  Create:
  \texttt{CakeBuild/}\allowbreak{}\texttt{Tasks/}\allowbreak{}\texttt{VerifyPackageTask.}\allowbreak{}\texttt{cs}
\item
  Create:
  \texttt{CakeBuild/}\allowbreak{}\texttt{Tasks/}\allowbreak{}\texttt{DefaultTask.}\allowbreak{}\texttt{cs}
\item
  Modify:
  \texttt{CakeBuild/}\allowbreak{}\texttt{CakeBuild.}\allowbreak{}\texttt{csproj:}\allowbreak{}\texttt{1}
\item
  Modify:
  \texttt{CakeBuild/}\allowbreak{}\texttt{BuildContext.}\allowbreak{}\texttt{cs:}\allowbreak{}\texttt{1}
\item
  Modify:
  \texttt{CakeBuild/}\allowbreak{}\texttt{Program.}\allowbreak{}\texttt{cs:}\allowbreak{}\texttt{1}
\item
  Modify:
  \texttt{CakeBuild/}\allowbreak{}\texttt{Tasks/}\allowbreak{}\texttt{BuildTask.}\allowbreak{}\texttt{cs:}\allowbreak{}\texttt{1}
\item
  Modify:
  \texttt{CakeBuild/}\allowbreak{}\texttt{Tasks/}\allowbreak{}\texttt{ValidateJsonTask.}\allowbreak{}\texttt{cs:}\allowbreak{}\texttt{1}
\item
  Modify:
  \texttt{CakeBuild/}\allowbreak{}\texttt{Tasks/}\allowbreak{}\texttt{DeployTask.}\allowbreak{}\texttt{cs:}\allowbreak{}\texttt{1}
\item
  Delete:
  \texttt{CakeBuild/}\allowbreak{}\texttt{Tasks/}\allowbreak{}\texttt{PackageModTask.}\allowbreak{}\texttt{cs}
\item
  Delete:
  \texttt{CakeBuild/}\allowbreak{}\texttt{Tasks/}\allowbreak{}\texttt{ZipModFolder.}\allowbreak{}\texttt{cs}
\item
  Modify:
  \texttt{First\ }\allowbreak{}\texttt{Gear\ }\allowbreak{}\texttt{Bank/}\allowbreak{}\texttt{modinfo.}\allowbreak{}\texttt{json:}\allowbreak{}\texttt{1}
\item
  Modify:
  \texttt{tests/}\allowbreak{}\texttt{FirstGearBank.}\allowbreak{}\texttt{Core.}\allowbreak{}\texttt{Tests/}\allowbreak{}\texttt{FirstGearBank.}\allowbreak{}\texttt{Core.}\allowbreak{}\texttt{Tests.}\allowbreak{}\texttt{csproj:}\allowbreak{}\texttt{1}
\item
  Modify:
  \texttt{tests/}\allowbreak{}\texttt{FirstGearBank.}\allowbreak{}\texttt{Core.}\allowbreak{}\texttt{Tests/}\allowbreak{}\texttt{packages.}\allowbreak{}\texttt{lock.}\allowbreak{}\texttt{json:}\allowbreak{}\texttt{1}
\item
  Modify:
  \texttt{tests/}\allowbreak{}\texttt{FirstGearBank.}\allowbreak{}\texttt{Adapter.}\allowbreak{}\texttt{Tests/}\allowbreak{}\texttt{FirstGearBank.}\allowbreak{}\texttt{Adapter.}\allowbreak{}\texttt{Tests.}\allowbreak{}\texttt{csproj:}\allowbreak{}\texttt{1}
\item
  Modify:
  \texttt{tests/}\allowbreak{}\texttt{FirstGearBank.}\allowbreak{}\texttt{Adapter.}\allowbreak{}\texttt{Tests/}\allowbreak{}\texttt{packages.}\allowbreak{}\texttt{lock.}\allowbreak{}\texttt{json:}\allowbreak{}\texttt{1}
\item
  Modify:
  \texttt{tests/}\allowbreak{}\texttt{CakeBuild.}\allowbreak{}\texttt{Tests/}\allowbreak{}\texttt{CakeBuild.}\allowbreak{}\texttt{Tests.}\allowbreak{}\texttt{csproj:}\allowbreak{}\texttt{1}
\item
  Modify:
  \texttt{tests/}\allowbreak{}\texttt{CakeBuild.}\allowbreak{}\texttt{Tests/}\allowbreak{}\texttt{packages.}\allowbreak{}\texttt{lock.}\allowbreak{}\texttt{json:}\allowbreak{}\texttt{1}
\item
  Modify:
  \texttt{FirstGearBank.}\allowbreak{}\texttt{Core/}\allowbreak{}\texttt{packages.}\allowbreak{}\texttt{lock.}\allowbreak{}\texttt{json:}\allowbreak{}\texttt{1}
\item
  Modify:
  \texttt{First\ }\allowbreak{}\texttt{Gear\ }\allowbreak{}\texttt{Bank/}\allowbreak{}\texttt{packages.}\allowbreak{}\texttt{lock.}\allowbreak{}\texttt{json:}\allowbreak{}\texttt{1}
\item
  Modify:
  \texttt{CakeBuild/}\allowbreak{}\texttt{packages.}\allowbreak{}\texttt{lock.}\allowbreak{}\texttt{json:}\allowbreak{}\texttt{1}
\item
  Modify: \passthrough{\lstinline!build.ps1:1!}
\item
  Modify: \passthrough{\lstinline!build.py:1!}
\item
  Modify: \passthrough{\lstinline!README.md:1!}
\end{itemize}

\textbf{Interfaces:}

\begin{itemize}
\tightlist
\item
  Consumes: the tested build primitives from Task 3.
\item
  Produces: unique \passthrough{\lstinline!Default!},
  \passthrough{\lstinline!Package!},
  \passthrough{\lstinline!VerifyPackage!}, and explicit
  \passthrough{\lstinline!Deploy!} targets with no implicit destructive
  action.
\end{itemize}

\begin{itemize}
\tightlist
\item[$\square$]
  \textbf{Step 1: Capture the current red integration baseline}
\end{itemize}

\begin{lstlisting}
dotnet run --project CakeBuild/CakeBuild.csproj -- --target=Build
\end{lstlisting}

Expected: FAIL while looking for
\texttt{First\_}\allowbreak{}\texttt{Gear\_}\allowbreak{}\texttt{Bank/}\allowbreak{}\texttt{modinfo.}\allowbreak{}\texttt{json},
proving that one identifier was incorrectly used as a directory, project
name, and mod ID.

\begin{itemize}
\tightlist
\item[$\square$]
  \textbf{Step 2: Centralize final package versions and remove
  build-only game dependencies}
\end{itemize}

Create:

\begin{lstlisting}[language=XML]
<Project>
  <PropertyGroup>
    <ManagePackageVersionsCentrally>true</ManagePackageVersionsCentrally>
  </PropertyGroup>
  <ItemGroup>
    <PackageVersion Include="Cake.Frosting" Version="6.2.0" />
    <PackageVersion Include="xunit.v3.mtp-v2" Version="4.0.0" />
  </ItemGroup>
</Project>
\end{lstlisting}

Remove explicit package versions from all project files.
\passthrough{\lstinline!CakeBuild.csproj!} keeps only
\passthrough{\lstinline!Cake.Frosting!} and removes Cake.Json,
Newtonsoft.Json, and \passthrough{\lstinline!VintagestoryAPI!}; Task
3\textquotesingle s \passthrough{\lstinline!System.Text.Json!} manifest
model replaces them.

Regenerate the committed lock files after this intentional dependency
change:

\begin{lstlisting}
dotnet restore "First Gear Bank.sln" --force-evaluate
\end{lstlisting}

\begin{itemize}
\tightlist
\item[$\square$]
  \textbf{Step 3: Correct the source manifest without inventing
  attribution}
\end{itemize}

\begin{lstlisting}
{
  "type": "code",
  "modid": "firstgearbank",
  "name": "First Gear Bank of the Old Roads",
  "authors": [
    "Unknown"
  ],
  "description": "A server-authoritative bank and economy for Vintage Story.",
  "version": "1.0.0",
  "dependencies": {
    "game": "1.22.7"
  }
}
\end{lstlisting}

Keep \passthrough{\lstinline!Unknown!} because neither the repository
nor the approved specification names a release author.

\begin{itemize}
\tightlist
\item[$\square$]
  \textbf{Step 4: Implement the unique nondestructive task graph}
\end{itemize}

Use this exact dependency graph:

\begin{lstlisting}
ValidateEnvironment -> ValidateMetadata -> ValidateJson -> Restore -> Build -> Test -> Default
                                                                   `-> Stage -> Package -> VerifyPackage -> Deploy
\end{lstlisting}

\passthrough{\lstinline!Default!} depends only on
\passthrough{\lstinline!Test!}. \passthrough{\lstinline!Stage!},
\passthrough{\lstinline!Package!},
\passthrough{\lstinline!VerifyPackage!}, and
\passthrough{\lstinline!Deploy!} run only when explicitly targeted.
\passthrough{\lstinline!BuildContext!} owns resolved
\passthrough{\lstinline!BuildPaths!}, configuration, optional bump
component, and optional deploy directory; it does not read a manifest in
its constructor. It resolves the repository through
\passthrough{\lstinline!RepositoryLocator!} from the current working
directory, with an optional validated
\passthrough{\lstinline!--repositoryRoot!} override for IDE invocation.

Use these target names exactly:

\begin{lstlisting}
[TaskName("Default")]
[IsDependentOn(typeof(TestTask))]
public sealed class DefaultTask : FrostingTask<BuildContext>
{
}

[TaskName("VerifyPackage")]
[IsDependentOn(typeof(PackageTask))]
public sealed class VerifyPackageTask : FrostingTask<BuildContext>
{
    public override void Run(BuildContext context)
    {
        var inspection = PackageArchiveInspector.Inspect(
            context.ReleaseArchive,
            context.PackageVersion);

        if (!inspection.IsValid)
        {
            throw new CakeException(string.Join(Environment.NewLine, inspection.Errors));
        }
    }
}
\end{lstlisting}

\begin{itemize}
\tightlist
\item[$\square$]
  \textbf{Step 5: Stage only the approved package root}
\end{itemize}

\passthrough{\lstinline!StageTask!} cleans only
\texttt{Artifacts/}\allowbreak{}\texttt{stage/}\allowbreak{}\texttt{firstgearbank},
publishes the adapter to
\texttt{Artifacts/}\allowbreak{}\texttt{publish}, copies only
\passthrough{\lstinline!FirstGearBank.dll!},
\passthrough{\lstinline!FirstGearBank.Core.dll!},
\passthrough{\lstinline!modinfo.json!},
\texttt{assets/}\allowbreak{}\texttt{**}, and optional
\passthrough{\lstinline!modicon.png!}, and writes the bumped version
only into the staged manifest. It rejects a publish output containing
either required DLL zero or multiple times.

\passthrough{\lstinline!PackageTask!} creates
\texttt{Releases/}\allowbreak{}\texttt{firstgearbank\_}\allowbreak{}\texttt{\textless{}}\allowbreak{}\texttt{version\textgreater{}}\allowbreak{}\texttt{.}\allowbreak{}\texttt{zip}
through \passthrough{\lstinline!PackageArchiveWriter!}; it never cleans
the complete \passthrough{\lstinline!Releases!} directory.

\begin{itemize}
\tightlist
\item[$\square$]
  \textbf{Step 6: Implement explicit hash-verified ZIP deployment}
\end{itemize}

Resolve the default only inside \passthrough{\lstinline!DeployTask!}:

\begin{lstlisting}
var appData = Environment.GetFolderPath(Environment.SpecialFolder.ApplicationData);
var defaultModsDirectory = Path.Combine(
    appData,
    "StoryForge",
    "installations",
    "working_test_world",
    "Mods");
\end{lstlisting}

Allow
\texttt{-}\allowbreak{}\texttt{-}\allowbreak{}\texttt{deployDirectory=}\allowbreak{}\texttt{\textless{}}\allowbreak{}\texttt{absolute\ }\allowbreak{}\texttt{path\textgreater{}}\allowbreak{}.
Reject a relative path, filesystem root, repository descendant, or any
build-owned source/\allowbreak{}stage/\allowbreak{}release directory.
\passthrough{\lstinline!DeployTask!} must obtain a successful
\texttt{DeploymentPathValidation.}\allowbreak{}\texttt{Directory} token
and pass that token---not the raw option string or
\passthrough{\lstinline!DirectoryInfo!}---to
\texttt{PackageDeploymentInstaller}; any invalid result returns before a
mutating file operation. Validate the source ZIP, copy it to a uniquely
named temporary file inside the resolved Mods directory, compare
SHA-256, and inspect only bounded root
\passthrough{\lstinline!modinfo.json!} entries when identifying existing
ZIPs. Move confirmed \passthrough{\lstinline!firstgearbank!} matches to
unique same-directory backup names, atomically rename the temporary
file, recheck SHA-256 and package validity, then remove only those
confirmed backups. If installation or verification fails, restore the
backups; if cleanup fails after success, retain the backup and warn.
Open File Explorer only after verified success. Warn and retain every
unreadable or foreign ZIP. Never recursively delete the Mods directory
or any unpacked directory.

\begin{itemize}
\tightlist
\item[$\square$]
  \textbf{Step 7: Simplify wrapper defaults and remove repository-wide
  staging}
\end{itemize}

Update usage text so no-argument invocation means
\passthrough{\lstinline!Default!}, package verification uses
\passthrough{\lstinline!--target=VerifyPackage!}, and deployment uses
\passthrough{\lstinline!--target=Deploy!}. Remove
\passthrough{\lstinline!--commit!}, Codex invocation, temporary prompt
files, and every \passthrough{\lstinline!git add -A!} path from both
wrappers. If a caller supplies the removed option, fail with a concise
migration message telling them to inspect and commit an explicit path
set themselves. The wrappers build; they never decide which repository
changes belong in a commit.

\begin{itemize}
\tightlist
\item[$\square$]
  \textbf{Step 8: Run focused build tests and all task smoke tests}
\end{itemize}

\begin{lstlisting}
dotnet test tests/CakeBuild.Tests/CakeBuild.Tests.csproj -c Release
dotnet run --project CakeBuild/CakeBuild.csproj -- --target=Default
dotnet run --project CakeBuild/CakeBuild.csproj -- --target=VerifyPackage
if (rg -n "git add -A|codex" build.py build.ps1) {
    throw "Build wrappers still contain repository-wide commit behavior."
}
$pythonCommitOutput = & python -B build.py --commit=test 2>&1
if ($LASTEXITCODE -eq 0 -or $pythonCommitOutput -notmatch "removed") {
    throw "build.py did not reject the removed --commit option."
}
$powershellCommitOutput = & pwsh -NoProfile -File .\build.ps1 --commit=test 2>&1
if ($LASTEXITCODE -eq 0 -or $powershellCommitOutput -notmatch "removed") {
    throw "build.ps1 did not reject the removed --commit option."
}
\end{lstlisting}

Expected: both targets succeed; \passthrough{\lstinline!Default!}
creates no release and performs no deployment;
\passthrough{\lstinline!VerifyPackage!} produces one validated ZIP with
both mod DLLs and no game/build DLL.

\begin{itemize}
\tightlist
\item[$\square$]
  \textbf{Step 9: Verify repeated packaging and commit}
\end{itemize}

\begin{lstlisting}
$first = (Get-FileHash -Algorithm SHA256 Releases/firstgearbank_1.0.0.zip).Hash
dotnet run --project CakeBuild/CakeBuild.csproj -- --target=VerifyPackage
$second = (Get-FileHash -Algorithm SHA256 Releases/firstgearbank_1.0.0.zip).Hash
if ($first -ne $second) { throw "Package is not reproducible." }
git diff --check
git add Directory.Packages.props CakeBuild "First Gear Bank/modinfo.json" build.ps1 build.py README.md `
  "First Gear Bank/packages.lock.json" FirstGearBank.Core/packages.lock.json `
  tests/FirstGearBank.Core.Tests/FirstGearBank.Core.Tests.csproj `
  tests/FirstGearBank.Core.Tests/packages.lock.json `
  tests/FirstGearBank.Adapter.Tests/FirstGearBank.Adapter.Tests.csproj `
  tests/FirstGearBank.Adapter.Tests/packages.lock.json tests/CakeBuild.Tests
git commit -m "build: make packaging reproducible and deployment explicit"
\end{lstlisting}

Expected: hashes are identical and
\passthrough{\lstinline!git status --short!} contains no release
artifact because build output is ignored.

\subsubsection{Task 5: Add Domain Identifiers, Fixed-Point Money, and
Financial
Instants}\label{task-5-add-domain-identifiers-fixed-point-money-and-financial-instants}

\textbf{Files:}

\begin{itemize}
\tightlist
\item
  Create:
  \texttt{FirstGearBank.}\allowbreak{}\texttt{Core/}\allowbreak{}\texttt{Domain/}\allowbreak{}\texttt{BankCurrency.}\allowbreak{}\texttt{cs}
\item
  Create:
  \texttt{FirstGearBank.}\allowbreak{}\texttt{Core/}\allowbreak{}\texttt{Domain/}\allowbreak{}\texttt{BankErrorCode.}\allowbreak{}\texttt{cs}
\item
  Create:
  \texttt{FirstGearBank.}\allowbreak{}\texttt{Core/}\allowbreak{}\texttt{Domain/}\allowbreak{}\texttt{BankFailure.}\allowbreak{}\texttt{cs}
\item
  Create:
  \texttt{FirstGearBank.}\allowbreak{}\texttt{Core/}\allowbreak{}\texttt{Domain/}\allowbreak{}\texttt{BankResult.}\allowbreak{}\texttt{cs}
\item
  Create:
  \texttt{FirstGearBank.}\allowbreak{}\texttt{Core/}\allowbreak{}\texttt{Domain/}\allowbreak{}\texttt{Identifiers.}\allowbreak{}\texttt{cs}
\item
  Create:
  \texttt{FirstGearBank.}\allowbreak{}\texttt{Core/}\allowbreak{}\texttt{Domain/}\allowbreak{}\texttt{BankMutationDomains.}\allowbreak{}\texttt{cs}
\item
  Create:
  \texttt{FirstGearBank.}\allowbreak{}\texttt{Core/}\allowbreak{}\texttt{Domain/}\allowbreak{}\texttt{BankRevisionVector.}\allowbreak{}\texttt{cs}
\item
  Create:
  \texttt{FirstGearBank.}\allowbreak{}\texttt{Core/}\allowbreak{}\texttt{Domain/}\allowbreak{}\texttt{BankInvocationKind.}\allowbreak{}\texttt{cs}
\item
  Create:
  \texttt{FirstGearBank.}\allowbreak{}\texttt{Core/}\allowbreak{}\texttt{Domain/}\allowbreak{}\texttt{BankInvocationKey.}\allowbreak{}\texttt{cs}
\item
  Create:
  \texttt{FirstGearBank.}\allowbreak{}\texttt{Core/}\allowbreak{}\texttt{Domain/}\allowbreak{}\texttt{BankInvocationIntent.}\allowbreak{}\texttt{cs}
\item
  Create:
  \texttt{FirstGearBank.}\allowbreak{}\texttt{Core/}\allowbreak{}\texttt{Domain/}\allowbreak{}\texttt{CommandPayloadDigest.}\allowbreak{}\texttt{cs}
\item
  Create:
  \texttt{FirstGearBank.}\allowbreak{}\texttt{Core/}\allowbreak{}\texttt{Domain/}\allowbreak{}\texttt{EvidenceDigest.}\allowbreak{}\texttt{cs}
\item
  Create:
  \texttt{FirstGearBank.}\allowbreak{}\texttt{Core/}\allowbreak{}\texttt{Money/}\allowbreak{}\texttt{BankUnits.}\allowbreak{}\texttt{cs}
\item
  Create:
  \texttt{FirstGearBank.}\allowbreak{}\texttt{Core/}\allowbreak{}\texttt{Money/}\allowbreak{}\texttt{MoneyLimits.}\allowbreak{}\texttt{cs}
\item
  Create:
  \texttt{FirstGearBank.}\allowbreak{}\texttt{Core/}\allowbreak{}\texttt{Money/}\allowbreak{}\texttt{MoneyMath.}\allowbreak{}\texttt{cs}
\item
  Create:
  \texttt{FirstGearBank.}\allowbreak{}\texttt{Core/}\allowbreak{}\texttt{Money/}\allowbreak{}\texttt{MoneyFormatter.}\allowbreak{}\texttt{cs}
\item
  Create:
  \texttt{FirstGearBank.}\allowbreak{}\texttt{Core/}\allowbreak{}\texttt{Time/}\allowbreak{}\texttt{FinancialInstant.}\allowbreak{}\texttt{cs}
\item
  Create:
  \texttt{FirstGearBank.}\allowbreak{}\texttt{Core/}\allowbreak{}\texttt{Time/}\allowbreak{}\texttt{FinancialTimestamp.}\allowbreak{}\texttt{cs}
\item
  Create:
  \texttt{FirstGearBank.}\allowbreak{}\texttt{Core/}\allowbreak{}\texttt{Diagnostics/}\allowbreak{}\texttt{BankLogLevel.}\allowbreak{}\texttt{cs}
\item
  Create:
  \texttt{FirstGearBank.}\allowbreak{}\texttt{Core/}\allowbreak{}\texttt{Diagnostics/}\allowbreak{}\texttt{IBankLog.}\allowbreak{}\texttt{cs}
\item
  Create:
  \texttt{FirstGearBank.}\allowbreak{}\texttt{Core/}\allowbreak{}\texttt{Security/}\allowbreak{}\texttt{IEntropySource.}\allowbreak{}\texttt{cs}
\item
  Create:
  \texttt{FirstGearBank.}\allowbreak{}\texttt{Core/}\allowbreak{}\texttt{Security/}\allowbreak{}\texttt{CryptographicEntropySource.}\allowbreak{}\texttt{cs}
\item
  Create:
  \texttt{tests/}\allowbreak{}\texttt{FirstGearBank.}\allowbreak{}\texttt{Core.}\allowbreak{}\texttt{Tests/}\allowbreak{}\texttt{Domain/}\allowbreak{}\texttt{IdentifierTests.}\allowbreak{}\texttt{cs}
\item
  Create:
  \texttt{tests/}\allowbreak{}\texttt{FirstGearBank.}\allowbreak{}\texttt{Core.}\allowbreak{}\texttt{Tests/}\allowbreak{}\texttt{Domain/}\allowbreak{}\texttt{BankRevisionVectorTests.}\allowbreak{}\texttt{cs}
\item
  Create:
  \texttt{tests/}\allowbreak{}\texttt{FirstGearBank.}\allowbreak{}\texttt{Core.}\allowbreak{}\texttt{Tests/}\allowbreak{}\texttt{Domain/}\allowbreak{}\texttt{BankInvocationKeyTests.}\allowbreak{}\texttt{cs}
\item
  Create:
  \texttt{tests/}\allowbreak{}\texttt{FirstGearBank.}\allowbreak{}\texttt{Core.}\allowbreak{}\texttt{Tests/}\allowbreak{}\texttt{Domain/}\allowbreak{}\texttt{CommandPayloadDigestTests.}\allowbreak{}\texttt{cs}
\item
  Create:
  \texttt{tests/}\allowbreak{}\texttt{FirstGearBank.}\allowbreak{}\texttt{Core.}\allowbreak{}\texttt{Tests/}\allowbreak{}\texttt{Money/}\allowbreak{}\texttt{BankUnitsTests.}\allowbreak{}\texttt{cs}
\item
  Create:
  \texttt{tests/}\allowbreak{}\texttt{FirstGearBank.}\allowbreak{}\texttt{Core.}\allowbreak{}\texttt{Tests/}\allowbreak{}\texttt{Money/}\allowbreak{}\texttt{MoneyMathTests.}\allowbreak{}\texttt{cs}
\item
  Create:
  \texttt{tests/}\allowbreak{}\texttt{FirstGearBank.}\allowbreak{}\texttt{Core.}\allowbreak{}\texttt{Tests/}\allowbreak{}\texttt{Money/}\allowbreak{}\texttt{MoneyFormatterTests.}\allowbreak{}\texttt{cs}
\item
  Create:
  \texttt{tests/}\allowbreak{}\texttt{FirstGearBank.}\allowbreak{}\texttt{Core.}\allowbreak{}\texttt{Tests/}\allowbreak{}\texttt{Time/}\allowbreak{}\texttt{FinancialInstantTests.}\allowbreak{}\texttt{cs}
\end{itemize}

\textbf{Interfaces:}

\begin{itemize}
\tightlist
\item
  Consumes: no engine type and no culture-sensitive global state.
\item
  Produces: the serialized value types used by journal, coordinator,
  settlement, persistence, and every later product.
\end{itemize}

\begin{itemize}
\tightlist
\item[$\square$]
  \textbf{Step 1: Write failing money conversion and overflow tests}
\end{itemize}

\begin{lstlisting}
[Theory]
[InlineData(0.5, 0)]
[InlineData(1.5, 2)]
[InlineData(-0.5, 0)]
[InlineData(-1.5, -2)]
public void Rate_result_conversion_rounds_midpoints_to_even(double exactBankUnits, long expectedUnits)
{
    var result = MoneyMath.RoundRateResultToUnits(exactBankUnits);

    Assert.True(result.IsSuccess);
    Assert.Equal(expectedUnits, result.Value.Value);
}

[Fact]
public void Checked_add_rejects_long_overflow()
{
    var result = new BankUnits(long.MaxValue).Add(new BankUnits(1));

    Assert.Equal(BankErrorCode.ArithmeticOverflow, result.Failure?.Code);
}
\end{lstlisting}

Add cases for exact one and quarter gear, negative values,
\passthrough{\lstinline!NaN!}, both infinities,
\passthrough{\lstinline!long.MinValue!} negation, exact-input
overprecision, and the customer cap
\passthrough{\lstinline!2\_147\_483\_647\_000\_000!}.

\begin{itemize}
\tightlist
\item[$\square$]
  \textbf{Step 2: Run the money tests and observe missing-type failures}
\end{itemize}

\begin{lstlisting}
dotnet test tests/FirstGearBank.Core.Tests/FirstGearBank.Core.Tests.csproj -c Release
\end{lstlisting}

Expected: FAIL because the money and result types do not exist.

\begin{itemize}
\tightlist
\item[$\square$]
  \textbf{Step 3: Implement typed results and checked fixed-point
  arithmetic}
\end{itemize}

Use these contracts:

\begin{lstlisting}
public enum BankCurrency : byte
{
    Rusty = 1,
    Temporal = 2
}

public sealed record BankFailure(BankErrorCode Code, string Diagnostic);

public sealed class BankResult<T>
{
    public bool IsSuccess { get; }
    public T Value { get; }
    public BankFailure? Failure { get; }

    public static BankResult<T> Success(T value);
    public static BankResult<T> Failed(BankErrorCode code, string diagnostic);
}

public readonly record struct BankUnits(long Value)
{
    public const long UnitsPerGear = 1_000_000;

    public BankResult<BankUnits> Add(BankUnits other);
    public BankResult<BankUnits> Subtract(BankUnits other);
    public BankResult<BankUnits> Negate();
}

public static class MoneyLimits
{
    public const long MaximumCustomerBalanceGears = int.MaxValue;
    public const long MaximumCustomerBalanceUnits = 2_147_483_647_000_000;
}

public static class MoneyMath
{
    public static BankResult<BankUnits> ParseExactGears(ReadOnlySpan<char> text);
    public static BankResult<BankUnits> RoundRateResultToUnits(double exactBankUnits);
    public static BankResult<BankUnits> ValidateCustomerBalance(BankUnits units);
}
\end{lstlisting}

Use invariant decimal parsing for player/configured exact amounts and
reject more than six meaningful fractional digits. Use
\passthrough{\lstinline!checked!} arithmetic and return typed failures
rather than wrapping. Floating-rate calculations pass their exact
bank-unit result to the central converter, which rejects
nonfinite/out-of-range input and applies
\passthrough{\lstinline!MidpointRounding.ToEven!} once.

Freeze the foundation values explicitly; later phases may append but
never renumber them:

\begin{lstlisting}
public enum BankErrorCode : ushort
{
    None = 0,
    InvalidAmount = 1,
    AmountOverprecision = 2,
    NonFiniteValue = 3,
    ArithmeticOverflow = 4,
    InsufficientFunds = 5,
    CustomerBalanceCapExceeded = 6,
    NegativeCustomerBalance = 7,
    InvalidFinancialInstant = 8,
    AccountNotFound = 9,
    InvalidConfiguration = 10,
    MalformedLedgerAccount = 11,
    UnbalancedJournal = 12,
    InvalidJournalSequence = 13,
    DuplicateOperation = 14,
    JournalIntegrityFailure = 15,
    InactiveRequestScope = 16,
    UnexpectedRequestSequence = 17,
    RequestPayloadMismatch = 18,
    AlreadyProcessedResponseExpired = 19,
    DomainUnavailable = 20,
    CorruptState = 21,
    UnsupportedSchema = 22,
    InvalidSettlementEvidence = 23,
    PersistenceFailure = 24,
    UndeclaredMutationDomain = 25,
    IdentifierContractViolation = 26
}
\end{lstlisting}

Player-facing text is not stored in this enum; later localization maps
codes to prose.

\begin{itemize}
\tightlist
\item[$\square$]
  \textbf{Step 4: Run the money tests and verify they pass}
\end{itemize}

\begin{lstlisting}
dotnet test tests/FirstGearBank.Core.Tests/FirstGearBank.Core.Tests.csproj -c Release
\end{lstlisting}

Expected: PASS with no culture-dependent result under
\passthrough{\lstinline!en-US!}, \passthrough{\lstinline!de-DE!}, and
invariant test cultures.

\begin{itemize}
\tightlist
\item[$\square$]
  \textbf{Step 5: Write failing formatting and financial-instant tests}
\end{itemize}

\begin{lstlisting}
[Theory]
[InlineData(0, "1")]
[InlineData(3, "1.235")]
[InlineData(6, "1.234568")]
public void Rusty_display_rounds_without_changing_ledger_units(int precision, string expected)
{
    var units = new BankUnits(1_234_567);

    Assert.Equal(expected, MoneyFormatter.FormatRusty(units, precision));
    Assert.Equal(1_234_567, units.Value);
}

[Fact]
public void Whole_month_tenor_addition_is_exact()
{
    var start = FinancialInstant.Create(41, 0.375m).Value;

    Assert.Equal(
        FinancialInstant.Create(53, 0.375m).Value,
        start.AddWholeMonths(12).Value);
}
\end{lstlisting}

Add precision bounds \passthrough{\lstinline!0..6!}, negative display,
temporal trailing-zero trimming, financial fraction bounds, ordering,
overflow, invariant \passthrough{\lstinline!G29!} fraction
serialization, and conversion to financial years only at the math
boundary.

\begin{lstlisting}
dotnet test tests/FirstGearBank.Core.Tests/FirstGearBank.Core.Tests.csproj -c Release
\end{lstlisting}

Expected: FAIL because formatting and normalized financial-time behavior
are not implemented.

\begin{itemize}
\tightlist
\item[$\square$]
  \textbf{Step 6: Implement presentation-only formatting and normalized
  time}
\end{itemize}

\begin{lstlisting}
public readonly record struct FinancialInstant : IComparable<FinancialInstant>
{
    public long CompletedMonths { get; }
    public decimal FractionOfMonth { get; }

    public static BankResult<FinancialInstant> Create(
        long completedMonths,
        decimal fractionOfMonth);

    public BankResult<FinancialInstant> AddWholeMonths(int months);
    public double ToFinancialYears();
    public int CompareTo(FinancialInstant other);
}

public readonly record struct FinancialTimestamp
{
    public FinancialInstant Financial { get; }
    public double WorldCalendarTotalDays { get; }

    public static BankResult<FinancialTimestamp> Create(
        FinancialInstant financial,
        double worldCalendarTotalDays);
}
\end{lstlisting}

\passthrough{\lstinline!Create!} accepts only
\passthrough{\lstinline!completedMonths >= 0!} and
\passthrough{\lstinline!0 <= fraction < 1!}.
\passthrough{\lstinline!FinancialTimestamp!} construction rejects a
nonfinite world-calendar value through a factory. The formatter uses
decimal division and ties-to-even presentation rounding but never
returns a modified \passthrough{\lstinline!BankUnits!} value.

\begin{itemize}
\tightlist
\item[$\square$]
  \textbf{Step 7: Write failing identifier and revision tests}
\end{itemize}

Test all identifier types named below for 16-byte binary roundtrip and
lowercase 32-hex diagnostic format. Test that a
\passthrough{\lstinline!PlayerId!} preserves exact case/content, rejects
empty and overlong UTF-8 input, and has no public display conversion.
Test \passthrough{\lstinline!CommandPayloadDigest!} for exact 32-byte
construction, defensive copying, structural equality, fixed-time
comparison, and command-kind/\allowbreak{}schema/\allowbreak{}payload
domain separation. Freeze
\texttt{NoticeId.}\allowbreak{}\texttt{ForAdminCorrection} with
operation bytes \texttt{000102030405060708090a0b0c0d0e0f}: SHA-256 over
the ASCII domain
\texttt{FGB/}\allowbreak{}\texttt{ADMIN-}\allowbreak{}\texttt{CORRECTION-}\allowbreak{}\texttt{NOTICE/}\allowbreak{}\texttt{1\textbackslash{}}\allowbreak{}\texttt{0}
followed by those canonical operation bytes yields the first-16-byte
notice ID \texttt{8c887b94a1b027db46e160bf54205382}. A different
operation or derivation domain must differ. Test
\passthrough{\lstinline!EvidenceDigest!} as an exact structural 32-byte
SHA-256 wrapper with defensive copies and fixed-time comparison.

\begin{lstlisting}
[Fact]
public void Only_changed_subsystem_revisions_advance()
{
    var before = BankRevisionVector.Initial;

    var after = before.Advance(BankMutationDomains.Finance | BankMutationDomains.Requests);

    Assert.Equal(1, after.Global.Value);
    Assert.Equal(1, after.Finance.Value);
    Assert.Equal(1, after.Requests.Value);
    Assert.Equal(0, after.Registry.Value);
    Assert.Equal(0, after.Branches.Value);
    Assert.Equal(0, after.Recovery.Value);
}
\end{lstlisting}

\begin{lstlisting}
dotnet test tests/FirstGearBank.Core.Tests/FirstGearBank.Core.Tests.csproj -c Release
\end{lstlisting}

Expected: FAIL because the canonical identifier and revision contracts
do not exist.

\begin{itemize}
\tightlist
\item[$\square$]
  \textbf{Step 8: Implement canonical identifiers and versioned stable
  hashes}
\end{itemize}

Create distinct record structs for:

\begin{lstlisting}
WorldInstanceId, PlayerId, OperationId, CommandId, RequestScopeNonce, RequestKey,
CdId, SettlementId, NoticeId, BranchId, RegistryEpochId, AutomaticOperationKey
\end{lstlisting}

GUID-backed IDs serialize as 16 canonical RFC-4122-order bytes in
persistence and lowercase \passthrough{\lstinline!N!} text in
diagnostics. \passthrough{\lstinline!PlayerId!} stores the exact
authenticated engine UID with a 256-byte UTF-8 bound.
\passthrough{\lstinline!RequestKey!} contains player, scope, and a
strictly positive signed 64-bit sequence.
\passthrough{\lstinline!AutomaticOperationKey!} contains a bounded ASCII
domain/version and stable canonical identifier; it never calls
process-randomized \passthrough{\lstinline!string.GetHashCode!}.
\texttt{NoticeId.}\allowbreak{}\texttt{ForAdminCorrection} performs the
exact domain-separated SHA-256 derivation tested above and copies the
first 16 digest bytes without GUID version/variant bit rewriting.
\passthrough{\lstinline!EvidenceDigest!} retains all 32 SHA-256 bytes
and never uses immutable-array reference equality.

Every journal-producing coordinator invocation carries exactly one typed
key:

\begin{lstlisting}
public enum BankInvocationKind : byte
{
    AuthenticatedRequest = 1,
    AutomaticOperation = 2
}

public readonly record struct BankInvocationKey
{
    public BankInvocationKind Kind { get; }
    public RequestKey? Request { get; }
    public AutomaticOperationKey? Automatic { get; }

    public static BankInvocationKey ForRequest(RequestKey request);
    public static BankInvocationKey ForAutomatic(AutomaticOperationKey automatic);
}

public enum BankInvocationIntentKind : byte
{
    Caller = 1,
    AutomaticOperation = 2
}

public readonly record struct BankInvocationIntent
{
    public BankInvocationIntentKind Kind { get; }
    public AutomaticOperationKey? AutomaticOperation { get; }
    public CommandPayloadDigest? AutomaticPayloadDigest { get; }

    public static BankInvocationIntent Caller { get; }
    public static BankInvocationIntent ForAutomatic(
        AutomaticOperationKey operation,
        CommandPayloadDigest payloadDigest);
}
\end{lstlisting}

The factories enforce one-and-only-one representation. Canonical
serialization emits the kind byte followed only by that
variant\textquotesingle s bytes.
\passthrough{\lstinline!CommandPayloadDigest!} is the structural
32-byte, domain/schema-separated SHA-256 value used by both request and
automatic intents. It hashes
\texttt{FGB/}\allowbreak{}\texttt{COMMAND-}\allowbreak{}\texttt{PAYLOAD/}\allowbreak{}\texttt{1\textbackslash{}}\allowbreak{}\texttt{0},
bounded command-kind UTF-8 bytes, schema version, and bounded canonical
payload bytes using the framing frozen in Task 9. Tests reject empty,
dual, or mismatched representations and freeze both variant encodings.

Use one revision vector:

\begin{lstlisting}
[Flags]
public enum BankMutationDomains : ushort
{
    None = 0,
    Finance = 1 << 0,
    Requests = 1 << 1,
    Registry = 1 << 2,
    Delivery = 1 << 3,
    Settlements = 1 << 4,
    Branches = 1 << 5,
    Configuration = 1 << 6,
    Recovery = 1 << 7
}

public sealed record BankRevisionVector(
    BankRevision Global,
    SubsystemRevision Finance,
    SubsystemRevision Requests,
    SubsystemRevision Registry,
    SubsystemRevision Delivery,
    SubsystemRevision Settlements,
    SubsystemRevision Branches,
    SubsystemRevision Configuration,
    SubsystemRevision Recovery)
{
    public static BankRevisionVector Initial { get; }
    public BankRevisionVector Advance(BankMutationDomains changedDomains);
}
\end{lstlisting}

Define the production entropy boundary at the same time:

\begin{lstlisting}
public interface IEntropySource
{
    void Fill(Span<byte> destination);
}

public sealed class CryptographicEntropySource : IEntropySource
{
    public void Fill(Span<byte> destination);
}
\end{lstlisting}

\texttt{CryptographicEntropySource} fills caller-owned spans through
\texttt{RandomNumberGenerator.}\allowbreak{}\texttt{Fill} with no
intermediate allocation. Deterministic entropy sources remain test-only;
production wiring never accepts a seed.

Define the caller-owned logging boundary at the same time:

\begin{lstlisting}
public enum BankLogLevel : byte
{
    Info = 1,
    Debug = 2,
    Warning = 3,
    Critical = 4
}

public interface IBankLog
{
    void Write(
        BankLogLevel level,
        string component,
        string message,
        Exception? exception = null);
}
\end{lstlisting}

The core emits through this interface only; it never creates a file or
calls an engine logger.

\begin{itemize}
\tightlist
\item[$\square$]
  \textbf{Step 9: Verify all primitive tests and commit}
\end{itemize}

\begin{lstlisting}
dotnet test tests/FirstGearBank.Core.Tests/FirstGearBank.Core.Tests.csproj -c Release
git diff --check
git add FirstGearBank.Core/Domain FirstGearBank.Core/Money FirstGearBank.Core/Time `
  FirstGearBank.Core/Diagnostics FirstGearBank.Core/Security tests/FirstGearBank.Core.Tests/Domain `
  tests/FirstGearBank.Core.Tests/Money tests/FirstGearBank.Core.Tests/Time
git commit -m "feat: add fixed-point banking primitives"
\end{lstlisting}

\subsubsection{Task 6: Parse and Validate the Complete Static JSONC
Configuration}\label{task-6-parse-and-validate-the-complete-static-jsonc-configuration}

\textbf{Files:}

\begin{itemize}
\tightlist
\item
  Create:
  \texttt{FirstGearBank.}\allowbreak{}\texttt{Core/}\allowbreak{}\texttt{Configuration/}\allowbreak{}\texttt{InterestTimeBasis.}\allowbreak{}\texttt{cs}
\item
  Create:
  \texttt{FirstGearBank.}\allowbreak{}\texttt{Core/}\allowbreak{}\texttt{Configuration/}\allowbreak{}\texttt{RiskFreeRateModel.}\allowbreak{}\texttt{cs}
\item
  Create:
  \texttt{FirstGearBank.}\allowbreak{}\texttt{Core/}\allowbreak{}\texttt{Configuration/}\allowbreak{}\texttt{RecipientSelectionMode.}\allowbreak{}\texttt{cs}
\item
  Create:
  \texttt{FirstGearBank.}\allowbreak{}\texttt{Core/}\allowbreak{}\texttt{Configuration/}\allowbreak{}\texttt{BankConfiguration.}\allowbreak{}\texttt{cs}
\item
  Create:
  \texttt{FirstGearBank.}\allowbreak{}\texttt{Core/}\allowbreak{}\texttt{Configuration/}\allowbreak{}\texttt{BankConfigurationDefaults.}\allowbreak{}\texttt{cs}
\item
  Create:
  \texttt{FirstGearBank.}\allowbreak{}\texttt{Core/}\allowbreak{}\texttt{Configuration/}\allowbreak{}\texttt{ConfigurationDiagnostic.}\allowbreak{}\texttt{cs}
\item
  Create:
  \texttt{FirstGearBank.}\allowbreak{}\texttt{Core/}\allowbreak{}\texttt{Configuration/}\allowbreak{}\texttt{ConfigurationParseResult.}\allowbreak{}\texttt{cs}
\item
  Create:
  \texttt{FirstGearBank.}\allowbreak{}\texttt{Core/}\allowbreak{}\texttt{Configuration/}\allowbreak{}\texttt{ValidatedEconomicConfiguration.}\allowbreak{}\texttt{cs}
\item
  Create:
  \texttt{FirstGearBank.}\allowbreak{}\texttt{Core/}\allowbreak{}\texttt{Configuration/}\allowbreak{}\texttt{BankConfigurationJsonc.}\allowbreak{}\texttt{cs}
\item
  Create:
  \texttt{FirstGearBank.}\allowbreak{}\texttt{Core/}\allowbreak{}\texttt{Configuration/}\allowbreak{}\texttt{JsoncSyntaxDocument.}\allowbreak{}\texttt{cs}
\item
  Create:
  \texttt{FirstGearBank.}\allowbreak{}\texttt{Core/}\allowbreak{}\texttt{Configuration/}\allowbreak{}\texttt{IEconomicConfigurationGate.}\allowbreak{}\texttt{cs}
\item
  Create:
  \texttt{FirstGearBank.}\allowbreak{}\texttt{Core/}\allowbreak{}\texttt{Configuration/}\allowbreak{}\texttt{IBankConfigurationStore.}\allowbreak{}\texttt{cs}
\item
  Create:
  \texttt{First\ }\allowbreak{}\texttt{Gear\ }\allowbreak{}\texttt{Bank/}\allowbreak{}\texttt{Configuration/}\allowbreak{}\texttt{firstgearbank.}\allowbreak{}\texttt{jsonc.}\allowbreak{}\texttt{template}
\item
  Create:
  \texttt{tests/}\allowbreak{}\texttt{FirstGearBank.}\allowbreak{}\texttt{Core.}\allowbreak{}\texttt{Tests/}\allowbreak{}\texttt{Configuration/}\allowbreak{}\texttt{ConfigurationDefaultsTests.}\allowbreak{}\texttt{cs}
\item
  Create:
  \texttt{tests/}\allowbreak{}\texttt{FirstGearBank.}\allowbreak{}\texttt{Core.}\allowbreak{}\texttt{Tests/}\allowbreak{}\texttt{Configuration/}\allowbreak{}\texttt{ConfigurationDomainTests.}\allowbreak{}\texttt{cs}
\item
  Create:
  \texttt{tests/}\allowbreak{}\texttt{FirstGearBank.}\allowbreak{}\texttt{Core.}\allowbreak{}\texttt{Tests/}\allowbreak{}\texttt{Configuration/}\allowbreak{}\texttt{JsoncRoundTripTests.}\allowbreak{}\texttt{cs}
\item
  Create:
  \texttt{tests/}\allowbreak{}\texttt{FirstGearBank.}\allowbreak{}\texttt{Core.}\allowbreak{}\texttt{Tests/}\allowbreak{}\texttt{Configuration/}\allowbreak{}\texttt{LegacyConfigurationMigrationTests.}\allowbreak{}\texttt{cs}
\item
  Modify:
  \texttt{First\ }\allowbreak{}\texttt{Gear\ }\allowbreak{}\texttt{Bank/}\allowbreak{}\texttt{First\ }\allowbreak{}\texttt{Gear\ }\allowbreak{}\texttt{Bank.}\allowbreak{}\texttt{csproj:}\allowbreak{}\texttt{1}
\end{itemize}

\textbf{Interfaces:}

\begin{itemize}
\tightlist
\item
  Consumes: exact money parsing and the enum/record identifiers from
  Task 5.
\item
  Produces: a fully typed configuration candidate, field-specific
  diagnostics, a lossless JSONC syntax document, and the
  derived-economic-validation gate consumed by the financial-core plan.
\end{itemize}

\begin{itemize}
\tightlist
\item[$\square$]
  \textbf{Step 1: Write failing locked-default tests}
\end{itemize}

\begin{lstlisting}
[Fact]
public void Defaults_match_the_locked_v03_configuration()
{
    var value = BankConfigurationDefaults.Value;

    Assert.Equal(InterestTimeBasis.InGame, value.InterestTimeBasis);
    Assert.Equal(RiskFreeRateModel.Cir, value.RiskFreeRateModel);
    Assert.Equal(3.0, value.TargetAnnualizedRiskFreeRate);
    Assert.Equal(1.0, value.Cir.MeanReversionSpeed);
    Assert.Equal(0.12, value.Cir.Volatility);
    Assert.Equal(3, value.Currency.RustyGearDisplayPrecision);
    Assert.Equal(new BankUnits(250_000), value.Cd.MinimumRustyGearPrincipal);
    Assert.Equal([1, 3, 6, 12], value.Cd.TenorsMonths);
    Assert.Equal(RecipientSelectionMode.ExactName, value.RecipientSelectionMode);
    Assert.Equal(1.0, value.TransferCooldownSeconds);
    Assert.Equal(2.0, value.Charter.ArrivalMinimumDays);
    Assert.Equal(5.0, value.Charter.ArrivalMaximumDays);
    Assert.Equal(3.0, value.BankerReplacement.MinimumDays);
    Assert.Equal(7.0, value.BankerReplacement.MaximumDays);
    Assert.Equal(32, value.MinimumCharterBranchSpacingBlocks);
    Assert.Equal(0.15, value.NaturalBranches.TraderCompanionProbability);
    Assert.False(value.NaturalBranches.BackfillExistingTraderLocations);
    Assert.Equal(10, value.Statements.RecentTransactionsOnPrintedStatement);
}
\end{lstlisting}

Also assert base-spread
\texttt{0.}\allowbreak{}\texttt{30/}\allowbreak{}\texttt{0.}\allowbreak{}\texttt{395},
level shock
\texttt{gaussian/}\allowbreak{}\texttt{0/}\allowbreak{}\texttt{0.}\allowbreak{}\texttt{05},
slope shock
\texttt{gaussian/}\allowbreak{}\texttt{0/}\allowbreak{}\texttt{0.}\allowbreak{}\texttt{04},
and liquidity
\texttt{0.}\allowbreak{}\texttt{25/}\allowbreak{}\texttt{0.}\allowbreak{}\texttt{15/}\allowbreak{}\texttt{0.}\allowbreak{}\texttt{10/}\allowbreak{}\texttt{1.}\allowbreak{}\texttt{0/}\allowbreak{}\texttt{0.}\allowbreak{}\texttt{10}
in the same test file.

\begin{itemize}
\tightlist
\item[$\square$]
  \textbf{Step 2: Run the default tests and observe missing-type
  failures}
\end{itemize}

\begin{lstlisting}
dotnet test tests/FirstGearBank.Core.Tests/FirstGearBank.Core.Tests.csproj -c Release
\end{lstlisting}

Expected: FAIL because the configuration records and defaults do not
exist.

\begin{itemize}
\tightlist
\item[$\square$]
  \textbf{Step 3: Implement immutable records for every section 18.1
  field}
\end{itemize}

Use these top-level contracts and nested records with exactly the
representative field names:

\begin{lstlisting}
public sealed record BankConfiguration(
    InterestTimeBasis InterestTimeBasis,
    RiskFreeRateModel RiskFreeRateModel,
    double TargetAnnualizedRiskFreeRate,
    CirConfiguration Cir,
    CurrencyConfiguration Currency,
    CdConfiguration Cd,
    RecipientSelectionMode RecipientSelectionMode,
    double TransferCooldownSeconds,
    CharterConfiguration Charter,
    BankerReplacementConfiguration BankerReplacement,
    int MinimumCharterBranchSpacingBlocks,
    NaturalBranchesConfiguration NaturalBranches,
    StatementsConfiguration Statements);

public static class BankConfigurationDefaults
{
    public static BankConfiguration Value { get; }
}
\end{lstlisting}

Copy every default literally from section 18.1. Store the CD minimum as
\passthrough{\lstinline!BankUnits!} after exact decimal conversion and
store tenor months as an immutable array.

\begin{itemize}
\tightlist
\item[$\square$]
  \textbf{Step 4: Verify defaults, then write failing field-domain
  tests}
\end{itemize}

\begin{lstlisting}
dotnet test tests/FirstGearBank.Core.Tests/FirstGearBank.Core.Tests.csproj -c Release
\end{lstlisting}

Add one theory row at each accepted/rejected boundary from section 18:
enum case sensitivity; target \passthrough{\lstinline!q!}; CIR
\passthrough{\lstinline!kappa!} and \passthrough{\lstinline!sigma!};
Gaussian distribution/zero average/nonnegative variation; spread
scale/exponent; funding ratio; adjustment bounds; strictly positive
width/half-life; exact positive minimum principal; unique positive
tenors; probability; precision; cooldown; spacing; ordered positive
timers; and printed count \passthrough{\lstinline!0..100!}.

\begin{lstlisting}
dotnet test tests/FirstGearBank.Core.Tests/FirstGearBank.Core.Tests.csproj -c Release
\end{lstlisting}

Expected: FAIL on the first unimplemented field-domain validation.

\begin{itemize}
\tightlist
\item[$\square$]
  \textbf{Step 5: Implement field-local parsing and diagnostics}
\end{itemize}

\begin{lstlisting}
public sealed record ConfigurationDiagnostic(
    string JsonPath,
    ConfigurationDiagnosticCode Code,
    string Message);

public sealed record ConfigurationParseResult(
    BankConfiguration Candidate,
    JsoncSyntaxDocument Source,
    ImmutableArray<ConfigurationDiagnostic> Diagnostics,
    bool UsesLegacyStochasticSpreadKey);

public static class BankConfigurationJsonc
{
    public static ConfigurationParseResult Parse(string jsonc);
}
\end{lstlisting}

Parse with comments and trailing commas enabled, but read every property
explicitly so exact enum spellings remain closed. An invalid field emits
one diagnostic containing its exact JSON path and uses only that
field\textquotesingle s documented default. Unknown properties are
retained in the syntax document and reported once as
\passthrough{\lstinline!UnknownProperty!}; they are not executed or
silently copied into typed state. A malformed document returns the
complete documented defaults and one root diagnostic.

\begin{itemize}
\tightlist
\item[$\square$]
  \textbf{Step 6: Run all domain-boundary tests}
\end{itemize}

\begin{lstlisting}
dotnet test tests/FirstGearBank.Core.Tests/FirstGearBank.Core.Tests.csproj -c Release
\end{lstlisting}

Expected: PASS; unrelated valid fields survive every single-field
invalid fixture.

\begin{itemize}
\tightlist
\item[$\square$]
  \textbf{Step 7: Write failing lossless JSONC ownership tests}
\end{itemize}

\begin{lstlisting}
[Fact]
public void Reading_valid_comment_free_json_does_not_request_a_rewrite()
{
    var result = BankConfigurationJsonc.Parse("{\"InterestTimeBasis\":\"InGame\"}");

    Assert.False(result.Source.RequiresOwnedWrite);
}

[Fact]
public void Owned_scalar_edit_preserves_comments_unknown_properties_and_line_endings()
{
    const string source = "{\r\n  // operator note\r\n  \"TransferCooldownSeconds\": 1.0,\r\n  \"OperatorTag\": \"keep\"\r\n}\r\n";
    var document = BankConfigurationJsonc.Parse(source).Source;

    var edited = document.ReplaceScalar("$.TransferCooldownSeconds", "2.5");

    Assert.Contains("// operator note\r\n", edited);
    Assert.Contains("\"OperatorTag\": \"keep\"", edited);
    Assert.Contains("\"TransferCooldownSeconds\": 2.5", edited);
}
\end{lstlisting}

Add tests for LF, CRLF, trailing commas, escaped strings, comment
markers inside strings, comments before/after a value, and no rewrite on
ordinary successful load. Also write the full legacy
\passthrough{\lstinline!StochasticCDSpread!} migration/idempotence
fixture described in Step 9 now, before
\texttt{MigrateLegacyStochasticSpread} exists.

\begin{lstlisting}
dotnet test tests/FirstGearBank.Core.Tests/FirstGearBank.Core.Tests.csproj -c Release
\end{lstlisting}

Expected: FAIL because lossless scalar editing and JSONC ownership
behavior are absent.

\begin{itemize}
\tightlist
\item[$\square$]
  \textbf{Step 8: Implement a trivia-preserving JSONC syntax document}
\end{itemize}

\begin{lstlisting}
public sealed class JsoncSyntaxDocument
{
    public string OriginalText { get; }
    public bool RequiresOwnedWrite { get; }

    public string ReplaceScalar(string jsonPath, string canonicalJsonToken);
    public string MigrateLegacyStochasticSpread();
}
\end{lstlisting}

Tokenize strings, punctuation, whitespace, line comments, and block
comments without discarding byte spans. Edits replace only the targeted
token span. For the historical root
\passthrough{\lstinline!StochasticCDSpread!}, move its complete
value/trivia span into the existing \passthrough{\lstinline!CD!} object
under \passthrough{\lstinline!StochasticSpread!}, remove only the legacy
property and necessary comma, and preserve all other source bytes. If
both legacy and canonical keys exist, canonical wins and the legacy key
is retained with a conflict diagnostic rather than guessing.

\begin{itemize}
\tightlist
\item[$\square$]
  \textbf{Step 9: Prove legacy migration and comment preservation}
\end{itemize}

Use a full fixture with a root
\passthrough{\lstinline!StochasticCDSpread!} object, an existing
\passthrough{\lstinline!CD!} object, comments inside the stochastic
object, unknown root/CD properties, and CRLF endings. Parse the migrated
text again and assert one canonical
\passthrough{\lstinline!CD.StochasticSpread!}, no legacy key, unchanged
comments/unknown properties, and byte-identical output on a second
migration.

\begin{lstlisting}
dotnet test tests/FirstGearBank.Core.Tests/FirstGearBank.Core.Tests.csproj -c Release
\end{lstlisting}

\begin{itemize}
\tightlist
\item[$\square$]
  \textbf{Step 10: Define the derived economic activation boundary}
\end{itemize}

\begin{lstlisting}
public interface IEconomicConfigurationGate
{
    BankResult<ValidatedEconomicConfiguration> Validate(
        BankConfiguration candidate,
        double carriedShortRate);
}
\end{lstlisting}

Phase 1 may parse and stage an economic candidate but does not activate
it for pricing. The financial-core plan supplies the curve/payoff
implementation that checks every tenor, shock extreme, liquidity bound,
\passthrough{\lstinline!r=theta!}, and carried rate before advancing the
configuration revision. The canonical defaults remain the foundation
state until that gate succeeds.

Expose \texttt{IBankConfigurationStore.}\allowbreak{}\texttt{Load} and
\passthrough{\lstinline!SaveOwnedEdits!} from core so the server host
depends on a narrow file-operation port rather than a Vintage Story or
filesystem implementation.

\begin{itemize}
\tightlist
\item[$\square$]
  \textbf{Step 11: Add the canonical commented template and commit}
\end{itemize}

Copy the complete section 18.1 JSONC block into
\texttt{firstgearbank.}\allowbreak{}\texttt{jsonc.}\allowbreak{}\texttt{template},
embed it in the adapter assembly, and assert in
\texttt{ConfigurationDefaultsTests} that parsing it equals
\texttt{BankConfigurationDefaults.}\allowbreak{}\texttt{Value} with no
diagnostic.

\begin{lstlisting}
dotnet test tests/FirstGearBank.Core.Tests/FirstGearBank.Core.Tests.csproj -c Release
git diff --check
git add FirstGearBank.Core/Configuration "First Gear Bank/Configuration" `
  "First Gear Bank/First Gear Bank.csproj" tests/FirstGearBank.Core.Tests/Configuration
git commit -m "feat: add lossless validated bank configuration"
\end{lstlisting}

\subsubsection{Task 7: Add the Typed, Independently Balanced
Journal}\label{task-7-add-the-typed-independently-balanced-journal}

\textbf{Files:}

\begin{itemize}
\tightlist
\item
  Create:
  \texttt{FirstGearBank.}\allowbreak{}\texttt{Core/}\allowbreak{}\texttt{Ledger/}\allowbreak{}\texttt{LedgerAccountKind.}\allowbreak{}\texttt{cs}
\item
  Create:
  \texttt{FirstGearBank.}\allowbreak{}\texttt{Core/}\allowbreak{}\texttt{Ledger/}\allowbreak{}\texttt{LedgerAccountId.}\allowbreak{}\texttt{cs}
\item
  Create:
  \texttt{FirstGearBank.}\allowbreak{}\texttt{Core/}\allowbreak{}\texttt{Ledger/}\allowbreak{}\texttt{TransactionType.}\allowbreak{}\texttt{cs}
\item
  Create:
  \texttt{FirstGearBank.}\allowbreak{}\texttt{Core/}\allowbreak{}\texttt{Ledger/}\allowbreak{}\texttt{Posting.}\allowbreak{}\texttt{cs}
\item
  Create:
  \texttt{FirstGearBank.}\allowbreak{}\texttt{Core/}\allowbreak{}\texttt{Ledger/}\allowbreak{}\texttt{JournalPartySnapshot.}\allowbreak{}\texttt{cs}
\item
  Create:
  \texttt{FirstGearBank.}\allowbreak{}\texttt{Core/}\allowbreak{}\texttt{Ledger/}\allowbreak{}\texttt{IJournalPayload.}\allowbreak{}\texttt{cs}
\item
  Create:
  \texttt{FirstGearBank.}\allowbreak{}\texttt{Core/}\allowbreak{}\texttt{Ledger/}\allowbreak{}\texttt{Payloads/}\allowbreak{}\texttt{PhysicalDepositPayload.}\allowbreak{}\texttt{cs}
\item
  Create:
  \texttt{FirstGearBank.}\allowbreak{}\texttt{Core/}\allowbreak{}\texttt{Ledger/}\allowbreak{}\texttt{Payloads/}\allowbreak{}\texttt{PhysicalWithdrawalPayload.}\allowbreak{}\texttt{cs}
\item
  Create:
  \texttt{FirstGearBank.}\allowbreak{}\texttt{Core/}\allowbreak{}\texttt{Ledger/}\allowbreak{}\texttt{Payloads/}\allowbreak{}\texttt{SavingsInterestPayload.}\allowbreak{}\texttt{cs}
\item
  Create:
  \texttt{FirstGearBank.}\allowbreak{}\texttt{Core/}\allowbreak{}\texttt{Ledger/}\allowbreak{}\texttt{Payloads/}\allowbreak{}\texttt{TemporalGrossInterestPayload.}\allowbreak{}\texttt{cs}
\item
  Create:
  \texttt{FirstGearBank.}\allowbreak{}\texttt{Core/}\allowbreak{}\texttt{Ledger/}\allowbreak{}\texttt{Payloads/}\allowbreak{}\texttt{TemporalStoragePayload.}\allowbreak{}\texttt{cs}
\item
  Create:
  \texttt{FirstGearBank.}\allowbreak{}\texttt{Core/}\allowbreak{}\texttt{Ledger/}\allowbreak{}\texttt{Payloads/}\allowbreak{}\texttt{RustyTransferPayload.}\allowbreak{}\texttt{cs}
\item
  Create:
  \texttt{FirstGearBank.}\allowbreak{}\texttt{Core/}\allowbreak{}\texttt{Ledger/}\allowbreak{}\texttt{Payloads/}\allowbreak{}\texttt{AdminCorrectionPayload.}\allowbreak{}\texttt{cs}
\item
  Create:
  \texttt{FirstGearBank.}\allowbreak{}\texttt{Core/}\allowbreak{}\texttt{Ledger/}\allowbreak{}\texttt{JournalDraft.}\allowbreak{}\texttt{cs}
\item
  Create:
  \texttt{FirstGearBank.}\allowbreak{}\texttt{Core/}\allowbreak{}\texttt{Ledger/}\allowbreak{}\texttt{JournalRecord.}\allowbreak{}\texttt{cs}
\item
  Create:
  \texttt{FirstGearBank.}\allowbreak{}\texttt{Core/}\allowbreak{}\texttt{Ledger/}\allowbreak{}\texttt{JournalRecordFactory.}\allowbreak{}\texttt{cs}
\item
  Create:
  \texttt{FirstGearBank.}\allowbreak{}\texttt{Core/}\allowbreak{}\texttt{Ledger/}\allowbreak{}\texttt{JournalCanonicalEncoder.}\allowbreak{}\texttt{cs}
\item
  Create:
  \texttt{FirstGearBank.}\allowbreak{}\texttt{Core/}\allowbreak{}\texttt{Ledger/}\allowbreak{}\texttt{JournalValidationAnchor.}\allowbreak{}\texttt{cs}
\item
  Create:
  \texttt{FirstGearBank.}\allowbreak{}\texttt{Core/}\allowbreak{}\texttt{Ledger/}\allowbreak{}\texttt{JournalBatchValidator.}\allowbreak{}\texttt{cs}
\item
  Create:
  \texttt{tests/}\allowbreak{}\texttt{FirstGearBank.}\allowbreak{}\texttt{Core.}\allowbreak{}\texttt{Tests/}\allowbreak{}\texttt{Ledger/}\allowbreak{}\texttt{LedgerAccountIdTests.}\allowbreak{}\texttt{cs}
\item
  Create:
  \texttt{tests/}\allowbreak{}\texttt{FirstGearBank.}\allowbreak{}\texttt{Core.}\allowbreak{}\texttt{Tests/}\allowbreak{}\texttt{Ledger/}\allowbreak{}\texttt{JournalBalanceTests.}\allowbreak{}\texttt{cs}
\item
  Create:
  \texttt{tests/}\allowbreak{}\texttt{FirstGearBank.}\allowbreak{}\texttt{Core.}\allowbreak{}\texttt{Tests/}\allowbreak{}\texttt{Ledger/}\allowbreak{}\texttt{JournalSequenceTests.}\allowbreak{}\texttt{cs}
\item
  Create:
  \texttt{tests/}\allowbreak{}\texttt{FirstGearBank.}\allowbreak{}\texttt{Core.}\allowbreak{}\texttt{Tests/}\allowbreak{}\texttt{Ledger/}\allowbreak{}\texttt{JournalHashTests.}\allowbreak{}\texttt{cs}
\item
  Create:
  \texttt{tests/}\allowbreak{}\texttt{FirstGearBank.}\allowbreak{}\texttt{Core.}\allowbreak{}\texttt{Tests/}\allowbreak{}\texttt{Ledger/}\allowbreak{}\texttt{JournalPayloadTests.}\allowbreak{}\texttt{cs}
\item
  Create:
  \texttt{tests/}\allowbreak{}\texttt{FirstGearBank.}\allowbreak{}\texttt{Core.}\allowbreak{}\texttt{Tests/}\allowbreak{}\texttt{TestSupport/}\allowbreak{}\texttt{JournalTestData.}\allowbreak{}\texttt{cs}
\item
  Create:
  \texttt{tests/}\allowbreak{}\texttt{FirstGearBank.}\allowbreak{}\texttt{Core.}\allowbreak{}\texttt{Tests/}\allowbreak{}\texttt{TestSupport/}\allowbreak{}\texttt{PlayerIds.}\allowbreak{}\texttt{cs}
\item
  Create:
  \texttt{tests/}\allowbreak{}\texttt{FirstGearBank.}\allowbreak{}\texttt{Core.}\allowbreak{}\texttt{Tests/}\allowbreak{}\texttt{TestSupport/}\allowbreak{}\texttt{FinancialTimestamps.}\allowbreak{}\texttt{cs}
\end{itemize}

\textbf{Interfaces:}

\begin{itemize}
\tightlist
\item
  Consumes: currencies, units, identifiers, timestamps, and failures
  from Task 5.
\item
  Produces: immutable
  \passthrough{\lstinline!JournalDraft!}/\passthrough{\lstinline!JournalRecord!}
  values and
  \texttt{JournalBatchValidator.}\allowbreak{}\texttt{Validate},
  consumed by projection, coordinator, and persistence.
\end{itemize}

\begin{itemize}
\tightlist
\item[$\square$]
  \textbf{Step 1: Write failing account-identity and per-currency
  balance tests}
\end{itemize}

\begin{lstlisting}
[Fact]
public void Rusty_and_temporal_postings_may_not_net_against_each_other()
{
    var record = JournalTestData.Record(
        JournalTestData.Posting(
            LedgerAccountId.RustyPhysicalCustody,
            BankCurrency.Rusty,
            new BankUnits(1_000_000)),
        JournalTestData.Posting(
            LedgerAccountId.TemporalPhysicalCustody,
            BankCurrency.Temporal,
            new BankUnits(-1_000_000)));

    var result = JournalBatchValidator.Validate([record], JournalValidationAnchor.Empty);

    Assert.Equal(BankErrorCode.UnbalancedJournal, result.Failure?.Code);
}

[Fact]
public void Sum_uses_wider_arithmetic_instead_of_overflowing_long()
{
    var record = JournalTestData.Record(
        JournalTestData.Posting(
            LedgerAccountId.RustyPhysicalCustody,
            BankCurrency.Rusty,
            new BankUnits(long.MaxValue)),
        JournalTestData.Posting(
            LedgerAccountId.RustySystemCorrection,
            BankCurrency.Rusty,
            new BankUnits(long.MaxValue)),
        JournalTestData.Posting(
            LedgerAccountId.RustyInterestExpense,
            BankCurrency.Rusty,
            new BankUnits(2)),
        JournalTestData.Posting(
            LedgerAccountId.RustyCapOverflow,
            BankCurrency.Rusty,
            new BankUnits(long.MinValue)),
        JournalTestData.Posting(
            LedgerAccountId.RustyCapOverflow,
            BankCurrency.Rusty,
            new BankUnits(long.MinValue)));

    Assert.True(JournalBatchValidator.Validate([record], JournalValidationAnchor.Empty).IsSuccess);
}
\end{lstlisting}

Add tests that customer accounts require a
\passthrough{\lstinline!PlayerId!}, CD liability requires a
\passthrough{\lstinline!CdId!}, system accounts reject owners, a posting
currency must match the account currency, and zero postings are rejected
rather than serialized.

\begin{itemize}
\tightlist
\item[$\square$]
  \textbf{Step 2: Run the balance tests and observe missing-type
  failures}
\end{itemize}

\begin{lstlisting}
dotnet test tests/FirstGearBank.Core.Tests/FirstGearBank.Core.Tests.csproj -c Release
\end{lstlisting}

Expected: FAIL because ledger accounts, postings, and the validator do
not exist.

\begin{itemize}
\tightlist
\item[$\square$]
  \textbf{Step 3: Implement stable ledger identities and overflow-safe
  balancing}
\end{itemize}

Freeze these numeric discriminators:

\begin{lstlisting}
public enum TransactionType : byte
{
    PhysicalDeposit = 1,
    PhysicalWithdrawal = 2,
    SavingsInterest = 3,
    TemporalGrossInterest = 4,
    TemporalStorage = 5,
    CdPurchase = 6,
    CdMaturity = 7,
    RustyTransfer = 8,
    AdminCorrection = 9
}

public enum LedgerAccountKind : byte
{
    CustomerRustySavingsLiability = 1,
    CustomerTemporalVaultLiability = 2,
    RustyPhysicalCustody = 3,
    TemporalPhysicalCustody = 4,
    RustyInterestExpense = 5,
    TemporalInterestExpense = 6,
    TemporalStorageIncome = 7,
    RustyCdPrincipalLiability = 8,
    RustyCdInterestExpense = 9,
    RustyCapOverflow = 10,
    RustySystemCorrection = 11,
    TemporalSystemCorrection = 12
}
\end{lstlisting}

\passthrough{\lstinline!LedgerAccountId!} has validated factories for
player, CD, and system accounts. \passthrough{\lstinline!Posting!}
construction is validated through a factory so no public path creates
zero, mismatched-currency, or malformed-owner postings. Sum signed units
into \passthrough{\lstinline!Int128!} per currency and require each sum
to equal zero.

\begin{itemize}
\tightlist
\item[$\square$]
  \textbf{Step 4: Verify balance and account tests pass}
\end{itemize}

\begin{lstlisting}
dotnet test tests/FirstGearBank.Core.Tests/FirstGearBank.Core.Tests.csproj -c Release
\end{lstlisting}

\begin{itemize}
\tightlist
\item[$\square$]
  \textbf{Step 5: Write failing sequence, uniqueness, payload, and
  hash-chain tests}
\end{itemize}

\begin{lstlisting}
[Fact]
public void One_bad_record_rejects_the_entire_correlated_batch()
{
    var first = JournalTestData.BalancedRecord(sequence: 7);
    var second = JournalTestData.UnbalancedRecord(sequence: 8);

    var result = JournalBatchValidator.Validate(
        [first, second],
        new JournalValidationAnchor(6, JournalTestData.PreviousHash));

    Assert.Equal(BankErrorCode.UnbalancedJournal, result.Failure?.Code);
}

[Fact]
public void Changing_a_frozen_name_snapshot_breaks_the_hash_chain()
{
    var record = JournalTestData.TransferRecord();
    var changedPayload = ((RustyTransferPayload)record.Payload) with { SenderDisplayName = "Changed" };
    var tampered = JournalTestData.CloneWithPayload(record, changedPayload);

    Assert.Equal(
        BankErrorCode.JournalIntegrityFailure,
        JournalBatchValidator.Validate([tampered], JournalValidationAnchor.Empty).Failure?.Code);
}
\end{lstlisting}

Add sequence-origin/\allowbreak{}gap/\allowbreak{}duplicate tests,
duplicate operation ID, mismatched shared
\passthrough{\lstinline!CommandId!}, invalid payload/type pairing,
overlong name/reason, missing or malformed invocation key, request-key
use on an automatic-only transaction, automatic-key use on a player-only
transaction, prior-hash mismatch, deterministic canonical bytes, and
changed
posting/\allowbreak{}party/\allowbreak{}timestamp/\allowbreak{}invocation-key
detection. Freeze the admin-correction
\passthrough{\lstinline!NoticeId!} derivation with a golden vector and
prove a different operation ID or notice domain cannot collide with it.

\begin{lstlisting}
dotnet test tests/FirstGearBank.Core.Tests/FirstGearBank.Core.Tests.csproj -c Release
\end{lstlisting}

Expected: FAIL on sequence, payload, or canonical-integrity behavior
that has not been implemented.

\begin{itemize}
\tightlist
\item[$\square$]
  \textbf{Step 6: Implement immutable records and typed payload codecs}
\end{itemize}

\begin{lstlisting}
public interface IJournalPayload
{
    TransactionType TransactionType { get; }
    ushort PayloadSchemaVersion { get; }
}

public sealed class JournalDraft
{
    public OperationId OperationId { get; }
    public BankInvocationIntent InvocationIntent { get; }
    public ushort RecordSchemaVersion { get; }
    public FinancialTimestamp Timestamp { get; }
    public ImmutableArray<JournalPartySnapshot> Parties { get; }
    public IJournalPayload Payload { get; }
    public ImmutableArray<Posting> Postings { get; }

    internal JournalDraft(
        OperationId operationId,
        BankInvocationIntent invocationIntent,
        ushort recordSchemaVersion,
        FinancialTimestamp timestamp,
        ImmutableArray<JournalPartySnapshot> parties,
        IJournalPayload payload,
        ImmutableArray<Posting> postings);
}

public sealed class JournalRecord
{
    public long Sequence { get; }
    public OperationId OperationId { get; }
    public CommandId CommandId { get; }
    public BankInvocationKey InvocationKey { get; }
    public ushort RecordSchemaVersion { get; }
    public FinancialTimestamp Timestamp { get; }
    public ImmutableArray<JournalPartySnapshot> Parties { get; }
    public IJournalPayload Payload { get; }
    public ImmutableArray<Posting> Postings { get; }
    public ImmutableArray<byte> PreviousRecordHash { get; }
    public ImmutableArray<byte> RecordHash { get; }

    internal JournalRecord(
        long sequence,
        OperationId operationId,
        CommandId commandId,
        BankInvocationKey invocationKey,
        ushort recordSchemaVersion,
        FinancialTimestamp timestamp,
        ImmutableArray<JournalPartySnapshot> parties,
        IJournalPayload payload,
        ImmutableArray<Posting> postings,
        ImmutableArray<byte> previousRecordHash,
        ImmutableArray<byte> recordHash)
    {
        Sequence = sequence;
        OperationId = operationId;
        CommandId = commandId;
        InvocationKey = invocationKey;
        RecordSchemaVersion = recordSchemaVersion;
        Timestamp = timestamp;
        Parties = parties;
        Payload = payload;
        Postings = postings;
        PreviousRecordHash = previousRecordHash;
        RecordHash = recordHash;
    }
}
\end{lstlisting}

Use sealed payload records, never an authoritative free-form dictionary.
Each payload contains exact units and identifiers needed for replay:
settlement ID for physical operations;
opening/\allowbreak{}closing/\allowbreak{}cap flag for accrual; stable
sender/recipient IDs plus frozen names and notice ID for transfers;
administrator/target IDs plus frozen names, a stable
\passthrough{\lstinline!NoticeId!}, mandatory bounded reason, and
optional original operation for corrections. Derive a correction notice
ID with the versioned domain
\texttt{FGB/}\allowbreak{}\texttt{ADMIN-}\allowbreak{}\texttt{CORRECTION-}\allowbreak{}\texttt{NOTICE/}\allowbreak{}\texttt{1\textbackslash{}}\allowbreak{}\texttt{0}
from its operation ID, and freeze that derivation with a golden test. CD
payload discriminators remain reserved; the certificates plan adds their
fully typed contractual records through a section/journal schema
migration before any CD can be issued.

Canonical encoding is length-prefixed binary with explicit enum widths
and the domain prefix
\texttt{FGB/}\allowbreak{}\texttt{JOURNAL/}\allowbreak{}\texttt{1\textbackslash{}}\allowbreak{}\texttt{0}.
Hash
\texttt{prefix\ }\allowbreak{}\texttt{\textbar{}}\allowbreak{}\texttt{\textbar{}}\allowbreak{}\texttt{\ }\allowbreak{}\texttt{previousHash\ }\allowbreak{}\texttt{\textbar{}}\allowbreak{}\texttt{\textbar{}}\allowbreak{}\texttt{\ }\allowbreak{}\texttt{canonicalRecordWithoutRecordHash}
with SHA-256. Use constant-time byte comparison when verifying hashes.

Every new record has exactly one
\passthrough{\lstinline!BankInvocationKey!}. The complete
pre-finalization \passthrough{\lstinline!JournalDraft!} shape is the one
above: it owns its already-reserved
\passthrough{\lstinline!OperationId!}, intent, schema, timestamp,
parties, typed payload, and postings, and deliberately has no sequence,
command ID, resolved invocation key, previous hash, record hash, or
revision. Its closed \passthrough{\lstinline!BankInvocationIntent!}
union is \passthrough{\lstinline!Caller!} or a particular automatic key
plus its canonical payload digest; it has no public way to install a
final invocation key. The coordinator resolves
\passthrough{\lstinline!Caller!} to the admitted request/automatic
context and admits every subordinate automatic intent before
constructing records. Physical deposit/withdrawal, transfer, and direct
correction records require the authenticated caller variant; clock
accrual and later CD maturity records require an automatic variant. A
multi-record catch-up batch shares the coordinator\textquotesingle s one
command ID but may contain deterministic subordinate automatic keys only
when the coordinator persists all of them in the same candidate; no
record may carry an unindexed automatic key.

Every production draft factory requires the exact
\passthrough{\lstinline!OperationId!} reserved from the current
coordinator-owned mutation identifier scope; factories never mint one
and the finalizer never replaces one.
\texttt{AdminCorrectionPayload.}\allowbreak{}\texttt{Create} derives its
\passthrough{\lstinline!NoticeId!} from that same operation ID, and
\texttt{JournalRecordFactory.}\allowbreak{}\texttt{AdminCorrection}
rejects a payload whose derived notice does not match the draft/header
operation. The golden notice-vector test therefore covers the actual
persisted record identity rather than an unrelated test input.

Create \passthrough{\lstinline!JournalTestData!},
\passthrough{\lstinline!PlayerIds!}, and
\passthrough{\lstinline!FinancialTimestamps!} as deterministic test-only
builders in the listed support files. They must call production
factories for valid records; only explicitly corrupt fixtures may bypass
a production factory through an internal test-visible constructor.

\begin{itemize}
\tightlist
\item[$\square$]
  \textbf{Step 7: Write and observe failing exact posting-template
  tests}
\end{itemize}

Add one test per foundation transaction type. For example, a rusty
deposit draft must contain exactly custody \passthrough{\lstinline!+A!}
and the same player\textquotesingle s rusty liability
\passthrough{\lstinline!-A!}, in canonical posting order, with no zero
entry:

\begin{lstlisting}
[Fact]
public void Physical_deposit_uses_the_locked_two_posting_template()
{
    var draft = JournalRecordFactory.PhysicalDeposit(
        JournalTestData.OperationOne,
        BankInvocationIntent.Caller,
        FinancialTimestamps.MonthFive,
        JournalTestData.RustyDepositPayload(PlayerIds.Ada, units: 250_000));

    Assert.Equal(JournalTestData.OperationOne, draft.OperationId);
    Assert.Equal(2, draft.Postings.Length);
    Assert.Equal(LedgerAccountId.RustyPhysicalCustody, draft.Postings[0].AccountId);
    Assert.Equal(250_000, draft.Postings[0].SignedUnits.Value);
    Assert.Equal(LedgerAccountId.CustomerRustySavings(PlayerIds.Ada), draft.Postings[1].AccountId);
    Assert.Equal(-250_000, draft.Postings[1].SignedUnits.Value);
}
\end{lstlisting}

\begin{lstlisting}
dotnet test tests/FirstGearBank.Core.Tests/FirstGearBank.Core.Tests.csproj -c Release
\end{lstlisting}

Expected: FAIL because the posting-template factories are not
implemented.

\begin{itemize}
\tightlist
\item[$\square$]
  \textbf{Step 8: Implement the exact posting templates already fixed by
  the specification}
\end{itemize}

\begin{lstlisting}
public static class JournalRecordFactory
{
    public static JournalDraft PhysicalDeposit(
        OperationId operationId,
        BankInvocationIntent invocationIntent,
        FinancialTimestamp timestamp,
        PhysicalDepositPayload payload);

    public static JournalDraft PhysicalWithdrawal(
        OperationId operationId,
        BankInvocationIntent invocationIntent,
        FinancialTimestamp timestamp,
        PhysicalWithdrawalPayload payload);

    public static JournalDraft RustyTransfer(
        OperationId operationId,
        BankInvocationIntent invocationIntent,
        FinancialTimestamp timestamp,
        RustyTransferPayload payload);

    public static JournalDraft AdminCorrection(
        OperationId operationId,
        BankInvocationIntent invocationIntent,
        FinancialTimestamp timestamp,
        AdminCorrectionPayload payload);
}
\end{lstlisting}

Deposit posts custody \passthrough{\lstinline!+A!} and customer
liability \passthrough{\lstinline!-A!}; withdrawal reverses both;
transfer posts sender liability \passthrough{\lstinline!+A!} and
recipient liability \passthrough{\lstinline!-A!}; correction posts the
per-currency system correction account \passthrough{\lstinline!+D!} and
target liability \passthrough{\lstinline!-D!}. Accrual factories accept
already rounded exact components and emit the section 6/9 posting
pattern without recalculating economic values. Add an exact
correction-factory test proving the draft operation ID, final record
header operation ID, and
\texttt{NoticeId.}\allowbreak{}\texttt{ForAdminCorrection(}\allowbreak{}\texttt{operationId)}\allowbreak{}
all originate from the same reserved ID.

\begin{itemize}
\tightlist
\item[$\square$]
  \textbf{Step 9: Verify all journal tests and commit}
\end{itemize}

\begin{lstlisting}
dotnet test tests/FirstGearBank.Core.Tests/FirstGearBank.Core.Tests.csproj -c Release
git diff --check
git add FirstGearBank.Core/Ledger tests/FirstGearBank.Core.Tests/Ledger `
  tests/FirstGearBank.Core.Tests/TestSupport/JournalTestData.cs `
  tests/FirstGearBank.Core.Tests/TestSupport/PlayerIds.cs `
  tests/FirstGearBank.Core.Tests/TestSupport/FinancialTimestamps.cs
git commit -m "feat: add typed balanced banking journal"
\end{lstlisting}

\subsubsection{Task 8: Rebuild Account and System Projections from the
Journal}\label{task-8-rebuild-account-and-system-projections-from-the-journal}

\textbf{Files:}

\begin{itemize}
\tightlist
\item
  Create:
  \texttt{FirstGearBank.}\allowbreak{}\texttt{Core/}\allowbreak{}\texttt{Projection/}\allowbreak{}\texttt{AccountTotals.}\allowbreak{}\texttt{cs}
\item
  Create:
  \texttt{FirstGearBank.}\allowbreak{}\texttt{Core/}\allowbreak{}\texttt{Projection/}\allowbreak{}\texttt{AccountHistoryDirection.}\allowbreak{}\texttt{cs}
\item
  Create:
  \texttt{FirstGearBank.}\allowbreak{}\texttt{Core/}\allowbreak{}\texttt{Projection/}\allowbreak{}\texttt{AccountHistoryEntry.}\allowbreak{}\texttt{cs}
\item
  Create:
  \texttt{FirstGearBank.}\allowbreak{}\texttt{Core/}\allowbreak{}\texttt{Projection/}\allowbreak{}\texttt{PlayerAccountProjection.}\allowbreak{}\texttt{cs}
\item
  Create:
  \texttt{FirstGearBank.}\allowbreak{}\texttt{Core/}\allowbreak{}\texttt{Projection/}\allowbreak{}\texttt{SystemLedgerProjection.}\allowbreak{}\texttt{cs}
\item
  Create:
  \texttt{FirstGearBank.}\allowbreak{}\texttt{Core/}\allowbreak{}\texttt{Projection/}\allowbreak{}\texttt{BankProjection.}\allowbreak{}\texttt{cs}
\item
  Create:
  \texttt{FirstGearBank.}\allowbreak{}\texttt{Core/}\allowbreak{}\texttt{Projection/}\allowbreak{}\texttt{ProjectionCheckpoint.}\allowbreak{}\texttt{cs}
\item
  Create:
  \texttt{FirstGearBank.}\allowbreak{}\texttt{Core/}\allowbreak{}\texttt{Projection/}\allowbreak{}\texttt{IJournalProjectionRule.}\allowbreak{}\texttt{cs}
\item
  Create:
  \texttt{FirstGearBank.}\allowbreak{}\texttt{Core/}\allowbreak{}\texttt{Projection/}\allowbreak{}\texttt{ProjectionRuleSet.}\allowbreak{}\texttt{cs}
\item
  Create:
  \texttt{FirstGearBank.}\allowbreak{}\texttt{Core/}\allowbreak{}\texttt{Projection/}\allowbreak{}\texttt{JournalProjector.}\allowbreak{}\texttt{cs}
\item
  Create:
  \texttt{FirstGearBank.}\allowbreak{}\texttt{Core/}\allowbreak{}\texttt{Projection/}\allowbreak{}\texttt{ProjectionValidator.}\allowbreak{}\texttt{cs}
\item
  Create:
  \texttt{tests/}\allowbreak{}\texttt{FirstGearBank.}\allowbreak{}\texttt{Core.}\allowbreak{}\texttt{Tests/}\allowbreak{}\texttt{Projection/}\allowbreak{}\texttt{AccountLifecycleProjectionTests.}\allowbreak{}\texttt{cs}
\item
  Create:
  \texttt{tests/}\allowbreak{}\texttt{FirstGearBank.}\allowbreak{}\texttt{Core.}\allowbreak{}\texttt{Tests/}\allowbreak{}\texttt{Projection/}\allowbreak{}\texttt{JournalProjectorTests.}\allowbreak{}\texttt{cs}
\item
  Create:
  \texttt{tests/}\allowbreak{}\texttt{FirstGearBank.}\allowbreak{}\texttt{Core.}\allowbreak{}\texttt{Tests/}\allowbreak{}\texttt{Projection/}\allowbreak{}\texttt{ProjectionInvariantTests.}\allowbreak{}\texttt{cs}
\item
  Create:
  \texttt{tests/}\allowbreak{}\texttt{FirstGearBank.}\allowbreak{}\texttt{Core.}\allowbreak{}\texttt{Tests/}\allowbreak{}\texttt{Projection/}\allowbreak{}\texttt{TransferHistoryProjectionTests.}\allowbreak{}\texttt{cs}
\end{itemize}

\textbf{Interfaces:}

\begin{itemize}
\tightlist
\item
  Consumes: a verified contiguous journal and registered typed payload
  rules from Task 7.
\item
  Produces: \passthrough{\lstinline!BankProjection!}, deterministic
  replay/checkpoints, and incremental affected-account validation for
  coordinator and persistence.
\end{itemize}

\begin{itemize}
\tightlist
\item[$\square$]
  \textbf{Step 1: Write failing lazy-account-lifecycle tests}
\end{itemize}

\begin{lstlisting}
[Fact]
public void Successful_physical_deposit_creates_the_first_account_projection()
{
    var record = JournalTestData.PhysicalDeposit(player: PlayerIds.Ada, units: 1_000_000);

    var result = JournalProjector.Foundation.Replay([record], ProjectionCheckpoint.Empty);

    Assert.True(result.IsSuccess);
    Assert.Equal(1_000_000, result.Value.Accounts[PlayerIds.Ada].RustySavings.Value);
}

[Fact]
public void Incoming_transfer_creates_an_offline_recipient_account()
{
    var records = JournalTestData.FinalizeBatch(
        JournalTestData.PhysicalDepositDraft(player: PlayerIds.Ada, units: 2_000_000),
        JournalTestData.RustyTransferDraft(
            sender: PlayerIds.Ada,
            recipient: PlayerIds.Byron,
            units: 500_000));

    var result = JournalProjector.Foundation.Replay(records, ProjectionCheckpoint.Empty);

    Assert.Equal(500_000, result.Value.Accounts[PlayerIds.Byron].RustySavings.Value);
}

[Fact]
public void Administrative_correction_cannot_open_an_account()
{
    var record = JournalTestData.AdminCorrection(player: PlayerIds.Ada, displayedDeltaUnits: 1_000_000);

    Assert.Equal(
        BankErrorCode.AccountNotFound,
        JournalProjector.Foundation.Replay([record], ProjectionCheckpoint.Empty).Failure?.Code);
}
\end{lstlisting}

The final case follows section 5\textquotesingle s exclusive rule: only
a successful deposit or received transfer opens the relationship.
\passthrough{\lstinline!FinalizeBatch!} is the Task 7 deterministic test
finalizer: it assigns the two records one contiguous sequence/hash chain
from \texttt{JournalValidationAnchor.}\allowbreak{}\texttt{Empty}. It
does not concatenate independently finalized records or keep hidden
global state.

\begin{itemize}
\tightlist
\item[$\square$]
  \textbf{Step 2: Run lifecycle tests and observe missing projection
  failures}
\end{itemize}

\begin{lstlisting}
dotnet test tests/FirstGearBank.Core.Tests/FirstGearBank.Core.Tests.csproj -c Release
\end{lstlisting}

Expected: FAIL because projection types do not exist.

\begin{itemize}
\tightlist
\item[$\square$]
  \textbf{Step 3: Implement immutable player/system projection records
  and rule registration}
\end{itemize}

\begin{lstlisting}
public sealed record PlayerAccountProjection(
    PlayerId PlayerId,
    BankUnits RustySavings,
    BankUnits TemporalVault,
    AccountTotals Totals,
    ImmutableArray<AccountHistoryEntry> History);

public sealed record BankProjection(
    ImmutableDictionary<PlayerId, PlayerAccountProjection> Accounts,
    SystemLedgerProjection SystemLedger,
    long AppliedJournalSequence,
    ImmutableArray<byte> AppliedJournalHash);

public interface IJournalProjectionRule
{
    TransactionType TransactionType { get; }
    BankResult<BankProjection> Apply(BankProjection current, JournalRecord record);
}
\end{lstlisting}

\texttt{ProjectionRuleSet.}\allowbreak{}\texttt{Foundation} registers
physical deposit/withdrawal, both accrual components, transfer, and
correction. Duplicate or missing transaction rules fail startup
validation. CD rules are added only with their fully typed payloads and
schema migration in the certificates plan.

\begin{itemize}
\tightlist
\item[$\square$]
  \textbf{Step 4: Implement liability-sign conversion and
  account-opening rules}
\end{itemize}

Compute displayed customer balances as the checked negation of
cumulative customer-liability ledger signs. Deposit may open its subject
account. Transfer requires an existing sender and may open only its
recipient. Withdrawal, interest, storage, and correction require an
existing target. Apply one record to a temporary immutable builder and
publish it only after all affected customer balances remain within
\texttt{[0,}\allowbreak{}\texttt{\ }\allowbreak{}\texttt{MaximumCustomerBalanceUnits]}.

\begin{itemize}
\tightlist
\item[$\square$]
  \textbf{Step 5: Verify account lifecycle tests pass}
\end{itemize}

\begin{lstlisting}
dotnet test tests/FirstGearBank.Core.Tests/FirstGearBank.Core.Tests.csproj -c Release
\end{lstlisting}

\begin{itemize}
\tightlist
\item[$\square$]
  \textbf{Step 6: Write failing replay, aggregate, and
  directional-history tests}
\end{itemize}

\begin{lstlisting}
[Fact]
public void Complete_replay_matches_incremental_application()
{
    var records = JournalTestData.MixedFoundationHistory();
    var projector = JournalProjector.Foundation;
    var incremental = ProjectionCheckpoint.Empty;

    foreach (var record in records)
    {
        incremental = projector.Apply(incremental, record).Value;
    }

    var replayed = projector.Replay(records, ProjectionCheckpoint.Empty).Value;

    var match = ProjectionValidator.MatchesReplay(incremental.Projection, replayed);

    Assert.True(match.IsSuccess);
    Assert.True(match.Value);
}

[Fact]
public void One_transfer_has_two_directional_account_views()
{
    var projection = JournalProjector.Foundation
        .Replay(JournalTestData.FundedTransferHistory(), ProjectionCheckpoint.Empty)
        .Value;

    Assert.Equal(AccountHistoryDirection.TransferOut, projection.Accounts[PlayerIds.Ada].History[^1].Direction);
    Assert.Equal(AccountHistoryDirection.TransferIn, projection.Accounts[PlayerIds.Byron].History[^1].Direction);
    Assert.Equal(
        projection.Accounts[PlayerIds.Ada].History[^1].OperationId,
        projection.Accounts[PlayerIds.Byron].History[^1].OperationId);
}
\end{lstlisting}

Add cumulative
deposit/\allowbreak{}withdrawal/\allowbreak{}interest/\allowbreak{}storage/\allowbreak{}transfer/\allowbreak{}correction
totals, signed negative system accounts, history order, unknown payload
rule, checkpoint hash mismatch, journal gap, negative customer result,
cap breach, and corrupted cached projection comparison.

\begin{lstlisting}
dotnet test tests/FirstGearBank.Core.Tests/FirstGearBank.Core.Tests.csproj -c Release
\end{lstlisting}

Expected: FAIL because complete replay, history, or cache comparison is
not implemented.

\begin{itemize}
\tightlist
\item[$\square$]
  \textbf{Step 7: Implement replay, checkpoints, totals, and cache
  comparison}
\end{itemize}

\begin{lstlisting}
public sealed class JournalProjector
{
    public static JournalProjector Foundation { get; }

    public BankResult<ProjectionCheckpoint> Apply(
        ProjectionCheckpoint checkpoint,
        JournalRecord nextRecord);

    public BankResult<BankProjection> Replay(
        IEnumerable<JournalRecord> records,
        ProjectionCheckpoint checkpoint);
}

public static class ProjectionValidator
{
    public static BankResult<bool> MatchesReplay(
        BankProjection cached,
        BankProjection replayed);
}
\end{lstlisting}

Replay first validates journal
continuity/\allowbreak{}hash/\allowbreak{}balance, then applies one
typed rule at a time. Incremental application touches only parties and
ledger accounts named by the record. Cache mismatch returns a repairable
diagnostic; it never creates a journal correction or promotes the cache
over replay.

\begin{itemize}
\tightlist
\item[$\square$]
  \textbf{Step 8: Verify all projection tests and commit}
\end{itemize}

\begin{lstlisting}
dotnet test tests/FirstGearBank.Core.Tests/FirstGearBank.Core.Tests.csproj -c Release
git diff --check
git add FirstGearBank.Core/Projection tests/FirstGearBank.Core.Tests/Projection
git commit -m "feat: add replayable bank projections"
\end{lstlisting}

\subsubsection{Task 9: Serialize Mutations Through the Idempotent
Banking
Coordinator}\label{task-9-serialize-mutations-through-the-idempotent-banking-coordinator}

\textbf{Files:}

\begin{itemize}
\tightlist
\item
  Create:
  \texttt{FirstGearBank.}\allowbreak{}\texttt{Core/}\allowbreak{}\texttt{State/}\allowbreak{}\texttt{BankState.}\allowbreak{}\texttt{cs}
\item
  Create:
  \texttt{FirstGearBank.}\allowbreak{}\texttt{Core/}\allowbreak{}\texttt{State/}\allowbreak{}\texttt{BankStateSnapshot.}\allowbreak{}\texttt{cs}
\item
  Create:
  \texttt{FirstGearBank.}\allowbreak{}\texttt{Core/}\allowbreak{}\texttt{State/}\allowbreak{}\texttt{BankStateRepository.}\allowbreak{}\texttt{cs}
\item
  Create:
  \texttt{FirstGearBank.}\allowbreak{}\texttt{Core/}\allowbreak{}\texttt{State/}\allowbreak{}\texttt{BankStateValidator.}\allowbreak{}\texttt{cs}
\item
  Create:
  \texttt{FirstGearBank.}\allowbreak{}\texttt{Core/}\allowbreak{}\texttt{State/}\allowbreak{}\texttt{BankDomainAvailability.}\allowbreak{}\texttt{cs}
\item
  Create:
  \texttt{FirstGearBank.}\allowbreak{}\texttt{Core/}\allowbreak{}\texttt{State/}\allowbreak{}\texttt{BankDomainSlot.}\allowbreak{}\texttt{cs}
\item
  Create:
  \texttt{FirstGearBank.}\allowbreak{}\texttt{Core/}\allowbreak{}\texttt{State/}\allowbreak{}\texttt{FinanceDomainState.}\allowbreak{}\texttt{cs}
\item
  Create:
  \texttt{FirstGearBank.}\allowbreak{}\texttt{Core/}\allowbreak{}\texttt{State/}\allowbreak{}\texttt{RequestControlState.}\allowbreak{}\texttt{cs}
\item
  Create:
  \texttt{FirstGearBank.}\allowbreak{}\texttt{Core/}\allowbreak{}\texttt{Recovery/}\allowbreak{}\texttt{QuarantineDomain.}\allowbreak{}\texttt{cs}
\item
  Create:
  \texttt{FirstGearBank.}\allowbreak{}\texttt{Core/}\allowbreak{}\texttt{Recovery/}\allowbreak{}\texttt{BankQuarantineState.}\allowbreak{}\texttt{cs}
\item
  Create:
  \texttt{FirstGearBank.}\allowbreak{}\texttt{Core/}\allowbreak{}\texttt{Recovery/}\allowbreak{}\texttt{QuarantineTransition.}\allowbreak{}\texttt{cs}
\item
  Create:
  \texttt{FirstGearBank.}\allowbreak{}\texttt{Core/}\allowbreak{}\texttt{Recovery/}\allowbreak{}\texttt{AuthorizedRecoveryContext.}\allowbreak{}\texttt{cs}
\item
  Create:
  \texttt{FirstGearBank.}\allowbreak{}\texttt{Core/}\allowbreak{}\texttt{Recovery/}\allowbreak{}\texttt{ArchivedRecoveryEvidence.}\allowbreak{}\texttt{cs}
\item
  Create:
  \texttt{FirstGearBank.}\allowbreak{}\texttt{Core/}\allowbreak{}\texttt{Recovery/}\allowbreak{}\texttt{RecoveryAuditEntry.}\allowbreak{}\texttt{cs}
\item
  Create:
  \texttt{FirstGearBank.}\allowbreak{}\texttt{Core/}\allowbreak{}\texttt{Recovery/}\allowbreak{}\texttt{RecoveryControlState.}\allowbreak{}\texttt{cs}
\item
  Create:
  \texttt{FirstGearBank.}\allowbreak{}\texttt{Core/}\allowbreak{}\texttt{Recovery/}\allowbreak{}\texttt{RecoveryLimits.}\allowbreak{}\texttt{cs}
\item
  Create:
  \texttt{FirstGearBank.}\allowbreak{}\texttt{Core/}\allowbreak{}\texttt{Requests/}\allowbreak{}\texttt{TerminalResponse.}\allowbreak{}\texttt{cs}
\item
  Create:
  \texttt{FirstGearBank.}\allowbreak{}\texttt{Core/}\allowbreak{}\texttt{Requests/}\allowbreak{}\texttt{ProcessedRequest.}\allowbreak{}\texttt{cs}
\item
  Create:
  \texttt{FirstGearBank.}\allowbreak{}\texttt{Core/}\allowbreak{}\texttt{Requests/}\allowbreak{}\texttt{OperationScope.}\allowbreak{}\texttt{cs}
\item
  Create:
  \texttt{FirstGearBank.}\allowbreak{}\texttt{Core/}\allowbreak{}\texttt{Requests/}\allowbreak{}\texttt{RequestRegistry.}\allowbreak{}\texttt{cs}
\item
  Create:
  \texttt{FirstGearBank.}\allowbreak{}\texttt{Core/}\allowbreak{}\texttt{Requests/}\allowbreak{}\texttt{TerminalResponseCache.}\allowbreak{}\texttt{cs}
\item
  Create:
  \texttt{FirstGearBank.}\allowbreak{}\texttt{Core/}\allowbreak{}\texttt{Requests/}\allowbreak{}\texttt{ProcessedAutomaticOperation.}\allowbreak{}\texttt{cs}
\item
  Create:
  \texttt{FirstGearBank.}\allowbreak{}\texttt{Core/}\allowbreak{}\texttt{Requests/}\allowbreak{}\texttt{AutomaticOperationRegistry.}\allowbreak{}\texttt{cs}
\item
  Create:
  \texttt{FirstGearBank.}\allowbreak{}\texttt{Core/}\allowbreak{}\texttt{Coordination/}\allowbreak{}\texttt{AuthenticatedCommandContext.}\allowbreak{}\texttt{cs}
\item
  Create:
  \texttt{FirstGearBank.}\allowbreak{}\texttt{Core/}\allowbreak{}\texttt{Coordination/}\allowbreak{}\texttt{AutomaticCommandContext.}\allowbreak{}\texttt{cs}
\item
  Create:
  \texttt{FirstGearBank.}\allowbreak{}\texttt{Core/}\allowbreak{}\texttt{Coordination/}\allowbreak{}\texttt{IBankMutation.}\allowbreak{}\texttt{cs}
\item
  Create:
  \texttt{FirstGearBank.}\allowbreak{}\texttt{Core/}\allowbreak{}\texttt{Coordination/}\allowbreak{}\texttt{IBankRecoveryMutation.}\allowbreak{}\texttt{cs}
\item
  Create:
  \texttt{FirstGearBank.}\allowbreak{}\texttt{Core/}\allowbreak{}\texttt{Coordination/}\allowbreak{}\texttt{IBankControlChange.}\allowbreak{}\texttt{cs}
\item
  Create:
  \texttt{FirstGearBank.}\allowbreak{}\texttt{Core/}\allowbreak{}\texttt{Coordination/}\allowbreak{}\texttt{BankControlChangeRuleSet.}\allowbreak{}\texttt{cs}
\item
  Create:
  \texttt{FirstGearBank.}\allowbreak{}\texttt{Core/}\allowbreak{}\texttt{Coordination/}\allowbreak{}\texttt{BankMutationReadView.}\allowbreak{}\texttt{cs}
\item
  Create:
  \texttt{FirstGearBank.}\allowbreak{}\texttt{Core/}\allowbreak{}\texttt{Coordination/}\allowbreak{}\texttt{MutationDisposition.}\allowbreak{}\texttt{cs}
\item
  Create:
  \texttt{FirstGearBank.}\allowbreak{}\texttt{Core/}\allowbreak{}\texttt{Coordination/}\allowbreak{}\texttt{MutationDecision.}\allowbreak{}\texttt{cs}
\item
  Create:
  \texttt{FirstGearBank.}\allowbreak{}\texttt{Core/}\allowbreak{}\texttt{Coordination/}\allowbreak{}\texttt{CoordinatorResponse.}\allowbreak{}\texttt{cs}
\item
  Create:
  \texttt{FirstGearBank.}\allowbreak{}\texttt{Core/}\allowbreak{}\texttt{Coordination/}\allowbreak{}\texttt{IFinancialTimestampSource.}\allowbreak{}\texttt{cs}
\item
  Create:
  \texttt{FirstGearBank.}\allowbreak{}\texttt{Core/}\allowbreak{}\texttt{Coordination/}\allowbreak{}\texttt{IIdentifierSource.}\allowbreak{}\texttt{cs}
\item
  Create:
  \texttt{FirstGearBank.}\allowbreak{}\texttt{Core/}\allowbreak{}\texttt{Coordination/}\allowbreak{}\texttt{IMutationIdentifierScope.}\allowbreak{}\texttt{cs}
\item
  Create:
  \texttt{FirstGearBank.}\allowbreak{}\texttt{Core/}\allowbreak{}\texttt{Coordination/}\allowbreak{}\texttt{BankingCoordinator.}\allowbreak{}\texttt{cs}
\item
  Create:
  \texttt{tests/}\allowbreak{}\texttt{FirstGearBank.}\allowbreak{}\texttt{Core.}\allowbreak{}\texttt{Tests/}\allowbreak{}\texttt{Requests/}\allowbreak{}\texttt{CommandPayloadDigestAdmissionTests.}\allowbreak{}\texttt{cs}
\item
  Create:
  \texttt{tests/}\allowbreak{}\texttt{FirstGearBank.}\allowbreak{}\texttt{Core.}\allowbreak{}\texttt{Tests/}\allowbreak{}\texttt{Requests/}\allowbreak{}\texttt{RequestRegistryTests.}\allowbreak{}\texttt{cs}
\item
  Create:
  \texttt{tests/}\allowbreak{}\texttt{FirstGearBank.}\allowbreak{}\texttt{Core.}\allowbreak{}\texttt{Tests/}\allowbreak{}\texttt{Requests/}\allowbreak{}\texttt{AutomaticOperationRegistryTests.}\allowbreak{}\texttt{cs}
\item
  Create:
  \texttt{tests/}\allowbreak{}\texttt{FirstGearBank.}\allowbreak{}\texttt{Core.}\allowbreak{}\texttt{Tests/}\allowbreak{}\texttt{Coordination/}\allowbreak{}\texttt{BankingCoordinatorTests.}\allowbreak{}\texttt{cs}
\item
  Create:
  \texttt{tests/}\allowbreak{}\texttt{FirstGearBank.}\allowbreak{}\texttt{Core.}\allowbreak{}\texttt{Tests/}\allowbreak{}\texttt{Recovery/}\allowbreak{}\texttt{QuarantineTransitionTests.}\allowbreak{}\texttt{cs}
\item
  Create:
  \texttt{tests/}\allowbreak{}\texttt{FirstGearBank.}\allowbreak{}\texttt{Core.}\allowbreak{}\texttt{Tests/}\allowbreak{}\texttt{State/}\allowbreak{}\texttt{BankStateRepositoryTests.}\allowbreak{}\texttt{cs}
\item
  Create:
  \texttt{tests/}\allowbreak{}\texttt{FirstGearBank.}\allowbreak{}\texttt{Core.}\allowbreak{}\texttt{Tests/}\allowbreak{}\texttt{TestSupport/}\allowbreak{}\texttt{RequestTestData.}\allowbreak{}\texttt{cs}
\item
  Create:
  \texttt{tests/}\allowbreak{}\texttt{FirstGearBank.}\allowbreak{}\texttt{Core.}\allowbreak{}\texttt{Tests/}\allowbreak{}\texttt{TestSupport/}\allowbreak{}\texttt{CoordinatorTestRig.}\allowbreak{}\texttt{cs}
\item
  Create:
  \texttt{tests/}\allowbreak{}\texttt{FirstGearBank.}\allowbreak{}\texttt{Core.}\allowbreak{}\texttt{Tests/}\allowbreak{}\texttt{TestSupport/}\allowbreak{}\texttt{CoordinatorTestMutations.}\allowbreak{}\texttt{cs}
\end{itemize}

\textbf{Interfaces:}

\begin{itemize}
\tightlist
\item
  Consumes: immutable journal/projection contracts and the revision
  vector.
\item
  Produces: authenticated and automatic
  \texttt{BankingCoordinator.}\allowbreak{}\texttt{Execute} paths,
  immutable \passthrough{\lstinline!BankStateSnapshot!}, staging-state
  semantics, and replay barriers used by all server mutations and
  persistence.
\end{itemize}

\begin{itemize}
\tightlist
\item[$\square$]
  \textbf{Step 1: Write failing canonical-digest and request-admission
  tests}
\end{itemize}

\begin{lstlisting}
[Fact]
public void Digest_is_domain_and_schema_separated()
{
    var payload = new byte[] { 10, 20, 30 };

    var v1 = CommandPayloadDigest.Compute("deposit", 1, payload);
    var v2 = CommandPayloadDigest.Compute("deposit", 2, payload);
    var other = CommandPayloadDigest.Compute("withdraw", 1, payload);

    Assert.False(v1.FixedTimeEquals(v2));
    Assert.False(v1.FixedTimeEquals(other));
}

[Fact]
public void Accepted_business_rejection_advances_high_water_without_a_journal_record()
{
    var registry = RequestRegistry.Empty.Open(RequestTestData.Scope);

    var result = registry.FinalizeNext(
        RequestTestData.Key(sequence: 1),
        RequestTestData.Digest,
        TerminalResponse.Rejected(BankErrorCode.InsufficientFunds),
        committedOperationIds: []);

    Assert.Equal(1, result.Value.Scopes[RequestTestData.Scope.Nonce].HighWaterSequence);
    Assert.Empty(result.Value.PermanentCommittedKeys);
}
\end{lstlisting}

Add inactive scope, sequence zero, sequence gap, identical cached retry,
changed-payload reuse, cache eviction, expired response, successful
permanent key, close, and restart invalidation cases. Add
automatic-operation cases proving that a successful deterministic key is
permanent, an identical duplicate returns its stored result without
evaluation, a changed operation cannot reuse the key, and save/reload
preserves the barrier.

\begin{itemize}
\tightlist
\item[$\square$]
  \textbf{Step 2: Run request tests and observe missing-type failures}
\end{itemize}

\begin{lstlisting}
dotnet test tests/FirstGearBank.Core.Tests/FirstGearBank.Core.Tests.csproj -c Release
\end{lstlisting}

Expected: FAIL because request-state types do not exist.

\begin{itemize}
\tightlist
\item[$\square$]
  \textbf{Step 3: Implement versioned digests, scopes, caches, and
  replay barriers}
\end{itemize}

Compute:

\begin{lstlisting}
SHA-256(
  ASCII "FGB/COMMAND-PAYLOAD/1\0" ||
  UInt16BigEndian(command-kind UTF-8 byte length) || command-kind bytes ||
  UInt16BigEndian(command schema version) ||
  UInt32BigEndian(payload length) || canonical payload bytes)
\end{lstlisting}

Use constant-time digest comparison.
\passthrough{\lstinline!OperationScope!} binds one authenticated
\passthrough{\lstinline!PlayerId!}, a nonreused nonce, the next
sequence, and bounded terminal responses. Production keeps at least
1,024 full terminal bodies per player; constructors accept a lower
positive limit for tests. Only successful financial requests enter
\passthrough{\lstinline!PermanentCommittedKeys!}. A processed sequence
whose body was pruned returns \texttt{AlreadyProcessedResponseExpired};
it never executes again.

\passthrough{\lstinline!CommandPayloadDigest!} is a 32-byte value type
with structural
\passthrough{\lstinline!Equals!}/\passthrough{\lstinline!GetHashCode!},
a \passthrough{\lstinline!FixedTimeEquals!} method used for admission,
and a defensive-copy byte export. It never inherits the shallow equality
behavior of an immutable array.

\texttt{AutomaticOperationRegistry} permanently maps each
\passthrough{\lstinline!AutomaticOperationKey!} to its committed
operation IDs and resulting revision together with the
domain/schema-separated digest of its canonical automatic-command
payload. Successful automatic financial work, including a deliberate
rounded-zero checkpoint, installs that mapping in the same revision. An
identical duplicate returns the stored result before capturing time or
issuing identifiers; the same key with a different digest returns
\passthrough{\lstinline!RequestPayloadMismatch!}. An invariant fault
does not install the key and may be retried after the fault is repaired.

\begin{itemize}
\tightlist
\item[$\square$]
  \textbf{Step 4: Verify all request-state tests pass}
\end{itemize}

\begin{lstlisting}
dotnet test tests/FirstGearBank.Core.Tests/FirstGearBank.Core.Tests.csproj -c Release
\end{lstlisting}

\begin{itemize}
\tightlist
\item[$\square$]
  \textbf{Step 5: Write failing coordinator atomicity and one-timestamp
  tests}
\end{itemize}

\begin{lstlisting}
[Fact]
public void Correlated_records_publish_together_with_one_timestamp()
{
    var clock = new CountingTimestampSource(FinancialTimestamps.MonthFive);
    var coordinator = CoordinatorTestRig.Create(clock);
    var mutation = new TwoRecordMutation(PlayerIds.Byron, PlayerIds.Ada);

    var response = coordinator.Execute(RequestTestData.Context(sequence: 1), mutation);
    var finance = coordinator.Current.Finance.TryGetVerified().Value;

    Assert.True(response.IsSuccess);
    Assert.Equal(1, clock.CaptureCount);
    Assert.Equal(2, finance.Journal.Length);
    Assert.All(
        finance.Journal,
        record => Assert.Equal(FinancialTimestamps.MonthFive, record.Timestamp));
}

[Fact]
public void Invalid_second_record_publishes_neither_record()
{
    var coordinator = CoordinatorTestRig.Create();
    var before = coordinator.Current;

    var response = coordinator.Execute(
        RequestTestData.Context(sequence: 1),
        new MutationWithUnbalancedSecondRecord());
    var currentFinance = coordinator.Current.Finance.TryGetVerified().Value;
    var currentRequests = coordinator.Current.RequestControl.TryGetVerified().Value;

    Assert.Equal(BankErrorCode.UnbalancedJournal, response.Failure?.Code);
    Assert.True(before.Finance.TryGetVerified().Value.Journal.AsSpan().SequenceEqual(currentFinance.Journal.AsSpan()));
    Assert.Same(before.Finance.TryGetVerified().Value.Projection, currentFinance.Projection);
    Assert.Equal(before.Revisions.Global.Value + 1, coordinator.Current.Revisions.Global.Value);
    Assert.Equal(1, currentRequests.Requests.GetHighWater(RequestTestData.Scope.Nonce));
}

[Fact]
public void Mutation_cannot_emit_finance_while_declaring_only_configuration()
{
    var coordinator = CoordinatorTestRig.Create();

    var response = coordinator.Execute(
        RequestTestData.Context(sequence: 1),
        new UnderDeclaredFinanceMutation());

    Assert.Equal(BankErrorCode.UndeclaredMutationDomain, response.Failure?.Code);
    Assert.Empty(coordinator.Current.Finance.TryGetVerified().Value.Journal);
    Assert.Equal(
        1,
        coordinator.Current.RequestControl.TryGetVerified().Value.Requests.GetHighWater(RequestTestData.Scope.Nonce));
}

[Fact]
public void Undeclared_quarantined_read_is_blocked_even_when_the_mutation_ignores_the_failure()
{
    var coordinator = CoordinatorTestRig.WithUnavailableFinance();

    var response = coordinator.Execute(
        RequestTestData.Context(sequence: 1),
        new ConfigurationMutationThatAttemptsFinanceRead());

    Assert.Equal(BankErrorCode.UndeclaredMutationDomain, response.Failure?.Code);
    Assert.True(coordinator.Current.Finance.IsUnavailable);
    Assert.False(coordinator.Current.RequestControl.TryGetVerified().Value.Requests.HasSuccessful(RequestTestData.Key(1)));
}
\end{lstlisting}

Both adversarial fixtures declare only
\passthrough{\lstinline!Configuration!}: the first emits a finance
draft, and the second calls \passthrough{\lstinline!ReadFinance!},
ignores its failure, and otherwise returns a configuration-only
decision. Add tests for ordinal \passthrough{\lstinline!PlayerId!}
participant ordering, duplicate operation IDs, cap/negative projection
rejection, quarantined required domain, terminal zero-record rejection,
replay without a new clock/identifier call, and one global plus
only-changed subsystem revision increment. Exercise the automatic
overload with a deterministic month-boundary key and prove duplicate and
restored duplicate calls do not recapture time or issue a
command/operation ID. Add identifier-contract cases for a draft that
uses a nonissued operation ID, two drafts that reuse one reserved ID, an
accepted decision that leaves a reserved operation ID unused, and final
records/terminal replay metadata that preserve the exact draft IDs in
canonical draft order. Each case must fail or succeed because of the
real coordinator result and state, not because a fake source was merely
called. In \texttt{QuarantineTransitionTests}, cover
undeclared/unavailable reads, ordinary-change attempts to touch
quarantine, wrong administrator, missing permission, empty/overlong
reason, stale evidence digest, clear without a verified replacement,
successful set and clear, exact raw-evidence archival, audit
identity/frozen name/reason, changed \passthrough{\lstinline!Recovery!}
plus target revisions, and rejection before mutation when the evidence
archive would exceed its structural bound. In
\texttt{BankStateRepositoryTests}, prove capture alone does not clear
\passthrough{\lstinline!HasUnstagedChanges!}, exact-revision staging
success clears only that flag, \texttt{RequiresEngineSaveRestage}
remains true, a newer publication cannot be cleared by an older staging
result, and staging failure retains both the state and diagnostic.

\begin{lstlisting}
dotnet test tests/FirstGearBank.Core.Tests/FirstGearBank.Core.Tests.csproj -c Release
\end{lstlisting}

Expected: FAIL because atomic coordinator publication and replay
behavior are not implemented.

\begin{itemize}
\tightlist
\item[$\square$]
  \textbf{Step 6: Implement the mutation and coordinator contracts}
\end{itemize}

\begin{lstlisting}
public interface IBankMutation
{
    // Every business domain the mutation may read or change. Requests is coordinator-owned and implicit.
    BankMutationDomains RequiredDomains { get; }
    ImmutableArray<PlayerId> Participants { get; }

    MutationDecision Evaluate(
        BankMutationReadView current,
        FinancialTimestamp timestamp,
        CommandId commandId,
        IMutationIdentifierScope identifiers);
}

public interface IBankRecoveryMutation : IBankMutation
{
    QuarantineTransition RequestedTransition { get; }
}

public interface IIdentifierSource
{
    CommandId NextCommandId();
    OperationId NextOperationId();
    SettlementId NextSettlementId();
}

public interface IMutationIdentifierScope
{
    OperationId ReserveOperationId();
    SettlementId ReserveSettlementId();
}

public sealed record AuthenticatedCommandContext(
    PlayerId AuthenticatedPlayer,
    RequestKey RequestKey,
    CommandPayloadDigest PayloadDigest);

public enum MutationDisposition : byte
{
    Accepted = 1,
    Rejected = 2
}

public interface IBankControlChange
{
    BankMutationDomains Domain { get; }
    BankInvocationIntent InvocationIntent { get; }
}

public sealed class BankMutationReadView
{
    public BankResult<FinanceDomainState> ReadFinance();
    public BankResult<BankConfiguration> ReadConfiguration();

    internal BankMutationDomains ObservedReadDomains { get; }
}

public sealed class MutationDecision
{
    public MutationDisposition Disposition { get; }
    public BankFailure? Failure { get; }
    public ImmutableArray<JournalDraft> JournalDrafts { get; }
    public ImmutableArray<IBankControlChange> ControlChanges { get; }

    public static MutationDecision Accepted(
        ImmutableArray<JournalDraft> journalDrafts,
        ImmutableArray<IBankControlChange> controlChanges);

    public static MutationDecision Rejected(BankFailure failure);

    private MutationDecision(
        MutationDisposition disposition,
        BankFailure? failure,
        ImmutableArray<JournalDraft> journalDrafts,
        ImmutableArray<IBankControlChange> controlChanges);
}

public enum BankDomainAvailability : byte
{
    Verified = 1,
    Unavailable = 2
}

public sealed class BankDomainSlot<T>
{
    public BankDomainAvailability Availability { get; }
    public bool IsUnavailable { get; }
    public BankResult<T> TryGetVerified();

    public static BankDomainSlot<T> Verified(T value);
    public static BankDomainSlot<T> Unavailable(QuarantineDomain domain, EvidenceDigest evidenceDigest);
}

public sealed record FinanceDomainState(
    ImmutableArray<JournalRecord> Journal,
    BankProjection Projection);

public sealed record RequestControlState(
    RequestRegistry Requests,
    AutomaticOperationRegistry AutomaticOperations);

public enum QuarantineTransitionKind : byte
{
    Set = 1,
    ClearAfterRepair = 2
}

public sealed record QuarantineTransition(
    QuarantineTransitionKind Kind,
    QuarantineDomain Domain,
    EvidenceDigest ExpectedEvidenceDigest,
    string MandatoryReason);

public sealed class AuthorizedRecoveryContext
{
    public PlayerId Administrator { get; }
    public string FrozenAdministratorName { get; }
    public string GrantedPermission { get; }

    internal AuthorizedRecoveryContext(
        PlayerId administrator,
        string frozenAdministratorName,
        string grantedPermission);
}

public sealed record RecoveryControlState(
    ImmutableDictionary<EvidenceDigest, ArchivedRecoveryEvidence> EvidenceArchive,
    ImmutableArray<RecoveryAuditEntry> AuditEntries);

public static class RecoveryLimits
{
    public const int MaximumArchivedEvidenceBytes = 192 * 1024 * 1024;
}

public sealed record BankState(
    BankDomainSlot<FinanceDomainState> Finance,
    BankDomainSlot<RequestControlState> RequestControl,
    BankDomainSlot<BankConfiguration> Configuration,
    RecoveryControlState RecoveryControl,
    BankQuarantineState Quarantine,
    BankRevisionVector Revisions);

public sealed record AutomaticCommandContext(
    AutomaticOperationKey OperationKey,
    CommandPayloadDigest PayloadDigest);

public sealed class BankingCoordinator
{
    public CoordinatorResponse Execute(
        AuthenticatedCommandContext context,
        IBankMutation mutation);

    public CoordinatorResponse Execute(
        AutomaticCommandContext context,
        IBankMutation mutation);

    public CoordinatorResponse ExecuteRecovery(
        AuthenticatedCommandContext context,
        AuthorizedRecoveryContext authorization,
        IBankRecoveryMutation mutation);
}

public sealed record CoordinatorResponse(
    BankInvocationKey InvocationKey,
    bool IsSuccess,
    BankFailure? Failure,
    ImmutableArray<OperationId> OperationIds,
    BankRevision ResultingRevision);
\end{lstlisting}

\passthrough{\lstinline!MutationDecision!} factories enforce that
Accepted has no failure and Rejected has one failure and no
drafts/control changes. \passthrough{\lstinline!RequiredDomains!} is the
complete set of business domains that \passthrough{\lstinline!Evaluate!}
may read or change; it excludes the coordinator-owned
\passthrough{\lstinline!Requests!} domain, which every admitted
invocation gets implicitly. The mutation never receives
\passthrough{\lstinline!BankState!}. The coordinator constructs a fresh
\passthrough{\lstinline!BankMutationReadView!} with only the declared,
verified immutable domain slices. Every \passthrough{\lstinline!Read*!}
call records its domain before returning; an undeclared or unavailable
read returns \passthrough{\lstinline!DomainUnavailable!} without
exposing a placeholder or underlying value. A mutation that ignores that
failure is still rejected afterward because the view records the
attempted domain. Later phases add explicit typed accessors as domains
become real; no generic escape hatch, service locator, reflection
accessor, or raw-state property is permitted.

Only registered sealed control-change types may run.
\passthrough{\lstinline!IBankControlChange!} carries data but has no
\passthrough{\lstinline!Apply(BankState)!} method.
\texttt{BankControlChangeRuleSet} owns the closed type switch, reads
only the change\textquotesingle s one declared non-Finance/non-Requests
domain, and applies its typed replacement/delta to a temporary
candidate. Explicit domain-specific semantic comparers prove that only
the declared top-level domain changed; this task cannot depend on the
persistence codecs introduced later. It validates the complete candidate
and rejects any declared/actual-domain mismatch.

Before evaluation, compute
\texttt{admittedDomains\ }\allowbreak{}\texttt{=}\allowbreak{}\texttt{\ }\allowbreak{}\texttt{mutation.}\allowbreak{}\texttt{RequiredDomains\ }\allowbreak{}\texttt{\textbar{}}\allowbreak{}\texttt{\ }\allowbreak{}\texttt{BankMutationDomains.}\allowbreak{}\texttt{Requests},
reject unknown bits, and prove none of those domains is quarantined or
unavailable. After evaluation and typed control-change application,
independently derive the intrinsic business-domain union: include
\passthrough{\lstinline!Finance!} when any journal draft exists, then
union every control change\textquotesingle s single declared domain.
Reject with \texttt{UndeclaredMutationDomain} unless both the
view\textquotesingle s observed-read domains and that intrinsic write
union are subsets of
\texttt{mutation.}\allowbreak{}\texttt{RequiredDomains}. Immediately
before publication, re-run the availability/quarantine check against the
observed-read and intrinsic-write union plus
\passthrough{\lstinline!Requests!} on the fully validated candidate.
This second check is mandatory defense in depth even though publication
remains under one lock.

\passthrough{\lstinline!BankDomainSlot<T>!} represents verified
authority or explicit unavailability; it never manufactures an empty
journal, registry, configuration, or other authoritative value after
corruption. \passthrough{\lstinline!TryGetVerified!} is the only
production value access. \passthrough{\lstinline!BankStateValidator!}
requires every unavailable slot to have a matching quarantine entry and
evidence digest and every verified slot to satisfy its domain validator.
This lets an unrelated verified domain advance while opaque corrupt
bytes remain outside the slot for exact persistence pass-through.

Ordinary mutations and control changes cannot alter
\passthrough{\lstinline!BankQuarantineState!} or
\passthrough{\lstinline!RecoveryControlState!}. Quarantine set/clear is
available only through \passthrough{\lstinline!ExecuteRecovery!}, an
authenticated permission-checked \texttt{AuthorizedRecoveryContext}, and
a registered sealed \passthrough{\lstinline!IBankRecoveryMutation!}. The
coordinator proves the administrator matches the authenticated caller,
the exact permission is granted, the mandatory reason is bounded, and
\passthrough{\lstinline!ExpectedEvidenceDigest!} matches the current
entry. A clear transition must install a fully validated repaired domain
value and move the byte-identical displaced raw evidence into the
authoritative \passthrough{\lstinline!RecoveryControlState!} archive
before removing the block; a set transition adds evidence and the block
atomically. Both append an immutable audit entry and advance the target
plus \passthrough{\lstinline!Recovery!} revisions. Stale evidence,
unregistered recovery types, missing permission, archive-size overflow,
or any attempt to clear without replacement rejects. If retaining both
repaired state and evidence would exceed persistence bounds, clearing
remains blocked until a later permission-gated export/prune workflow is
explicitly designed; evidence is never silently discarded.
\texttt{AuthorizedRecoveryContext} has an internal constructor visible
only to the trusted server adapter and core tests. The adapter
constructs it only after the engine\textquotesingle s permission API
accepts the authenticated server player; no network DTO carries or
selects \passthrough{\lstinline!GrantedPermission!}.
\passthrough{\lstinline!ExecuteRecovery!} still requires verified
\passthrough{\lstinline!RequestControl!} and
\passthrough{\lstinline!RecoveryControl!}, but it may admit its one
explicitly named unavailable target without exposing a fabricated target
value. Its recovery view supplies only the matching quarantine
metadata/evidence and the mutation-supplied replacement; all nontarget
domain reads remain gated normally. The replacement must pass the
target\textquotesingle s complete validator and any cross-domain checks
before the transition can clear the slot.

The coordinator obtains its one \passthrough{\lstinline!CommandId!}
directly from the injected \passthrough{\lstinline!IIdentifierSource!},
then gives \passthrough{\lstinline!Evaluate!} only a new
invocation-scoped tracking \texttt{IMutationIdentifierScope}. Every
factory receives an operation ID returned by
\passthrough{\lstinline!ReserveOperationId!}; an accepted decision must
use each reserved operation ID exactly once, may use no nonissued
operation ID, and may not duplicate one. Violation returns
\texttt{IdentifierContractViolation} and publishes no business state.
The coordinator finalizer copies each draft\textquotesingle s operation
ID unchanged into its record header and stores the same IDs, in
canonical draft order, in the terminal request or automatic-operation
entry and \passthrough{\lstinline!CoordinatorResponse!}. Task 10 extends
the tracker: each reserved settlement ID must be consumed by one
matching physical payload/control-change pair, and that
pair\textquotesingle s associated operation-ID array must equal the
exact finalized record IDs.

Under one private coordinator lock: pre-admit the authenticated request
or automatic root key; sort distinct participants with
\passthrough{\lstinline!StringComparer.Ordinal!} over the exact UID
value; capture one timestamp and command ID; create the tracking
identifier scope; evaluate one immutable decision; validate its
identifier and intrinsic-domain declarations; pre-admit every
subordinate automatic intent and reject any other request key; resolve
invocation intents; finalize sequences/hashes while preserving the
reserved operation IDs; replay/validate every new record and affected
projection; validate typed control changes; recheck candidate
quarantine; advance the revision vector; finalize the request and every
automatic replay barrier; then exchange the repository reference once.
No service may mutate a projection or journal outside this path.

Authenticated pre-admission first proves that
\passthrough{\lstinline!RequestKey.PlayerId!}, its scope owner, and
\passthrough{\lstinline!AuthenticatedPlayer!} are the same exact
\passthrough{\lstinline!PlayerId!}; a mismatch is preadmission failure
and never reaches the mutation.

If an admitted business request rejects, publish one valid zero-journal
control revision that advances Requests only. If a candidate has an
arithmetic/\allowbreak{}time/\allowbreak{}journal invariant fault,
discard it and build the terminal authenticated-request control revision
from the prior state. Preadmission failures such as inactive scope,
wrong sequence, changed digest, or a duplicate automatic key do not
advance state. Automatic invariant failures publish nothing; successful
automatic decisions install their permanent key atomically, even when
their validated financial result has zero journal records.

The listed coordinator support files contain deterministic fake
clocks/identifier sources and the two-record, rejecting, and
deliberately corrupt \passthrough{\lstinline!IBankMutation!}
implementations used by these tests. They do not duplicate coordinator
logic.

\begin{itemize}
\tightlist
\item[$\square$]
  \textbf{Step 7: Implement immutable snapshot and staging semantics}
\end{itemize}

\begin{lstlisting}
public sealed class BankStateRepository
{
    public bool HasUnstagedChanges { get; }
    public bool RequiresEngineSaveRestage { get; }
    public BankRevision? LastStagedRevision { get; }

    public BankStateSnapshot CaptureForStaging();
    public void Publish(BankState next);
    public void MarkStagingSucceeded(BankRevision snapshottedRevision);
    public void MarkStagingFailed(BankRevision snapshottedRevision, Exception exception);
}
\end{lstlisting}

\passthrough{\lstinline!CaptureForStaging!} returns the current
immutable state and revision without clearing
\passthrough{\lstinline!HasUnstagedChanges!}.
\passthrough{\lstinline!MarkStagingSucceeded!} clears that flag only
when the current global revision still equals the revision accepted by
both world-data staging calls. It records
\passthrough{\lstinline!LastStagedRevision!} but never calls it durably
saved. A separate \texttt{RequiresEngineSaveRestage} flag remains true
throughout the live host because Vintage Story exposes no post-disk-save
acknowledgment; the adapter restages the current snapshot on every
\passthrough{\lstinline!GameWorldSave!}. A staging failure or newer
publish retains \passthrough{\lstinline!HasUnstagedChanges!} as well.

\begin{itemize}
\tightlist
\item[$\square$]
  \textbf{Step 8: Verify coordinator/repository behavior under
  concurrency}
\end{itemize}

Run 100 parallel test calls with valid independently issued scopes and
assert journal sequences are contiguous, operation IDs unique, no lost
updates occur, and final projection equals replay. The test uses a
barrier to overlap callers but no wall-clock sleeps.

\begin{lstlisting}
dotnet test tests/FirstGearBank.Core.Tests/FirstGearBank.Core.Tests.csproj -c Release
\end{lstlisting}

\begin{itemize}
\tightlist
\item[$\square$]
  \textbf{Step 9: Run all state tests and commit}
\end{itemize}

\begin{lstlisting}
dotnet test tests/FirstGearBank.Core.Tests/FirstGearBank.Core.Tests.csproj -c Release
git diff --check
git add FirstGearBank.Core/State FirstGearBank.Core/Recovery FirstGearBank.Core/Requests `
  FirstGearBank.Core/Coordination tests/FirstGearBank.Core.Tests/Requests `
  tests/FirstGearBank.Core.Tests/Coordination tests/FirstGearBank.Core.Tests/Recovery `
  tests/FirstGearBank.Core.Tests/State `
  tests/FirstGearBank.Core.Tests/TestSupport/RequestTestData.cs `
  tests/FirstGearBank.Core.Tests/TestSupport/CoordinatorTestRig.cs `
  tests/FirstGearBank.Core.Tests/TestSupport/CoordinatorTestMutations.cs
git commit -m "feat: add atomic idempotent banking coordinator"
\end{lstlisting}

\subsubsection{Task 10: Freeze the Physical-Settlement Recovery Evidence
Contract}\label{task-10-freeze-the-physical-settlement-recovery-evidence-contract}

\textbf{Files:}

\begin{itemize}
\tightlist
\item
  Create:
  \texttt{FirstGearBank.}\allowbreak{}\texttt{Core/}\allowbreak{}\texttt{Settlement/}\allowbreak{}\texttt{PhysicalSettlementDirection.}\allowbreak{}\texttt{cs}
\item
  Create:
  \texttt{FirstGearBank.}\allowbreak{}\texttt{Core/}\allowbreak{}\texttt{Settlement/}\allowbreak{}\texttt{PhysicalSettlementPhase.}\allowbreak{}\texttt{cs}
\item
  Create:
  \texttt{FirstGearBank.}\allowbreak{}\texttt{Core/}\allowbreak{}\texttt{Settlement/}\allowbreak{}\texttt{SettlementManifestLine.}\allowbreak{}\texttt{cs}
\item
  Create:
  \texttt{FirstGearBank.}\allowbreak{}\texttt{Core/}\allowbreak{}\texttt{Settlement/}\allowbreak{}\texttt{SettlementManifest.}\allowbreak{}\texttt{cs}
\item
  Create:
  \texttt{FirstGearBank.}\allowbreak{}\texttt{Core/}\allowbreak{}\texttt{Settlement/}\allowbreak{}\texttt{InventoryFingerprint.}\allowbreak{}\texttt{cs}
\item
  Create:
  \texttt{FirstGearBank.}\allowbreak{}\texttt{Core/}\allowbreak{}\texttt{Settlement/}\allowbreak{}\texttt{InventorySlotAddress.}\allowbreak{}\texttt{cs}
\item
  Create:
  \texttt{FirstGearBank.}\allowbreak{}\texttt{Core/}\allowbreak{}\texttt{Settlement/}\allowbreak{}\texttt{InventorySlotDeltaV1.}\allowbreak{}\texttt{cs}
\item
  Create:
  \texttt{FirstGearBank.}\allowbreak{}\texttt{Core/}\allowbreak{}\texttt{Settlement/}\allowbreak{}\texttt{InventoryMutationPlanV1.}\allowbreak{}\texttt{cs}
\item
  Create:
  \texttt{FirstGearBank.}\allowbreak{}\texttt{Core/}\allowbreak{}\texttt{Settlement/}\allowbreak{}\texttt{PhysicalSettlementLimits.}\allowbreak{}\texttt{cs}
\item
  Create:
  \texttt{FirstGearBank.}\allowbreak{}\texttt{Core/}\allowbreak{}\texttt{Settlement/}\allowbreak{}\texttt{PhysicalSettlementRecoveryPlanV1.}\allowbreak{}\texttt{cs}
\item
  Create:
  \texttt{FirstGearBank.}\allowbreak{}\texttt{Core/}\allowbreak{}\texttt{Settlement/}\allowbreak{}\texttt{PhysicalSettlementRecoveryPlanCodec.}\allowbreak{}\texttt{cs}
\item
  Create:
  \texttt{FirstGearBank.}\allowbreak{}\texttt{Core/}\allowbreak{}\texttt{Settlement/}\allowbreak{}\texttt{PhysicalSettlementDomainState.}\allowbreak{}\texttt{cs}
\item
  Create:
  \texttt{FirstGearBank.}\allowbreak{}\texttt{Core/}\allowbreak{}\texttt{Settlement/}\allowbreak{}\texttt{PhysicalSettlementRecord.}\allowbreak{}\texttt{cs}
\item
  Create:
  \texttt{FirstGearBank.}\allowbreak{}\texttt{Core/}\allowbreak{}\texttt{Settlement/}\allowbreak{}\texttt{SettlementRecoveryCapsule.}\allowbreak{}\texttt{cs}
\item
  Create:
  \texttt{FirstGearBank.}\allowbreak{}\texttt{Core/}\allowbreak{}\texttt{Settlement/}\allowbreak{}\texttt{RecoveryAuthenticationKey.}\allowbreak{}\texttt{cs}
\item
  Create:
  \texttt{FirstGearBank.}\allowbreak{}\texttt{Core/}\allowbreak{}\texttt{Settlement/}\allowbreak{}\texttt{SettlementCapsuleAuthenticator.}\allowbreak{}\texttt{cs}
\item
  Create:
  \texttt{FirstGearBank.}\allowbreak{}\texttt{Core/}\allowbreak{}\texttt{Settlement/}\allowbreak{}\texttt{SettlementRecoveryEvidence.}\allowbreak{}\texttt{cs}
\item
  Create:
  \texttt{FirstGearBank.}\allowbreak{}\texttt{Core/}\allowbreak{}\texttt{Settlement/}\allowbreak{}\texttt{SettlementRecoveryDecision.}\allowbreak{}\texttt{cs}
\item
  Create:
  \texttt{FirstGearBank.}\allowbreak{}\texttt{Core/}\allowbreak{}\texttt{Settlement/}\allowbreak{}\texttt{SettlementReconciler.}\allowbreak{}\texttt{cs}
\item
  Create:
  \texttt{FirstGearBank.}\allowbreak{}\texttt{Core/}\allowbreak{}\texttt{Settlement/}\allowbreak{}\texttt{VerifiedPhysicalRecoveryContext.}\allowbreak{}\texttt{cs}
\item
  Create:
  \texttt{FirstGearBank.}\allowbreak{}\texttt{Core/}\allowbreak{}\texttt{Settlement/}\allowbreak{}\texttt{PhysicalSettlementControlChange.}\allowbreak{}\texttt{cs}
\item
  Create:
  \texttt{FirstGearBank.}\allowbreak{}\texttt{Core/}\allowbreak{}\texttt{Settlement/}\allowbreak{}\texttt{IPlayerRecoveryCapsuleStore.}\allowbreak{}\texttt{cs}
\item
  Create:
  \texttt{tests/}\allowbreak{}\texttt{FirstGearBank.}\allowbreak{}\texttt{Core.}\allowbreak{}\texttt{Tests/}\allowbreak{}\texttt{Settlement/}\allowbreak{}\texttt{SettlementCapsuleAuthenticatorTests.}\allowbreak{}\texttt{cs}
\item
  Create:
  \texttt{tests/}\allowbreak{}\texttt{FirstGearBank.}\allowbreak{}\texttt{Core.}\allowbreak{}\texttt{Tests/}\allowbreak{}\texttt{Settlement/}\allowbreak{}\texttt{PhysicalSettlementRecoveryPlanCodecTests.}\allowbreak{}\texttt{cs}
\item
  Create:
  \texttt{tests/}\allowbreak{}\texttt{FirstGearBank.}\allowbreak{}\texttt{Core.}\allowbreak{}\texttt{Tests/}\allowbreak{}\texttt{Settlement/}\allowbreak{}\texttt{SettlementRecoveryMatrixTests.}\allowbreak{}\texttt{cs}
\item
  Create:
  \texttt{tests/}\allowbreak{}\texttt{FirstGearBank.}\allowbreak{}\texttt{Core.}\allowbreak{}\texttt{Tests/}\allowbreak{}\texttt{Settlement/}\allowbreak{}\texttt{CapsuleOnlyCoordinatorRecoveryTests.}\allowbreak{}\texttt{cs}
\item
  Create:
  \texttt{tests/}\allowbreak{}\texttt{FirstGearBank.}\allowbreak{}\texttt{Core.}\allowbreak{}\texttt{Tests/}\allowbreak{}\texttt{Settlement/}\allowbreak{}\texttt{SettlementContractBoundsTests.}\allowbreak{}\texttt{cs}
\item
  Create:
  \texttt{tests/}\allowbreak{}\texttt{FirstGearBank.}\allowbreak{}\texttt{Core.}\allowbreak{}\texttt{Tests/}\allowbreak{}\texttt{TestSupport/}\allowbreak{}\texttt{SettlementTestData.}\allowbreak{}\texttt{cs}
\item
  Create:
  \texttt{tests/}\allowbreak{}\texttt{FirstGearBank.}\allowbreak{}\texttt{Core.}\allowbreak{}\texttt{Tests/}\allowbreak{}\texttt{TestSupport/}\allowbreak{}\texttt{Evidence.}\allowbreak{}\texttt{cs}
\item
  Create:
  \texttt{tests/}\allowbreak{}\texttt{FirstGearBank.}\allowbreak{}\texttt{Core.}\allowbreak{}\texttt{Tests/}\allowbreak{}\texttt{TestSupport/}\allowbreak{}\texttt{TestKeys.}\allowbreak{}\texttt{cs}
\item
  Modify:
  \texttt{FirstGearBank.}\allowbreak{}\texttt{Core/}\allowbreak{}\texttt{State/}\allowbreak{}\texttt{BankState.}\allowbreak{}\texttt{cs:}\allowbreak{}\texttt{1}
\item
  Modify:
  \texttt{FirstGearBank.}\allowbreak{}\texttt{Core/}\allowbreak{}\texttt{State/}\allowbreak{}\texttt{BankStateValidator.}\allowbreak{}\texttt{cs:}\allowbreak{}\texttt{1}
\item
  Modify:
  \texttt{FirstGearBank.}\allowbreak{}\texttt{Core/}\allowbreak{}\texttt{Requests/}\allowbreak{}\texttt{RequestRegistry.}\allowbreak{}\texttt{cs:}\allowbreak{}\texttt{1}
\item
  Modify:
  \texttt{FirstGearBank.}\allowbreak{}\texttt{Core/}\allowbreak{}\texttt{Coordination/}\allowbreak{}\texttt{BankingCoordinator.}\allowbreak{}\texttt{cs:}\allowbreak{}\texttt{1}
\item
  Modify:
  \texttt{tests/}\allowbreak{}\texttt{FirstGearBank.}\allowbreak{}\texttt{Core.}\allowbreak{}\texttt{Tests/}\allowbreak{}\texttt{Coordination/}\allowbreak{}\texttt{BankingCoordinatorTests.}\allowbreak{}\texttt{cs:}\allowbreak{}\texttt{1}
\item
  Modify:
  \texttt{docs/}\allowbreak{}\texttt{architecture/}\allowbreak{}\texttt{0001-}\allowbreak{}\texttt{physical-}\allowbreak{}\texttt{settlement-}\allowbreak{}\texttt{recovery-}\allowbreak{}\texttt{threat-}\allowbreak{}\texttt{model.}\allowbreak{}\texttt{md:}\allowbreak{}\texttt{1}
\item
  Modify:
  \texttt{docs/}\allowbreak{}\texttt{architecture/}\allowbreak{}\texttt{0001-}\allowbreak{}\texttt{physical-}\allowbreak{}\texttt{settlement-}\allowbreak{}\texttt{recovery-}\allowbreak{}\texttt{threat-}\allowbreak{}\texttt{model.}\allowbreak{}\texttt{tex:}\allowbreak{}\texttt{1}
  through regeneration
\item
  Generate but do not commit:
  \texttt{docs/}\allowbreak{}\texttt{architecture/}\allowbreak{}\texttt{0001-}\allowbreak{}\texttt{physical-}\allowbreak{}\texttt{settlement-}\allowbreak{}\texttt{recovery-}\allowbreak{}\texttt{threat-}\allowbreak{}\texttt{model.}\allowbreak{}\texttt{pdf}
\end{itemize}

\textbf{Interfaces:}

\begin{itemize}
\tightlist
\item
  Consumes: canonical IDs, money, hashes, and revision values from Tasks
  5 and 9.
\item
  Produces: the authenticated record/capsule formats and exact
  reconciliation decision matrix persisted by Task 11 and implemented
  against inventory in the physical-settlement plan.
\end{itemize}

\begin{itemize}
\tightlist
\item[$\square$]
  \textbf{Step 1: Write failing capsule roundtrip and tamper tests}
\end{itemize}

\begin{lstlisting}
[Fact]
public void Capsule_authenticates_only_for_the_exact_world_player_and_manifest()
{
    using var key = RecoveryAuthenticationKey.FromBytes(TestKeys.ThirtyTwoBytes);
    var capsule = SettlementTestData.UnsignedCapsule();
    var signed = SettlementCapsuleAuthenticator.Sign(key, capsule);

    Assert.True(SettlementCapsuleAuthenticator.Verify(key, signed));
    Assert.False(SettlementCapsuleAuthenticator.Verify(
        key,
        signed with { Plan = signed.Plan with { PlayerId = PlayerIds.Byron } }));
    Assert.False(SettlementCapsuleAuthenticator.Verify(
        key,
        signed with
        {
            Plan = signed.Plan with { Manifest = SettlementTestData.DifferentManifest }
        }));
}
\end{lstlisting}

Independently tamper every recovery-plan field: world binding, player,
settlement ID, request key, request payload digest, command ID,
financial/world timestamp, direction, manifest line order/content,
before/after fingerprint, inventory address, slot index, before/after
stack bytes, draft
operation/\allowbreak{}intent/\allowbreak{}schema/\allowbreak{}party/\allowbreak{}payload/\allowbreak{}postings,
associated operation IDs, terminal response, and phase. Test a wrong
key, short/long key, short/long tag, duplicate item lines/slot
addresses, overlong asset or inventory code, zero/negative count,
invalid stack-codec version, arithmetic overflow, and every exact size
boundary.

\begin{itemize}
\tightlist
\item[$\square$]
  \textbf{Step 2: Run capsule tests and observe missing-type failures}
\end{itemize}

\begin{lstlisting}
dotnet test tests/FirstGearBank.Core.Tests/FirstGearBank.Core.Tests.csproj -c Release
\end{lstlisting}

Expected: FAIL because the settlement evidence types do not exist.

\begin{itemize}
\tightlist
\item[$\square$]
  \textbf{Step 3: Implement bounded canonical evidence and HMAC
  authentication}
\end{itemize}

Use only the four persisted phases locked by section 10.1:

\begin{lstlisting}
public enum PhysicalSettlementPhase : byte
{
    Prepared = 1,
    InventoryApplied = 2,
    BankCommitted = 3,
    Finalized = 4
}

public sealed record InventorySlotAddress(
    string InventoryCode,
    int SlotIndex);

public sealed record InventorySlotDeltaV1(
    InventorySlotAddress Address,
    ImmutableArray<byte> CanonicalBeforeStack,
    ImmutableArray<byte> CanonicalAfterStack);

public sealed record InventoryMutationPlanV1(
    ushort StackCodecVersion,
    ImmutableArray<InventorySlotDeltaV1> SlotDeltas);

public static class PhysicalSettlementLimits
{
    public const int MaximumJournalDrafts = 16;
    public const int MaximumInventorySlotDeltas = 64;
    public const int MaximumCanonicalStackBytes = 64 * 1024;
    public const int MaximumRecoveryPlanBytes = 4 * 1024 * 1024;
}

public sealed record PhysicalSettlementRecoveryPlanV1(
    ushort PlanSchemaVersion,
    EvidenceDigest WorldBindingHash,
    PlayerId PlayerId,
    SettlementId SettlementId,
    RequestKey RequestKey,
    CommandPayloadDigest PayloadDigest,
    CommandId CommandId,
    FinancialTimestamp Timestamp,
    PhysicalSettlementDirection Direction,
    SettlementManifest Manifest,
    InventoryFingerprint BeforeStateFingerprint,
    InventoryFingerprint AfterStateFingerprint,
    InventoryMutationPlanV1 InventoryMutation,
    ImmutableArray<JournalDraft> PredeclaredBankDrafts,
    ImmutableArray<OperationId> AssociatedOperationIds,
    TerminalResponse SuccessfulTerminalResponse);

public sealed record SettlementRecoveryCapsule(
    ushort CapsuleSchemaVersion,
    PhysicalSettlementRecoveryPlanV1 Plan,
    PhysicalSettlementPhase ProvenPhase,
    ImmutableArray<byte> AuthenticationTag);
\end{lstlisting}

\texttt{PhysicalSettlementRecord} stores this same complete recovery
plan plus phase; a capsule never points at Bank-only data.
\texttt{PhysicalSettlementRecoveryPlanCodec} encodes with
\texttt{FGB/}\allowbreak{}\texttt{PHYSICAL-}\allowbreak{}\texttt{SETTLEMENT-}\allowbreak{}\texttt{RECOVERY-}\allowbreak{}\texttt{PLAN/}\allowbreak{}\texttt{1\textbackslash{}}\allowbreak{}\texttt{0},
explicit numeric widths, length-prefixed UTF-8/bytes, sorted unique slot
addresses, and the Task 7 canonical encodings of every complete draft
field. Empty stack bytes mean an empty slot. The phase-3 adapter must
prove a stable player-inventory address and versioned canonical
Bank-item stack codec before it may advertise co-serialization or
automatically apply a slot delta; until then capability remains
\passthrough{\lstinline!UnprovedCoSerialization!}.

Validate the plan as one closed contract: request/player agree; every
draft is \passthrough{\lstinline!Caller!}, shares the exact stored
timestamp, uses a physical payload with the same
settlement/\allowbreak{}player/\allowbreak{}direction/\allowbreak{}manifest
facts, and is independently balanced; its associated operation-ID array
equals the draft IDs in canonical order; and its bounded terminal
response is successful and describes those exact IDs. The stored command
ID and request digest are the original admitted values. Sequence
numbers, previous hashes, record hashes, and resulting revision are
deliberately absent because the coordinator assigns them against the
verified current Bank state during initial commit or recovery.

Hash the canonical draft array with
\texttt{FGB/}\allowbreak{}\texttt{PHYSICAL-}\allowbreak{}\texttt{SETTLEMENT-}\allowbreak{}\texttt{BANK-}\allowbreak{}\texttt{BATCH/}\allowbreak{}\texttt{1\textbackslash{}}\allowbreak{}\texttt{0};
expose that derived digest for diagnostics and cross-checks, but retain
the drafts themselves because a digest is not reconstructive evidence.
Canonicalize the outer capsule with
\texttt{FGB/}\allowbreak{}\texttt{SETTLEMENT-}\allowbreak{}\texttt{CAPSULE/}\allowbreak{}\texttt{1\textbackslash{}}\allowbreak{}\texttt{0}
over schema, complete plan bytes, and phase, excluding the
authentication-tag field itself. Generate a random 32-byte key through
the Task 5 \passthrough{\lstinline!IEntropySource!}; production uses
\texttt{RandomNumberGenerator.}\allowbreak{}\texttt{Fill}. Never expose
it in logs, clear owned key memory on dispose, sign with HMAC-SHA-256,
and verify with
\texttt{CryptographicOperations.}\allowbreak{}\texttt{FixedTimeEquals}.

Bound a plan to 16 journal drafts, 64 slot deltas, 64 KiB per canonical
stack, and 4 MiB total canonical plan bytes. Reject limits before
allocating or iterating nested payloads. Golden fixtures independently
hand-code one deposit and one withdrawal plan and assert byte-identical
plan encoding, batch digest, capsule bytes, and authentication tag.

\begin{itemize}
\tightlist
\item[$\square$]
  \textbf{Step 4: Verify capsule and bounds tests pass}
\end{itemize}

\begin{lstlisting}
dotnet test tests/FirstGearBank.Core.Tests/FirstGearBank.Core.Tests.csproj -c Release
\end{lstlisting}

\begin{itemize}
\tightlist
\item[$\square$]
  \textbf{Step 5: Write the exact failing reconciliation-matrix tests}
\end{itemize}

\begin{lstlisting}
[Theory]
[MemberData(nameof(RecoveryCases))]
public void Recovery_matrix_has_one_conservative_action(
    SettlementRecoveryEvidence evidence,
    SettlementRecoveryAction expected)
{
    Assert.Equal(expected, SettlementReconciler.Decide(evidence).Action);
}

public static TheoryData<SettlementRecoveryEvidence, SettlementRecoveryAction> RecoveryCases => new()
{
    { Evidence.NeitherSide, SettlementRecoveryAction.DiscardOrphanPrepared },
    { Evidence.ProvedCoSerializedCapsuleOnly, SettlementRecoveryAction.AppendPredeclaredBankBatch },
    { Evidence.UnprovedCoSerializedCapsuleOnly, SettlementRecoveryAction.Quarantine },
    { Evidence.BankOnlyFingerprintMatches, SettlementRecoveryAction.ApplyInventoryAndCapsule },
    { Evidence.BothSides, SettlementRecoveryAction.MarkFinalized },
    { Evidence.BankOnlyFingerprintChanged, SettlementRecoveryAction.Quarantine },
    { Evidence.InvalidCapsuleOnly, SettlementRecoveryAction.Quarantine },
    { Evidence.ContradictoryPhases, SettlementRecoveryAction.Quarantine }
};

[Fact]
public void Proved_capsule_only_recovery_commits_the_exact_plan_once_across_restart()
{
    var rig = CapsuleOnlyRecoveryRig.WithNoBankPreparedRecord();
    var verified = rig.VerifyCapsule(RecoveryCapsuleCapability.ProvedCoSerialization).Value;

    var first = rig.Coordinator.RecoverPhysical(verified);

    Assert.True(first.IsSuccess);
    Assert.True(rig.Plan.AssociatedOperationIds.AsSpan().SequenceEqual(first.OperationIds.AsSpan()));
    var committed = rig.State.Finance.TryGetVerified().Value.Journal[^1];
    Assert.Equal(rig.Plan.CommandId, committed.CommandId);
    Assert.Equal(rig.Plan.Timestamp, committed.Timestamp);
    Assert.True(rig.Plan.PredeclaredBankDrafts[0].Parties.AsSpan().SequenceEqual(committed.Parties.AsSpan()));
    Assert.True(rig.HasPermanentRequestBarrier(rig.Plan.RequestKey, rig.Plan.PayloadDigest));

    var restored = rig.SaveAndRestore();
    var duplicate = restored.Coordinator.RecoverPhysical(restored.VerifySameCapsule().Value);

    Assert.True(first.OperationIds.AsSpan().SequenceEqual(duplicate.OperationIds.AsSpan()));
    Assert.Single(restored.State.Finance.TryGetVerified().Value.Journal);
    Assert.Equal(0, restored.TimestampCaptureCount);
    Assert.Equal(0, restored.IdentifierIssueCount);
}
\end{lstlisting}

The evidence fixture carries the adapter\textquotesingle s
\texttt{RecoveryCapsuleCapability}; a cryptographically valid capsule is
not sufficient for capsule-only completion while that capability is
\passthrough{\lstinline!UnprovedCoSerialization!}. Add the
withdrawal-capacity case: a valid Bank-only withdrawal whose exact
output no longer fits returns \texttt{WaitForOriginalDeliveryCapacity},
retaining the same \passthrough{\lstinline!SettlementId!} and never
proposing another debit. Add a coordinator-integrated contract case in
which one physical mutation reserves its settlement and operation IDs,
emits the journal draft plus \texttt{PhysicalSettlementControlChange},
and proves the persisted record header, settlement record, and capsule
\passthrough{\lstinline!AssociatedOperationIds!} contain the same
operation ID. Mutate each association independently: an unreserved
settlement ID, reordered/\allowbreak{}extra/\allowbreak{}missing
operation ID, payload/control settlement mismatch, or a second consumer
of either ID must return \texttt{IdentifierContractViolation} without
publishing finance or settlement state.

\begin{lstlisting}
dotnet test tests/FirstGearBank.Core.Tests/FirstGearBank.Core.Tests.csproj -c Release
\end{lstlisting}

Expected: FAIL because the conservative reconciliation decision table is
not implemented.

\begin{itemize}
\tightlist
\item[$\square$]
  \textbf{Step 6: Implement a pure decision table with no inventory
  heuristics}
\end{itemize}

\begin{lstlisting}
public static class SettlementReconciler
{
    public static SettlementRecoveryDecision Decide(SettlementRecoveryEvidence evidence);
}

public interface IPlayerRecoveryCapsuleStore
{
    RecoveryCapsuleCapability Capability { get; }
    BankResult<SettlementRecoveryCapsule?> Read(PlayerId playerId, SettlementId settlementId);
    BankResult<bool> Stage(PlayerId playerId, SettlementRecoveryCapsule capsule);
    BankResult<bool> Remove(PlayerId playerId, SettlementId settlementId);
}

public static class SettlementCapsuleAuthenticator
{
    public static BankResult<VerifiedPhysicalRecoveryContext> VerifyAndOpen(
        RecoveryAuthenticationKey key,
        SettlementRecoveryCapsule capsule,
        EvidenceDigest expectedWorldBinding,
        PlayerId expectedPlayer,
        RecoveryCapsuleCapability capability);
}

public sealed class BankingCoordinator
{
    public CoordinatorResponse RecoverPhysical(VerifiedPhysicalRecoveryContext context);
}
\end{lstlisting}

The decision table consumes authenticated evidence flags; it never
examines fungible stack totals. An unavailable capsule store reports
\passthrough{\lstinline!UnprovedCoSerialization!}, which forces
ambiguous one-sided evidence to quarantine. It must not masquerade as a
valid absent capsule.

\texttt{VerifiedPhysicalRecoveryContext} has no public constructor and
exposes no mutable plan. \passthrough{\lstinline!VerifyAndOpen!} returns
it only after HMAC, canonical-plan validation, expected world/player
binding, and exact \passthrough{\lstinline!ProvedCoSerialization!}
capability all pass. \passthrough{\lstinline!RecoverPhysical!} is a
closed coordinator path, not a general mutation hook: it accepts only
physical draft types from that context, uses the plan\textquotesingle s
original command
ID/\allowbreak{}timestamp/\allowbreak{}digest/\allowbreak{}parties/\allowbreak{}payload/\allowbreak{}postings/\allowbreak{}operation
IDs, and assigns only current
sequence/\allowbreak{}hash/\allowbreak{}revision fields. It never opens
an operation scope and never captures new time or identifiers.

For capsule-only evidence, the coordinator checks that neither a
conflicting settlement nor permanent request entry exists, appends the
predeclared batch, installs the matching settlement tombstone, and calls
an internal
\texttt{RequestRegistry.}\allowbreak{}\texttt{FinalizeRecoveredSuccess}
that can install the exact original
\passthrough{\lstinline!RequestKey!}, payload digest, terminal body,
operation IDs, and actual resulting revision without an active scope.
That internal method is callable only from this verified recovery path.
An identical existing barrier/tombstone returns its stored result; any
digest, plan, response, or operation-ID disagreement quarantines the
settlement. The end-to-end test starts with no Bank-side
\passthrough{\lstinline!Prepared!} record and retries after another
save/restart, proving the capsule contains reconstructive data rather
than only a digest.

\texttt{PhysicalSettlementControlChange} is the first production
\passthrough{\lstinline!IBankControlChange!}. It replaces only the
immutable settlement map, declares
\texttt{BankMutationDomains.}\allowbreak{}\texttt{Settlements}, carries
the complete recovery plan and the same caller intent as its physical
journal draft, and is rejected if any other state domain differs. Tests
prove it can join a finance/request revision atomically without gaining
a path to mutate journal, projections, replay registries, or revisions
directly.

This task extends \passthrough{\lstinline!BankState!} with
\texttt{BankDomainSlot\textless{}}\allowbreak{}\texttt{PhysicalSettlementDomainState\textgreater{}}\allowbreak{}\texttt{\ }\allowbreak{}\texttt{Settlements}
and adds the gated
\texttt{BankMutationReadView.}\allowbreak{}\texttt{ReadSettlements(}\allowbreak{}\texttt{)}\allowbreak{}
accessor. A corrupt settlement section therefore has no fabricated empty
map and blocks all player banking as required, while verified unrelated
read-only state remains available.

Extend the Task 9 tracking scope validation without changing draft IDs:
each accepted physical settlement consumes exactly one ID returned by
\passthrough{\lstinline!ReserveSettlementId!}, and its control record
identifies the exact physical draft set for that settlement. The capsule
and control record\textquotesingle s
\passthrough{\lstinline!AssociatedOperationIds!} must equal, in
canonical draft order, the operation IDs that the coordinator copies
into the finalized journal headers. No later phase may regenerate,
substitute, or infer these IDs from array position.

\passthrough{\lstinline!SettlementTestData!},
\passthrough{\lstinline!Evidence!}, and
\passthrough{\lstinline!TestKeys!} provide deterministic bounded
fixtures only. The production key generator is never replaced or seeded
outside tests.

\begin{itemize}
\tightlist
\item[$\square$]
  \textbf{Step 7: Record the public-API gate result honestly}
\end{itemize}

Update the threat model with:

\begin{lstlisting}
Gate result: the 1.22.7 public API exposes IServerPlayer.SetModData/GetModData and world-save blobs, but exposes no contract
proving that player mod data and inventory are durably co-serialized or ordered against the Bank world blob.  The shipping
capability therefore defaults to UnprovedCoSerialization.  A controlled engine crash matrix may promote the adapter only
after repeatable evidence; until then, every ambiguous one-sided cross-blob state quarantines as specified.
\end{lstlisting}

This is a supported conservative outcome under section 10.1, not a claim
of disk atomicity. The threat model must also record that
\texttt{PhysicalSettlementRecoveryPlanV1}, not its digest, is the
reconstructive authority; that
sequence/\allowbreak{}hash/\allowbreak{}revision are assigned only when
the plan is committed; and that the inactive-scope recovery exception is
available solely through an HMAC-verified, world/player-bound capsule
and the closed \passthrough{\lstinline!RecoverPhysical!} coordinator
path.

Regenerate the threat model\textquotesingle s same-basename LaTeX and
PDF from Markdown, render every PDF page, and verify the revised gate
text before staging the \passthrough{\lstinline!.md!} and
\passthrough{\lstinline!.tex!}. Do not edit either generated format
directly.

\begin{itemize}
\tightlist
\item[$\square$]
  \textbf{Step 8: Verify the recovery matrix and commit}
\end{itemize}

\begin{lstlisting}
dotnet test tests/FirstGearBank.Core.Tests/FirstGearBank.Core.Tests.csproj -c Release
git diff --check
git add FirstGearBank.Core/Settlement FirstGearBank.Core/State `
  FirstGearBank.Core/Requests/RequestRegistry.cs `
  FirstGearBank.Core/Coordination/BankingCoordinator.cs `
  tests/FirstGearBank.Core.Tests/Settlement `
  tests/FirstGearBank.Core.Tests/Coordination/BankingCoordinatorTests.cs `
  tests/FirstGearBank.Core.Tests/TestSupport/SettlementTestData.cs `
  tests/FirstGearBank.Core.Tests/TestSupport/Evidence.cs `
  tests/FirstGearBank.Core.Tests/TestSupport/TestKeys.cs `
  docs/architecture/0001-physical-settlement-recovery-threat-model.md `
  docs/architecture/0001-physical-settlement-recovery-threat-model.tex
git commit -m "feat: define authenticated settlement recovery evidence"
\end{lstlisting}

\subsubsection{Task 11: Persist Independently Checksummed Sections with
Migrations and
Quarantine}\label{task-11-persist-independently-checksummed-sections-with-migrations-and-quarantine}

\textbf{Files:}

\begin{itemize}
\tightlist
\item
  Create:
  \texttt{FirstGearBank.}\allowbreak{}\texttt{Core/}\allowbreak{}\texttt{Persistence/}\allowbreak{}\texttt{BankSectionKind.}\allowbreak{}\texttt{cs}
\item
  Create:
  \texttt{FirstGearBank.}\allowbreak{}\texttt{Core/}\allowbreak{}\texttt{Persistence/}\allowbreak{}\texttt{BankSectionDescriptor.}\allowbreak{}\texttt{cs}
\item
  Create:
  \texttt{FirstGearBank.}\allowbreak{}\texttt{Core/}\allowbreak{}\texttt{Persistence/}\allowbreak{}\texttt{BankEnvelopeManifest.}\allowbreak{}\texttt{cs}
\item
  Create:
  \texttt{FirstGearBank.}\allowbreak{}\texttt{Core/}\allowbreak{}\texttt{Persistence/}\allowbreak{}\texttt{DecodedBankEnvelope.}\allowbreak{}\texttt{cs}
\item
  Create:
  \texttt{FirstGearBank.}\allowbreak{}\texttt{Core/}\allowbreak{}\texttt{Persistence/}\allowbreak{}\texttt{QuarantinedSection.}\allowbreak{}\texttt{cs}
\item
  Create:
  \texttt{FirstGearBank.}\allowbreak{}\texttt{Core/}\allowbreak{}\texttt{Persistence/}\allowbreak{}\texttt{SectionQuarantineReason.}\allowbreak{}\texttt{cs}
\item
  Create:
  \texttt{FirstGearBank.}\allowbreak{}\texttt{Core/}\allowbreak{}\texttt{Persistence/}\allowbreak{}\texttt{BankEnvelopeCodec.}\allowbreak{}\texttt{cs}
\item
  Create:
  \texttt{FirstGearBank.}\allowbreak{}\texttt{Core/}\allowbreak{}\texttt{Persistence/}\allowbreak{}\texttt{BankPersistenceLimits.}\allowbreak{}\texttt{cs}
\item
  Create:
  \texttt{FirstGearBank.}\allowbreak{}\texttt{Core/}\allowbreak{}\texttt{Persistence/}\allowbreak{}\texttt{CanonicalJson.}\allowbreak{}\texttt{cs}
\item
  Create:
  \texttt{FirstGearBank.}\allowbreak{}\texttt{Core/}\allowbreak{}\texttt{Persistence/}\allowbreak{}\texttt{ISectionCodec.}\allowbreak{}\texttt{cs}
\item
  Create:
  \texttt{FirstGearBank.}\allowbreak{}\texttt{Core/}\allowbreak{}\texttt{Persistence/}\allowbreak{}\texttt{ISectionMigration.}\allowbreak{}\texttt{cs}
\item
  Create:
  \texttt{FirstGearBank.}\allowbreak{}\texttt{Core/}\allowbreak{}\texttt{Persistence/}\allowbreak{}\texttt{SectionMigrationRegistry.}\allowbreak{}\texttt{cs}
\item
  Create:
  \texttt{FirstGearBank.}\allowbreak{}\texttt{Core/}\allowbreak{}\texttt{Persistence/}\allowbreak{}\texttt{BankStateSerializer.}\allowbreak{}\texttt{cs}
\item
  Create:
  \texttt{FirstGearBank.}\allowbreak{}\texttt{Core/}\allowbreak{}\texttt{Persistence/}\allowbreak{}\texttt{BankStateLoader.}\allowbreak{}\texttt{cs}
\item
  Create:
  \texttt{FirstGearBank.}\allowbreak{}\texttt{Core/}\allowbreak{}\texttt{Persistence/}\allowbreak{}\texttt{BankLoadOutcome.}\allowbreak{}\texttt{cs}
\item
  Create:
  \texttt{FirstGearBank.}\allowbreak{}\texttt{Core/}\allowbreak{}\texttt{Persistence/}\allowbreak{}\texttt{InstallMarker.}\allowbreak{}\texttt{cs}
\item
  Create:
  \texttt{FirstGearBank.}\allowbreak{}\texttt{Core/}\allowbreak{}\texttt{Persistence/}\allowbreak{}\texttt{InstallMarkerCodec.}\allowbreak{}\texttt{cs}
\item
  Create:
  \texttt{FirstGearBank.}\allowbreak{}\texttt{Core/}\allowbreak{}\texttt{Persistence/}\allowbreak{}\texttt{IBankWorldStateStore.}\allowbreak{}\texttt{cs}
\item
  Create:
  \texttt{FirstGearBank.}\allowbreak{}\texttt{Core/}\allowbreak{}\texttt{Persistence/}\allowbreak{}\texttt{BankPersistenceService.}\allowbreak{}\texttt{cs}
\item
  Create:
  \texttt{FirstGearBank.}\allowbreak{}\texttt{Core/}\allowbreak{}\texttt{Persistence/}\allowbreak{}\texttt{Sections/}\allowbreak{}\texttt{WorldAuthoritySectionV1.}\allowbreak{}\texttt{cs}
\item
  Create:
  \texttt{FirstGearBank.}\allowbreak{}\texttt{Core/}\allowbreak{}\texttt{Persistence/}\allowbreak{}\texttt{Sections/}\allowbreak{}\texttt{JournalSectionV1.}\allowbreak{}\texttt{cs}
\item
  Create:
  \texttt{FirstGearBank.}\allowbreak{}\texttt{Core/}\allowbreak{}\texttt{Persistence/}\allowbreak{}\texttt{Sections/}\allowbreak{}\texttt{ProjectionCacheSectionV1.}\allowbreak{}\texttt{cs}
\item
  Create:
  \texttt{FirstGearBank.}\allowbreak{}\texttt{Core/}\allowbreak{}\texttt{Persistence/}\allowbreak{}\texttt{Sections/}\allowbreak{}\texttt{RequestControlSectionV1.}\allowbreak{}\texttt{cs}
\item
  Create:
  \texttt{FirstGearBank.}\allowbreak{}\texttt{Core/}\allowbreak{}\texttt{Persistence/}\allowbreak{}\texttt{Sections/}\allowbreak{}\texttt{ConfigurationSectionV1.}\allowbreak{}\texttt{cs}
\item
  Create:
  \texttt{FirstGearBank.}\allowbreak{}\texttt{Core/}\allowbreak{}\texttt{Persistence/}\allowbreak{}\texttt{Sections/}\allowbreak{}\texttt{RecipientRegistrySectionV1.}\allowbreak{}\texttt{cs}
\item
  Create:
  \texttt{FirstGearBank.}\allowbreak{}\texttt{Core/}\allowbreak{}\texttt{Persistence/}\allowbreak{}\texttt{Sections/}\allowbreak{}\texttt{DeliveryControlSectionV1.}\allowbreak{}\texttt{cs}
\item
  Create:
  \texttt{FirstGearBank.}\allowbreak{}\texttt{Core/}\allowbreak{}\texttt{Persistence/}\allowbreak{}\texttt{Sections/}\allowbreak{}\texttt{PhysicalSettlementSectionV1.}\allowbreak{}\texttt{cs}
\item
  Create:
  \texttt{FirstGearBank.}\allowbreak{}\texttt{Core/}\allowbreak{}\texttt{Persistence/}\allowbreak{}\texttt{Sections/}\allowbreak{}\texttt{BranchControlSectionV1.}\allowbreak{}\texttt{cs}
\item
  Create:
  \texttt{FirstGearBank.}\allowbreak{}\texttt{Core/}\allowbreak{}\texttt{Persistence/}\allowbreak{}\texttt{Sections/}\allowbreak{}\texttt{DiagnosticsCacheSectionV1.}\allowbreak{}\texttt{cs}
\item
  Create:
  \texttt{FirstGearBank.}\allowbreak{}\texttt{Core/}\allowbreak{}\texttt{Persistence/}\allowbreak{}\texttt{Sections/}\allowbreak{}\texttt{RecoveryControlSectionV1.}\allowbreak{}\texttt{cs}
\item
  Create:
  \texttt{tests/}\allowbreak{}\texttt{FirstGearBank.}\allowbreak{}\texttt{Core.}\allowbreak{}\texttt{Tests/}\allowbreak{}\texttt{Persistence/}\allowbreak{}\texttt{BankEnvelopeCodecTests.}\allowbreak{}\texttt{cs}
\item
  Create:
  \texttt{tests/}\allowbreak{}\texttt{FirstGearBank.}\allowbreak{}\texttt{Core.}\allowbreak{}\texttt{Tests/}\allowbreak{}\texttt{Persistence/}\allowbreak{}\texttt{BankStateRoundTripTests.}\allowbreak{}\texttt{cs}
\item
  Create:
  \texttt{tests/}\allowbreak{}\texttt{FirstGearBank.}\allowbreak{}\texttt{Core.}\allowbreak{}\texttt{Tests/}\allowbreak{}\texttt{Persistence/}\allowbreak{}\texttt{SectionIsolationTests.}\allowbreak{}\texttt{cs}
\item
  Create:
  \texttt{tests/}\allowbreak{}\texttt{FirstGearBank.}\allowbreak{}\texttt{Core.}\allowbreak{}\texttt{Tests/}\allowbreak{}\texttt{Persistence/}\allowbreak{}\texttt{SectionMigrationRegistryTests.}\allowbreak{}\texttt{cs}
\item
  Create:
  \texttt{tests/}\allowbreak{}\texttt{FirstGearBank.}\allowbreak{}\texttt{Core.}\allowbreak{}\texttt{Tests/}\allowbreak{}\texttt{Persistence/}\allowbreak{}\texttt{InstallMarkerTests.}\allowbreak{}\texttt{cs}
\item
  Create:
  \texttt{tests/}\allowbreak{}\texttt{FirstGearBank.}\allowbreak{}\texttt{Core.}\allowbreak{}\texttt{Tests/}\allowbreak{}\texttt{Persistence/}\allowbreak{}\texttt{BankPersistenceServiceTests.}\allowbreak{}\texttt{cs}
\item
  Create:
  \texttt{tests/}\allowbreak{}\texttt{FirstGearBank.}\allowbreak{}\texttt{Core.}\allowbreak{}\texttt{Tests/}\allowbreak{}\texttt{TestSupport/}\allowbreak{}\texttt{PersistenceTestData.}\allowbreak{}\texttt{cs}
\item
  Create:
  \texttt{tests/}\allowbreak{}\texttt{FirstGearBank.}\allowbreak{}\texttt{Core.}\allowbreak{}\texttt{Tests/}\allowbreak{}\texttt{TestSupport/}\allowbreak{}\texttt{RecordingTestMigration.}\allowbreak{}\texttt{cs}
\item
  Create:
  \texttt{tests/}\allowbreak{}\texttt{FirstGearBank.}\allowbreak{}\texttt{Core.}\allowbreak{}\texttt{Tests/}\allowbreak{}\texttt{TestSupport/}\allowbreak{}\texttt{DeterministicEntropySource.}\allowbreak{}\texttt{cs}
\item
  Create:
  \texttt{tests/}\allowbreak{}\texttt{FirstGearBank.}\allowbreak{}\texttt{Core.}\allowbreak{}\texttt{Tests/}\allowbreak{}\texttt{Fixtures/}\allowbreak{}\texttt{Persistence/}\allowbreak{}\texttt{schema-}\allowbreak{}\texttt{1-}\allowbreak{}\texttt{empty.}\allowbreak{}\texttt{hex}
\item
  Create:
  \texttt{tests/}\allowbreak{}\texttt{FirstGearBank.}\allowbreak{}\texttt{Core.}\allowbreak{}\texttt{Tests/}\allowbreak{}\texttt{Fixtures/}\allowbreak{}\texttt{Persistence/}\allowbreak{}\texttt{schema-}\allowbreak{}\texttt{1-}\allowbreak{}\texttt{one-}\allowbreak{}\texttt{deposit.}\allowbreak{}\texttt{hex}
\item
  Modify:
  \texttt{FirstGearBank.}\allowbreak{}\texttt{Core/}\allowbreak{}\texttt{State/}\allowbreak{}\texttt{BankState.}\allowbreak{}\texttt{cs:}\allowbreak{}\texttt{1}
\item
  Modify:
  \texttt{FirstGearBank.}\allowbreak{}\texttt{Core/}\allowbreak{}\texttt{State/}\allowbreak{}\texttt{BankStateSnapshot.}\allowbreak{}\texttt{cs:}\allowbreak{}\texttt{1}
\item
  Modify:
  \texttt{FirstGearBank.}\allowbreak{}\texttt{Core/}\allowbreak{}\texttt{State/}\allowbreak{}\texttt{BankStateRepository.}\allowbreak{}\texttt{cs:}\allowbreak{}\texttt{1}
\item
  Modify:
  \texttt{FirstGearBank.}\allowbreak{}\texttt{Core/}\allowbreak{}\texttt{State/}\allowbreak{}\texttt{BankStateValidator.}\allowbreak{}\texttt{cs:}\allowbreak{}\texttt{1}
\item
  Modify:
  \texttt{FirstGearBank.}\allowbreak{}\texttt{Core/}\allowbreak{}\texttt{Coordination/}\allowbreak{}\texttt{BankingCoordinator.}\allowbreak{}\texttt{cs:}\allowbreak{}\texttt{1}
\item
  Modify:
  \texttt{tests/}\allowbreak{}\texttt{FirstGearBank.}\allowbreak{}\texttt{Core.}\allowbreak{}\texttt{Tests/}\allowbreak{}\texttt{Recovery/}\allowbreak{}\texttt{QuarantineTransitionTests.}\allowbreak{}\texttt{cs:}\allowbreak{}\texttt{1}
\item
  Modify:
  \texttt{docs/}\allowbreak{}\texttt{architecture/}\allowbreak{}\texttt{0004-}\allowbreak{}\texttt{persistence-}\allowbreak{}\texttt{sections-}\allowbreak{}\texttt{and-}\allowbreak{}\texttt{revisions.}\allowbreak{}\texttt{md:}\allowbreak{}\texttt{1}
\item
  Modify:
  \texttt{docs/}\allowbreak{}\texttt{architecture/}\allowbreak{}\texttt{0004-}\allowbreak{}\texttt{persistence-}\allowbreak{}\texttt{sections-}\allowbreak{}\texttt{and-}\allowbreak{}\texttt{revisions.}\allowbreak{}\texttt{tex:}\allowbreak{}\texttt{1}
  through regeneration
\item
  Generate but do not commit:
  \texttt{docs/}\allowbreak{}\texttt{architecture/}\allowbreak{}\texttt{0004-}\allowbreak{}\texttt{persistence-}\allowbreak{}\texttt{sections-}\allowbreak{}\texttt{and-}\allowbreak{}\texttt{revisions.}\allowbreak{}\texttt{pdf}
\end{itemize}

\textbf{Interfaces:}

\begin{itemize}
\tightlist
\item
  Consumes: immutable state, replay, request control, configuration,
  settlement evidence, and revision contracts.
\item
  Produces:
  \texttt{BankPersistenceService.}\allowbreak{}\texttt{Load}/\passthrough{\lstinline!Save!},
  canonical section codecs, migration registry, raw quarantine evidence,
  and the \passthrough{\lstinline!IBankWorldStateStore!} port
  implemented by Task 12.
\end{itemize}

\begin{itemize}
\tightlist
\item[$\square$]
  \textbf{Step 1: Write failing deterministic frame and corruption
  tests}
\end{itemize}

\begin{lstlisting}
[Fact]
public void Same_sections_encode_to_the_same_bytes()
{
    var envelope = PersistenceTestData.FoundationEnvelope;

    var first = BankEnvelopeCodec.Encode(envelope).Value;
    var second = BankEnvelopeCodec.Encode(envelope).Value;

    Assert.True(first.Span.SequenceEqual(second.Span));
}

[Fact]
public void Corrupt_registry_payload_does_not_invalidate_verified_journal_payload()
{
    var bytes = BankEnvelopeCodec.Encode(PersistenceTestData.FoundationEnvelope).Value.ToArray();
    PersistenceTestData.FlipPayloadByte(bytes, BankSectionKind.RecipientRegistry);

    var result = BankEnvelopeCodec.Decode(bytes);

    Assert.True(result.VerifiedSections.ContainsKey(BankSectionKind.Journal));
    Assert.Equal(
        QuarantineDomain.RecipientRegistry,
        Assert.Single(result.QuarantinedSections).Domain);
}

[Fact]
public void Unknown_raw_v1_kind_quarantines_the_whole_envelope_without_reserializing_it()
{
    var stored = PersistenceTestData.EnvelopeWithUnknownRawKind(0x7fff);

    var result = BankEnvelopeCodec.Decode(stored);

    Assert.True(result.IsWholeEnvelopeQuarantined);
    Assert.True(stored.AsSpan().SequenceEqual(result.RawEnvelopeEvidence.Span));
    Assert.Empty(result.VerifiedSections);
    Assert.False(result.CanReserialize);
}
\end{lstlisting}

Add magic, container version, global revision, descriptor checksum, sort
order, duplicate kind, overlap, gap, offset/length overflow, truncation,
payload checksum, trailing bytes, section count, complete-size, and
per-section-size tests. Run the unknown-kind case once with a
checksum-valid payload and once with a checksum-invalid payload; both
retain the byte-identical complete envelope and expose no partially
decoded state to a reserializer.

\begin{itemize}
\tightlist
\item[$\square$]
  \textbf{Step 2: Run codec tests and observe missing-type failures}
\end{itemize}

\begin{lstlisting}
dotnet test tests/FirstGearBank.Core.Tests/FirstGearBank.Core.Tests.csproj -c Release
\end{lstlisting}

Expected: FAIL because the envelope codec and section catalog do not
exist.

\begin{itemize}
\tightlist
\item[$\square$]
  \textbf{Step 3: Implement the binary container and per-section
  integrity rules}
\end{itemize}

Use this exact big-endian container layout:

\begin{lstlisting}
4 bytes   ASCII magic "FGBK"
2 bytes   unsigned container version = 1
8 bytes   signed global revision
2 bytes   unsigned section count
4 bytes   unsigned descriptor-table byte length
N bytes   descriptor table, entries sorted by numeric section kind
32 bytes  SHA-256 of descriptor table
M bytes   contiguous section payloads in descriptor order

descriptor entry:
2 bytes   unsigned raw section-kind code
2 bytes   unsigned section schema version
8 bytes   signed subsystem revision
8 bytes   unsigned payload offset relative to first payload byte
4 bytes   unsigned payload length
32 bytes  SHA-256 of exact payload bytes
\end{lstlisting}

Require offsets to be contiguous and nonoverlapping, kinds
unique/increasing, revisions nonnegative, and no trailing bytes. An
invalid magic/\allowbreak{}header/\allowbreak{}descriptor
checksum/layout causes whole-envelope quarantine. If the descriptor
table is trusted, an invalid payload checksum quarantines only the
section\textquotesingle s code-owned domain and preserves its exact raw
bytes.

Use these stable kinds:

\begin{lstlisting}
public enum BankSectionKind : ushort
{
    WorldAuthority = 1,
    Journal = 2,
    ProjectionCache = 3,
    RequestControl = 4,
    Configuration = 5,
    RecipientRegistry = 6,
    DeliveryControl = 7,
    PhysicalSettlement = 8,
    BranchControl = 9,
    DiagnosticsCache = 10,
    FinancialClock = 11,
    FinancialMarket = 12,
    RecoveryControl = 13
}
\end{lstlisting}

The financial-clock and market slots are reserved now and become
required only through explicit migrations in the financial-core plan.
\passthrough{\lstinline!RecoveryControl!} is required in the foundation
schema and stores authoritative quarantine/audit and resolved-evidence
archival state; corruption in it blocks all mutation because recovery
status can no longer be trusted.
\passthrough{\lstinline!BankSectionDescriptor!} retains the raw
\passthrough{\lstinline!ushort!} kind code before attempting enum
conversion. Container version 1 has no trustworthy criticality bit, so a
genuinely unknown kind causes whole-envelope read-only quarantine and
exact raw envelope retention; it is never silently treated as an
optional section. A later container may add an authenticated criticality
rule only through an explicit migration. A known kind with an
unsupported version follows the dependency matrix below and is never
discarded or resaved as if understood.

\begin{lstlisting}
public sealed record QuarantinedSection(
    ushort RawKind,
    ushort SchemaVersion,
    long SubsystemRevision,
    ImmutableArray<byte> DeclaredSha256,
    ImmutableArray<byte> ExactPayload,
    QuarantineDomain Domain,
    SectionQuarantineReason Reason);
\end{lstlisting}

The descriptor offset may be recomputed when healthy sections change,
but every other field above and the corrupt payload bytes pass through
unchanged. Re-encoding therefore preserves the failed declared hash
rather than blessing corrupt bytes with a new checksum.

Use this foundation dependency matrix; later schema plans may only make
it more specific with evidence:

{\def\LTcaptype{none} % do not increment counter
\begin{longtable}[]{@{}
  >{\raggedright\arraybackslash}p{(\linewidth - 4\tabcolsep) * \real{0.2900}}
  >{\raggedright\arraybackslash}p{(\linewidth - 4\tabcolsep) * \real{0.3400}}
  >{\raggedright\arraybackslash}p{(\linewidth - 4\tabcolsep) * \real{0.3700}}@{}}
\toprule\noalign{}
\begin{minipage}[b]{\linewidth}\raggedright
Section problem
\end{minipage} & \begin{minipage}[b]{\linewidth}\raggedright
Services blocked
\end{minipage} & \begin{minipage}[b]{\linewidth}\raggedright
Services still permitted
\end{minipage} \\
\midrule\noalign{}
\endhead
\bottomrule\noalign{}
\endlastfoot
\passthrough{\lstinline!WorldAuthority!},
\passthrough{\lstinline!RecoveryControl!}, invalid container, descriptor
table, or unknown kind & Entire Bank; preserve/export only &
Permission-gated diagnostics/export \\
\passthrough{\lstinline!Journal!} & All finance and automatic finance
mutations & Diagnostics/export; branch repair only when independent \\
\passthrough{\lstinline!ProjectionCache!} & None after verified replay
repair & All services allowed by other sections \\
\passthrough{\lstinline!RequestControl!} & Every authenticated and
automatic mutation & Read-only account/history and diagnostics \\
\passthrough{\lstinline!Configuration!} & Config activation/edit; retain
last verified active state & Services using independently verified
active market state \\
\passthrough{\lstinline!RecipientRegistry!} & Name resolution,
transfers, name-based corrections & Own-account services and unrelated
automatic work \\
\passthrough{\lstinline!DeliveryControl!} & Transfers/corrections that
require durable notices & Own deposits, withdrawals, CDs, statements,
diagnostics \\
\passthrough{\lstinline!PhysicalSettlement!} & All player banking and
finance mutation until affected players can be proved & Read-only views
and diagnostics \\
\passthrough{\lstinline!BranchControl!} & Banker sessions and branch
operations & Permission-gated nonbranch diagnostics/export \\
\passthrough{\lstinline!DiagnosticsCache!} & None; rebuild from verified
authority & All services allowed by other sections \\
\passthrough{\lstinline!FinancialClock!} or
\passthrough{\lstinline!FinancialMarket!} once required & Finance
mutation and quotes & Read-only views, branch repair, diagnostics \\
\end{longtable}
}

\begin{itemize}
\tightlist
\item[$\square$]
  \textbf{Step 4: Set and test structural limits}
\end{itemize}

\begin{lstlisting}
public static class BankPersistenceLimits
{
    public const int MaximumContainerBytes = 256 * 1024 * 1024;
    public const ushort MaximumSections = 64;
    public const int MaximumSectionBytes = 192 * 1024 * 1024;
    public const int MaximumJournalRecords = 5_000_000;
    public const int MaximumPostingsPerRecord = 16;
    public const int MaximumPartiesPerRecord = 8;
    public const int MaximumOperationIdsPerSettlement = 16;
    public const int MaximumInventorySlotDeltasPerSettlement = PhysicalSettlementLimits.MaximumInventorySlotDeltas;
    public const int MaximumCanonicalStackBytes = PhysicalSettlementLimits.MaximumCanonicalStackBytes;
    public const int MaximumRecoveryPlanBytes = PhysicalSettlementLimits.MaximumRecoveryPlanBytes;
    public const int MaximumArchivedRecoveryEvidenceBytes = RecoveryLimits.MaximumArchivedEvidenceBytes;
    public const int MaximumUtf8NameBytes = 256;
    public const int MaximumUtf8ReasonBytes = 4_096;
}
\end{lstlisting}

Reject sizes before allocating buffers or enumerating contents. Limits
are structural safety invariants, not user config.

\begin{itemize}
\tightlist
\item[$\square$]
  \textbf{Step 5: Write failing canonical section and replay-repair
  tests}
\end{itemize}

\begin{lstlisting}
[Fact]
public void Bad_projection_cache_is_rebuilt_from_verified_journal()
{
    var stored = PersistenceTestData.StateWithValidJournalAndWrongProjectionCache;

    var outcome = BankStateLoader.Load(stored, PersistenceTestData.ExpectedWorldBinding);

    Assert.False(outcome.Quarantine.IsBlocked(QuarantineDomain.Finance));
    Assert.True(outcome.RequiresResave);
    var finance = outcome.State.Finance.TryGetVerified().Value;
    var replayed = JournalProjector.Foundation.Replay(
        finance.Journal,
        ProjectionCheckpoint.Empty).Value;
    var match = ProjectionValidator.MatchesReplay(finance.Projection, replayed);

    Assert.True(match.IsSuccess);
    Assert.True(match.Value);
}

[Fact]
public void Corrupt_journal_preserves_raw_bytes_and_never_uses_cache_as_authority()
{
    var stored = PersistenceTestData.EnvelopeWithCorruptJournalAndValidCache;

    var outcome = BankStateLoader.Load(stored, PersistenceTestData.ExpectedWorldBinding);

    Assert.True(outcome.Quarantine.IsBlocked(QuarantineDomain.Finance));
    Assert.True(stored.AsSpan().SequenceEqual(outcome.RawEnvelopeEvidence.Span));
    Assert.False(outcome.CanMutateFinance);
    Assert.True(outcome.State.Finance.IsUnavailable);
    Assert.False(outcome.State.Finance.TryGetVerified().IsSuccess);
}

[Fact]
public void Corrupt_journal_does_not_require_a_fake_empty_finance_state_for_an_independent_config_commit()
{
    var rig = PersistenceRecoveryRig.Load(PersistenceTestData.EnvelopeWithCorruptJournal);
    var originalEvidence = rig.RawSection(BankSectionKind.Journal).ToArray();

    var response = rig.Coordinator.Execute(
        RequestTestData.Context(sequence: 1),
        new ConfigurationOnlyMutation(ConfigurationTestData.ValidDisplayPrecisionFour));
    rig.StageAndReload();

    Assert.True(response.IsSuccess);
    Assert.True(rig.State.Finance.IsUnavailable);
    Assert.Equal(4, rig.State.Configuration.TryGetVerified().Value.DisplayPrecision);
    Assert.True(originalEvidence.AsSpan().SequenceEqual(rig.RawSection(BankSectionKind.Journal).Span));
}
\end{lstlisting}

Add roundtrips for every foundation journal payload and invocation
variant, request high-water/\allowbreak{}cache/\allowbreak{}permanent
request key, permanent automatic
key/\allowbreak{}digest/\allowbreak{}result, configuration candidate,
empty current registry epoch, empty delivery/branch state, all
settlement phases/complete recovery plans/capsules, recovery-control
evidence/audit, every verified/unavailable domain-slot variant, revision
vector, and raw diagnostics. A loader may never satisfy a missing
authoritative slot with an empty object. Add the decisive pass-through
scenario: load a corrupt registry alongside a valid journal, commit a
healthy permitted financial revision, stage and reload the envelope,
then prove the journal advanced while the registry remains quarantined
with byte-identical corrupt payload, declared hash, schema version, raw
kind, and subsystem revision.

\begin{lstlisting}
dotnet test tests/FirstGearBank.Core.Tests/FirstGearBank.Core.Tests.csproj -c Release
\end{lstlisting}

Expected: FAIL because section codecs, replay repair, and complete state
reconciliation are not implemented.

\begin{itemize}
\tightlist
\item[$\square$]
  \textbf{Step 6: Implement explicit DTO codecs and state
  reconciliation}
\end{itemize}

Each section codec maps domain state to a versioned DTO; domain
interfaces are never passed directly to
\passthrough{\lstinline!JsonSerializer!}. Canonical JSON uses
declaration-ordered properties, UTF-8 without BOM, no indentation,
invariant integer numbers, \passthrough{\lstinline!G29!} strings for
decimals, \passthrough{\lstinline!R!} strings for finite doubles,
lowercase hex IDs/hashes, and sorted arrays instead of dictionaries.
Deserialize with strict duplicate-property, enum, size, and
trailing-token checks.

\passthrough{\lstinline!WorldAuthoritySectionV1!} contains the
\passthrough{\lstinline!WorldInstanceId!}, the SHA-256 world binding,
and the 32-byte recovery authentication key. Its canonical binding is
\texttt{SHA-}\allowbreak{}\texttt{256(}\allowbreak{}\texttt{"FGB/}\allowbreak{}\texttt{WORLD-}\allowbreak{}\texttt{BINDING/}\allowbreak{}\texttt{1\textbackslash{}}\allowbreak{}\texttt{0"\ }\allowbreak{}\texttt{\textbar{}}\allowbreak{}\texttt{\textbar{}}\allowbreak{}\texttt{\ }\allowbreak{}\texttt{UInt32BigEndian(}\allowbreak{}\texttt{identifierLength)}\allowbreak{}\texttt{\ }\allowbreak{}\texttt{\textbar{}}\allowbreak{}\texttt{\textbar{}}\allowbreak{}\texttt{\ }\allowbreak{}\texttt{UTF8(}\allowbreak{}\texttt{SavegameIdentifier)}\allowbreak{}\texttt{)}\allowbreak{}.
Corruption or mismatch in this section quarantines the complete envelope
because no other section may be imported into a different world.

This task extends \passthrough{\lstinline!BankStateSnapshot!} and its
repository-owned immutable publication unit with a bounded
\texttt{ImmutableArray\textless{}}\allowbreak{}\texttt{QuarantinedSection\textgreater{}}\allowbreak{}
opaque-section set. Coordinator candidates carry the same set by
reference and cannot read, remove, or rewrite it; only the loader and
the typed recovery path may create or consume an entry. Snapshot capture
therefore cannot lose pass-through evidence while a healthy domain
publishes concurrently.
\texttt{BankingCoordinator.}\allowbreak{}\texttt{ExecuteRecovery}
receives an internal repository capability that can select only the one
opaque entry whose domain and \passthrough{\lstinline!EvidenceDigest!}
match its \passthrough{\lstinline!QuarantineTransition!}; ordinary
\passthrough{\lstinline!BankMutationReadView!} never receives that
capability. A successful clear moves those exact bytes into
\passthrough{\lstinline!RecoveryControlState!} and removes the opaque
entry in the same single snapshot publication. Add a test that changes
the opaque evidence between authorization and commit and prove the stale
recovery rejects without replacing the unavailable slot.

\passthrough{\lstinline!BankStateLoader!} validates the world binding
and authority key, decodes/migrates each section, verifies the journal,
replays from a verified checkpoint or sequence one, compares the cache,
repairs only derivable caches, clears active operation scopes for
restart, retains permanent request and automatic-operation barriers, and
assembles the quarantine state before returning any service access.
Every authoritative domain becomes either
\texttt{BankDomainSlot.}\allowbreak{}\texttt{Verified(}\allowbreak{}\texttt{value)}\allowbreak{}
or
\texttt{BankDomainSlot.}\allowbreak{}\texttt{Unavailable(}\allowbreak{}\texttt{domain,}\allowbreak{}\texttt{\ }\allowbreak{}\texttt{evidenceDigest)}\allowbreak{};
the loader never substitutes an empty value for unreadable authority.
\texttt{RecoveryControlSectionV1} is required, structurally bounded, and
cross-checks every unavailable
slot/\allowbreak{}quarantine/\allowbreak{}archive entry.
\passthrough{\lstinline!RequestControlSectionV1!} owns both registries
so no finance path can load without the exact deduplication state
required for its invocation keys. Startup cross-checks every journal
request/automatic invocation against the corresponding permanent entry,
digest, operation-ID set, and revision; a missing or contradictory
barrier is authoritative corruption, not a cache repair.
\passthrough{\lstinline!BankStateSerializer!} receives both healthy
domain state and retained \passthrough{\lstinline!QuarantinedSection!}
values. It emits quarantined known sections as opaque pass-through
payloads with their original declared hashes and metadata while healthy
section descriptors and the table checksum advance normally. It
serializes no payload for an unavailable slot except that
slot\textquotesingle s exact opaque pass-through section. Once an
authorized repair clears a slot, the old opaque bytes must already exist
byte-identically in
\texttt{RecoveryControlState.}\allowbreak{}\texttt{EvidenceArchive};
only then may the repaired section replace the pass-through entry.
Whole-envelope quarantine never enters this reserializer.

Fresh-state creation requests its world ID and 32-byte recovery key from
the Task 5 \passthrough{\lstinline!IEntropySource!}. Task 12 wires
\texttt{CryptographicEntropySource}; golden fixtures inject only the
listed deterministic test implementation.

\begin{itemize}
\tightlist
\item[$\square$]
  \textbf{Step 7: Write failing migration and unknown-version tests}
\end{itemize}

\begin{lstlisting}
[Fact]
public void Migration_validates_old_and_new_sections_before_replacement()
{
    var migration = new RecordingTestMigration(fromVersion: 1, toVersion: 2);
    var registry = SectionMigrationRegistry.Create([migration]);

    var result = registry.Migrate(
        BankSectionKind.DiagnosticsCache,
        sourceVersion: 1,
        PersistenceTestData.ValidDiagnosticsV1);

    Assert.True(result.IsSuccess);
    Assert.Equal(["validate-old", "transform", "validate-new"], migration.Calls);
}
\end{lstlisting}

Test missing link, duplicate edge, cycle, nondeterministic output
detection, invalid old/new data, unknown newer container, unknown newer
critical section, and preservation of its raw bytes.

\begin{lstlisting}
dotnet test tests/FirstGearBank.Core.Tests/FirstGearBank.Core.Tests.csproj -c Release
\end{lstlisting}

Expected: FAIL because deterministic migration-chain validation is not
implemented.

\begin{itemize}
\tightlist
\item[$\square$]
  \textbf{Step 8: Implement deterministic per-section migration chains}
\end{itemize}

\begin{lstlisting}
public interface ISectionMigration
{
    BankSectionKind SectionKind { get; }
    ushort FromVersion { get; }
    ushort ToVersion { get; }

    BankResult<ReadOnlyMemory<byte>> Migrate(ReadOnlyMemory<byte> verifiedSource);
}
\end{lstlisting}

The registry accepts only \passthrough{\lstinline!n -> n+1!} edges, one
edge per \passthrough{\lstinline!(kind, from)!}, acyclic complete paths,
and byte-identical output when the same migration is invoked twice in
tests. Source validation happens before transformation and target
validation before candidate installation. No migration writes live state
in place.

\begin{itemize}
\tightlist
\item[$\square$]
  \textbf{Step 9: Write failing install-marker and world-save-staging
  tests}
\end{itemize}

Use two namespaced world keys so first installation can be distinguished
from deleted authority:

\begin{lstlisting}
firstgearbank:installed
firstgearbank:world-state
\end{lstlisting}

Test this matrix:

\begin{lstlisting}
marker absent + state absent  -> create fresh unstaged state; block mutations until initial snapshot is staged
marker absent + state valid   -> load state and repair marker on the next world-save staging callback
marker present + state absent -> whole-Bank quarantine; do not create defaults
marker/state binding mismatch -> whole-Bank quarantine
marker present + state valid  -> normal load
\end{lstlisting}

Also test store exception, serialization exception, exact-revision
unstaged-flag clearing, a newer publish during staging, the always-on
engine-restage flag, and that a first-save marker-stage failure after a
successful envelope stage reloads the verified state and repairs only
the marker.

\begin{lstlisting}
dotnet test tests/FirstGearBank.Core.Tests/FirstGearBank.Core.Tests.csproj -c Release
\end{lstlisting}

Expected: FAIL because the install marker and world-save staging
protocol are not implemented.

\begin{itemize}
\tightlist
\item[$\square$]
  \textbf{Step 10: Implement the persistence service and two-key
  protocol}
\end{itemize}

\begin{lstlisting}
public interface IBankWorldStateStore
{
    ReadOnlyMemory<byte>? ReadInstallMarker();
    ReadOnlyMemory<byte>? ReadEnvelope();
    void StageInstallMarker(ReadOnlyMemory<byte> bytes);
    void StageEnvelope(ReadOnlyMemory<byte> bytes);
}

public sealed class BankPersistenceService
{
    public BankLoadOutcome Load(string savegameIdentifier);
    public BankResult<BankRevision> StageForWorldSave(BankStateRepository repository);
}
\end{lstlisting}

On save, serialize an immutable snapshot first, then stage the envelope
before the install marker. On first installation, the marker is the
final commit flag; a crash after the envelope stage is therefore the
recoverable
\texttt{marker\ }\allowbreak{}\texttt{absent\ }\allowbreak{}\texttt{+\ }\allowbreak{}\texttt{state\ }\allowbreak{}\texttt{valid}
case instead of the unrecoverable
\texttt{marker\ }\allowbreak{}\texttt{present\ }\allowbreak{}\texttt{+\ }\allowbreak{}\texttt{state\ }\allowbreak{}\texttt{absent}
case. The marker contains container version,
\passthrough{\lstinline!WorldInstanceId!}, and world-binding hash but
never the recovery HMAC key. If a crash leaves marker-only state, the
next load quarantines safely. If state exists without a marker, its
verified embedded world authority permits marker repair.
\passthrough{\lstinline!HasUnstagedChanges!} clears only after both
staging calls return for the exact captured revision;
\texttt{RequiresEngineSaveRestage} does not. This result means accepted
by the engine\textquotesingle s pre-save world-data staging API, not
written to disk. The service calls
\passthrough{\lstinline!MarkStagingSucceeded!}/\passthrough{\lstinline!MarkStagingFailed!};
the host restages on every later \passthrough{\lstinline!GameWorldSave!}
because 1.22.7 exposes no post-save acknowledgment.

\begin{itemize}
\tightlist
\item[$\square$]
  \textbf{Step 11: Add stable golden fixtures}
\end{itemize}

Encode one empty state and one single-deposit state using a
deterministic test entropy source, write uppercase byte-pair hex with
one line per 32 bytes, then make tests decode and compare against
freshly encoded bytes. Fixture regeneration must be an explicit test
utility argument and ordinary test runs must fail on a byte change.

\begin{lstlisting}
dotnet test tests/FirstGearBank.Core.Tests/FirstGearBank.Core.Tests.csproj -c Release
\end{lstlisting}

Expected: PASS for roundtrip, deterministic bytes, migration, cache
repair, isolation, limits, and staging behavior.

\begin{itemize}
\tightlist
\item[$\square$]
  \textbf{Step 12: Commit the persistence schema and decision record}
\end{itemize}

Regenerate the decision record\textquotesingle s same-basename
\passthrough{\lstinline!.tex!} and \passthrough{\lstinline!.pdf!} from
Markdown and visually verify the rendered PDF before staging only the
Markdown and LaTeX sources.

\begin{lstlisting}
git diff --check
git add FirstGearBank.Core/Persistence FirstGearBank.Core/State `
  FirstGearBank.Core/Coordination/BankingCoordinator.cs `
  tests/FirstGearBank.Core.Tests/Persistence `
  tests/FirstGearBank.Core.Tests/Recovery/QuarantineTransitionTests.cs `
  tests/FirstGearBank.Core.Tests/TestSupport/PersistenceTestData.cs `
  tests/FirstGearBank.Core.Tests/TestSupport/RecordingTestMigration.cs `
  tests/FirstGearBank.Core.Tests/TestSupport/DeterministicEntropySource.cs `
  tests/FirstGearBank.Core.Tests/Fixtures/Persistence `
  docs/architecture/0004-persistence-sections-and-revisions.md `
  docs/architecture/0004-persistence-sections-and-revisions.tex
git commit -m "feat: add sectioned bank state persistence"
\end{lstlisting}

\subsubsection{Task 12: Wire the Thin Vintage Story Server Lifecycle,
Configuration, Capsule, and Log
Adapters}\label{task-12-wire-the-thin-vintage-story-server-lifecycle-configuration-capsule-and-log-adapters}

\textbf{Files:}

\begin{itemize}
\tightlist
\item
  Create:
  \texttt{First\ }\allowbreak{}\texttt{Gear\ }\allowbreak{}\texttt{Bank/}\allowbreak{}\texttt{FirstGearBankModSystem.}\allowbreak{}\texttt{cs}
\item
  Create:
  \texttt{First\ }\allowbreak{}\texttt{Gear\ }\allowbreak{}\texttt{Bank/}\allowbreak{}\texttt{Runtime/}\allowbreak{}\texttt{ServerHostState.}\allowbreak{}\texttt{cs}
\item
  Create:
  \texttt{First\ }\allowbreak{}\texttt{Gear\ }\allowbreak{}\texttt{Bank/}\allowbreak{}\texttt{Runtime/}\allowbreak{}\texttt{FirstGearBankServerHost.}\allowbreak{}\texttt{cs}
\item
  Create:
  \texttt{First\ }\allowbreak{}\texttt{Gear\ }\allowbreak{}\texttt{Bank/}\allowbreak{}\texttt{Persistence/}\allowbreak{}\texttt{VintageStoryWorldStateStore.}\allowbreak{}\texttt{cs}
\item
  Create:
  \texttt{First\ }\allowbreak{}\texttt{Gear\ }\allowbreak{}\texttt{Bank/}\allowbreak{}\texttt{Configuration/}\allowbreak{}\texttt{VintageStoryConfigurationStore.}\allowbreak{}\texttt{cs}
\item
  Create:
  \texttt{First\ }\allowbreak{}\texttt{Gear\ }\allowbreak{}\texttt{Bank/}\allowbreak{}\texttt{Settlement/}\allowbreak{}\texttt{VintageStoryPlayerCapsuleStore.}\allowbreak{}\texttt{cs}
\item
  Create:
  \texttt{First\ }\allowbreak{}\texttt{Gear\ }\allowbreak{}\texttt{Bank/}\allowbreak{}\texttt{Diagnostics/}\allowbreak{}\texttt{FgbFileLogger.}\allowbreak{}\texttt{cs}
\item
  Create:
  \texttt{First\ }\allowbreak{}\texttt{Gear\ }\allowbreak{}\texttt{Bank/}\allowbreak{}\texttt{Diagnostics/}\allowbreak{}\texttt{VintageStoryBankLogger.}\allowbreak{}\texttt{cs}
\item
  Delete:
  \texttt{First\ }\allowbreak{}\texttt{Gear\ }\allowbreak{}\texttt{Bank/}\allowbreak{}\texttt{First\ }\allowbreak{}\texttt{Gear\ }\allowbreak{}\texttt{BankModSystem.}\allowbreak{}\texttt{cs}
\item
  Create:
  \texttt{tests/}\allowbreak{}\texttt{FirstGearBank.}\allowbreak{}\texttt{Adapter.}\allowbreak{}\texttt{Tests/}\allowbreak{}\texttt{Runtime/}\allowbreak{}\texttt{FirstGearBankServerHostTests.}\allowbreak{}\texttt{cs}
\item
  Create:
  \texttt{tests/}\allowbreak{}\texttt{FirstGearBank.}\allowbreak{}\texttt{Adapter.}\allowbreak{}\texttt{Tests/}\allowbreak{}\texttt{Configuration/}\allowbreak{}\texttt{VintageStoryConfigurationStoreTests.}\allowbreak{}\texttt{cs}
\item
  Create:
  \texttt{tests/}\allowbreak{}\texttt{FirstGearBank.}\allowbreak{}\texttt{Adapter.}\allowbreak{}\texttt{Tests/}\allowbreak{}\texttt{Diagnostics/}\allowbreak{}\texttt{FgbFileLoggerTests.}\allowbreak{}\texttt{cs}
\item
  Create:
  \texttt{tests/}\allowbreak{}\texttt{FirstGearBank.}\allowbreak{}\texttt{Adapter.}\allowbreak{}\texttt{Tests/}\allowbreak{}\texttt{Persistence/}\allowbreak{}\texttt{WorldStateStoreContractTests.}\allowbreak{}\texttt{cs}
\item
  Create:
  \texttt{tests/}\allowbreak{}\texttt{FirstGearBank.}\allowbreak{}\texttt{Adapter.}\allowbreak{}\texttt{Tests/}\allowbreak{}\texttt{Settlement/}\allowbreak{}\texttt{PlayerCapsuleStoreContractTests.}\allowbreak{}\texttt{cs}
\item
  Create:
  \texttt{tests/}\allowbreak{}\texttt{FirstGearBank.}\allowbreak{}\texttt{Adapter.}\allowbreak{}\texttt{Tests/}\allowbreak{}\texttt{TestSupport/}\allowbreak{}\texttt{FakeBankWorldStateStore.}\allowbreak{}\texttt{cs}
\item
  Create:
  \texttt{tests/}\allowbreak{}\texttt{FirstGearBank.}\allowbreak{}\texttt{Adapter.}\allowbreak{}\texttt{Tests/}\allowbreak{}\texttt{TestSupport/}\allowbreak{}\texttt{ServerHostTestRig.}\allowbreak{}\texttt{cs}
\item
  Create:
  \texttt{tests/}\allowbreak{}\texttt{FirstGearBank.}\allowbreak{}\texttt{Adapter.}\allowbreak{}\texttt{Tests/}\allowbreak{}\texttt{TestSupport/}\allowbreak{}\texttt{TemporaryDirectory.}\allowbreak{}\texttt{cs}
\item
  Create:
  \texttt{tests/}\allowbreak{}\texttt{FirstGearBank.}\allowbreak{}\texttt{Adapter.}\allowbreak{}\texttt{Tests/}\allowbreak{}\texttt{TestSupport/}\allowbreak{}\texttt{TestClock.}\allowbreak{}\texttt{cs}
\item
  Create:
  \texttt{tests/}\allowbreak{}\texttt{FirstGearBank.}\allowbreak{}\texttt{Adapter.}\allowbreak{}\texttt{Tests/}\allowbreak{}\texttt{TestSupport/}\allowbreak{}\texttt{TestTemplate.}\allowbreak{}\texttt{cs}
\item
  Create:
  \texttt{tests/}\allowbreak{}\texttt{FirstGearBank.}\allowbreak{}\texttt{Adapter.}\allowbreak{}\texttt{Tests/}\allowbreak{}\texttt{TestSupport/}\allowbreak{}\texttt{LogTestReader.}\allowbreak{}\texttt{cs}
\item
  Modify:
  \texttt{First\ }\allowbreak{}\texttt{Gear\ }\allowbreak{}\texttt{Bank/}\allowbreak{}\texttt{First\ }\allowbreak{}\texttt{Gear\ }\allowbreak{}\texttt{Bank.}\allowbreak{}\texttt{csproj:}\allowbreak{}\texttt{1}
\item
  Modify: \passthrough{\lstinline!README.md:1!}
\item
  Modify:
  \texttt{docs/}\allowbreak{}\texttt{testing/}\allowbreak{}\texttt{foundation-}\allowbreak{}\texttt{engine-}\allowbreak{}\texttt{smoke.}\allowbreak{}\texttt{md:}\allowbreak{}\texttt{1}
\item
  Modify:
  \texttt{docs/}\allowbreak{}\texttt{testing/}\allowbreak{}\texttt{foundation-}\allowbreak{}\texttt{engine-}\allowbreak{}\texttt{smoke.}\allowbreak{}\texttt{tex:}\allowbreak{}\texttt{1}
  through regeneration
\item
  Generate but do not commit:
  \texttt{docs/}\allowbreak{}\texttt{testing/}\allowbreak{}\texttt{foundation-}\allowbreak{}\texttt{engine-}\allowbreak{}\texttt{smoke.}\allowbreak{}\texttt{pdf}
\end{itemize}

\textbf{Interfaces:}

\begin{itemize}
\tightlist
\item
  Consumes: \passthrough{\lstinline!IBankWorldStateStore!},
  \passthrough{\lstinline!BankPersistenceService!},
  \texttt{IPlayerRecoveryCapsuleStore}, configuration parser, and
  \passthrough{\lstinline!IBankLog!} from the core.
\item
  Produces: the only direct Vintage Story lifecycle/storage adapters and
  a foundation mod that loads under 1.22.7.
\end{itemize}

\begin{itemize}
\tightlist
\item[$\square$]
  \textbf{Step 1: Write failing host lifecycle tests behind fake stores}
\end{itemize}

\begin{lstlisting}
[Fact]
public void Fresh_world_stays_mutation_blocked_until_its_initial_state_is_staged()
{
    var store = new FakeBankWorldStateStore(marker: null, envelope: null);
    using var host = ServerHostTestRig.Create(store);

    host.OnSaveGameLoaded();

    Assert.Equal(ServerHostState.AwaitingInitialSave, host.State);
    Assert.False(host.CanAcceptMutations);

    host.OnGameWorldSave();

    Assert.Equal(ServerHostState.Ready, host.State);
    Assert.True(host.CanAcceptMutations);
    Assert.NotNull(store.Marker);
    Assert.NotNull(store.Envelope);
}

[Fact]
public void Repeated_load_callbacks_do_not_replace_live_state()
{
    var store = FakeBankWorldStateStore.WithValidEnvelope();
    using var host = ServerHostTestRig.Create(store);

    host.OnSaveGameLoaded();
    var first = host.Repository;
    host.OnSaveGameLoaded();

    Assert.Same(first, host.Repository);
    Assert.Equal(1, store.EnvelopeReadCount);
}
\end{lstlisting}

Add valid restore, marker-only quarantine, corrupt finance, isolated
registry corruption, staging exception, newer revision during staging,
repeated restaging without a new revision, create/load event order, save
callback before load, and disposed-host callback tests.

\begin{itemize}
\tightlist
\item[$\square$]
  \textbf{Step 2: Run host tests and observe missing-type failures}
\end{itemize}

\begin{lstlisting}
dotnet test tests/FirstGearBank.Adapter.Tests/FirstGearBank.Adapter.Tests.csproj -c Release
\end{lstlisting}

Expected: FAIL because the host does not exist.

\begin{itemize}
\tightlist
\item[$\square$]
  \textbf{Step 3: Implement a testable host with explicit lifecycle
  states}
\end{itemize}

\begin{lstlisting}
public sealed class FirstGearBankServerHost : IDisposable
{
    public ServerHostState State { get; }
    public bool CanAcceptMutations { get; }
    public BankStateRepository? Repository { get; }

    public void OnSaveGameCreated();
    public void OnSaveGameLoaded();
    public void OnGameWorldSave();
    public void Dispose();
}
\end{lstlisting}

The host is constructed with the savegame identifier, world-state store,
configuration store, entropy source, section codecs/migrations, and
logger. The first create/load callback performs exactly one load. Fresh
state is view-only until the first successful staging of marker and
envelope. Quarantine enables only domains allowed by the load outcome.
\passthrough{\lstinline!OnGameWorldSave!} captures one immutable
revision and calls \passthrough{\lstinline!StageForWorldSave!} even when
no newer revision exists. It logs "staged for engine world save" or the
exact staging failure, never "durably saved". The host cannot clear
\texttt{RequiresEngineSaveRestage} because the public event fires before
disk writing and exposes no completion callback.

The listed test-support types implement only project-owned ports, fixed
clocks/templates, and scoped temporary directories. They must not mock
or reimplement Vintage Story behavior.

\begin{itemize}
\tightlist
\item[$\square$]
  \textbf{Step 4: Verify host tests pass}
\end{itemize}

\begin{lstlisting}
dotnet test tests/FirstGearBank.Adapter.Tests/FirstGearBank.Adapter.Tests.csproj -c Release
\end{lstlisting}

\begin{itemize}
\tightlist
\item[$\square$]
  \textbf{Step 5: Write failing file-log and configuration-file tests}
\end{itemize}

\begin{lstlisting}
[Fact]
public void File_log_uses_only_the_four_exact_labels()
{
    using var directory = TemporaryDirectory.Create();
    using var log = new FgbFileLogger(directory.Path, maximumBytes: 2048, retainedFiles: 2, TestClock.Utc);

    log.Write(BankLogLevel.Info, "Host", "started");
    log.Write(BankLogLevel.Debug, "Host", "detail");
    log.Write(BankLogLevel.Warning, "Host", "degraded");
    log.Write(BankLogLevel.Critical, "Host", "blocked");

    Assert.Equal(
        ["INFO", "DEBG", "WARN", "CRIT"],
        LogTestReader.ReadLabels(
            Path.Combine(directory.Path, "logs", "firstgearbank.log")));
}

[Fact]
public void Valid_existing_configuration_is_read_without_rewrite()
{
    using var directory = TemporaryDirectory.Create();
    var path = directory.Write("firstgearbank.jsonc", "{\"TransferCooldownSeconds\":2.0}");
    var before = File.GetLastWriteTimeUtc(path);

    var result = new VintageStoryConfigurationStore(path, TestTemplate.Text).Load();

    Assert.Equal(2.0, result.Candidate.TransferCooldownSeconds);
    Assert.Equal(before, File.GetLastWriteTimeUtc(path));
}
\end{lstlisting}

Add missing-config canonical creation, invalid-config no rewrite,
explicit owned scalar save, one-time legacy migration backup, atomic
temp-file replacement, failed write preserving original, exact 2 MiB
rotation threshold, five-file retention, concurrent writes, exception
rendering, and no UID/balance fields in host lifecycle log calls.

\begin{lstlisting}
dotnet test tests/FirstGearBank.Adapter.Tests/FirstGearBank.Adapter.Tests.csproj -c Release
\end{lstlisting}

Expected: FAIL because bounded file logging and owned config-file
behavior are not implemented.

\begin{itemize}
\tightlist
\item[$\square$]
  \textbf{Step 6: Implement bounded logging and comment-preserving
  configuration storage}
\end{itemize}

Write logs beneath
\texttt{\textless{}}\allowbreak{}\texttt{runtime-}\allowbreak{}\texttt{data\textgreater{}}\allowbreak{}\texttt{/}\allowbreak{}\texttt{FirstGearBank/}\allowbreak{}\texttt{logs/}\allowbreak{}\texttt{firstgearbank.}\allowbreak{}\texttt{log},
rotate before a write would exceed 2 MiB, retain exactly five completed
files, and format each line as:

\begin{lstlisting}
2026-09-08T12:34:56.7890000Z INFO Host message
\end{lstlisting}

Escape embedded CR/LF so one event remains one record.
\passthrough{\lstinline!VintageStoryBankLogger!} writes to the file
first and maps \passthrough{\lstinline!INFO -> Notification!},
\passthrough{\lstinline!DEBG -> Debug!},
\passthrough{\lstinline!WARN -> Warning!}, and
\passthrough{\lstinline!CRIT -> Fatal!} on the native logger.

Place the live config at
\texttt{\textless{}}\allowbreak{}\texttt{ModConfig\textgreater{}}\allowbreak{}\texttt{/}\allowbreak{}\texttt{FirstGearBank/}\allowbreak{}\texttt{firstgearbank.}\allowbreak{}\texttt{jsonc}.
On absence, atomically copy the embedded canonical template. Ordinary
load never writes. An explicit owned edit applies
\passthrough{\lstinline!JsoncSyntaxDocument!} spans and replaces through
a same-directory temporary file. The one legacy migration writes a
\passthrough{\lstinline!.pre-fgb-migration!} backup once before atomic
replace.

\begin{itemize}
\tightlist
\item[$\square$]
  \textbf{Step 7: Verify log/config tests pass}
\end{itemize}

\begin{lstlisting}
dotnet test tests/FirstGearBank.Adapter.Tests/FirstGearBank.Adapter.Tests.csproj -c Release
\end{lstlisting}

\begin{itemize}
\tightlist
\item[$\square$]
  \textbf{Step 8: Write and run failing engine-storage adapter contract
  tests}
\end{itemize}

Use the internal delegate seams to prove key selection and buffer
ownership without pretending to emulate engine saves:

\begin{lstlisting}
[Fact]
public void World_store_clones_bytes_read_from_and_staged_to_the_engine_boundary()
{
    var engineBytes = new byte[] { 1, 2, 3 };
    byte[]? stagedBytes = null;
    var store = new VintageStoryWorldStateStore(
        key => key == VintageStoryWorldStateStore.EnvelopeKey ? engineBytes : null,
        (key, bytes) => stagedBytes = bytes);

    var read = store.ReadEnvelope()!.Value.ToArray();
    read[0] = 9;
    var input = new byte[] { 4, 5, 6 };
    store.StageEnvelope(input);
    input[0] = 9;

    Assert.Equal(1, engineBytes[0]);
    Assert.Equal(4, stagedBytes![0]);
}
\end{lstlisting}

Add capsule-store cases for exact player/key routing, multiple sorted
settlement IDs, stage/\allowbreak{}read/\allowbreak{}remove roundtrip,
clone ownership, duplicate IDs, structural bounds,
corrupt/unknown-version bytes, and the default
\passthrough{\lstinline!UnprovedCoSerialization!} capability.

\begin{lstlisting}
dotnet test tests/FirstGearBank.Adapter.Tests/FirstGearBank.Adapter.Tests.csproj -c Release
\end{lstlisting}

Expected: FAIL because the Vintage Story storage adapters do not exist.

\begin{itemize}
\tightlist
\item[$\square$]
  \textbf{Step 9: Write the Vintage Story storage adapters with exact
  namespaced keys}
\end{itemize}

\begin{lstlisting}
public sealed class VintageStoryWorldStateStore : IBankWorldStateStore
{
    public const string InstallMarkerKey = "firstgearbank:installed";
    public const string EnvelopeKey = "firstgearbank:world-state";

    public VintageStoryWorldStateStore(ISaveGame saveGame);
    public ReadOnlyMemory<byte>? ReadInstallMarker();
    public ReadOnlyMemory<byte>? ReadEnvelope();
    public void StageInstallMarker(ReadOnlyMemory<byte> bytes);
    public void StageEnvelope(ReadOnlyMemory<byte> bytes);
}
\end{lstlisting}

Use
\texttt{ISaveGame.}\allowbreak{}\texttt{GetData(}\allowbreak{}\texttt{string)}\allowbreak{}
and
\texttt{StoreData(}\allowbreak{}\texttt{string,}\allowbreak{}\texttt{\ }\allowbreak{}\texttt{byte[])}\allowbreak{}
only. Clone input/output buffers at the adapter boundary so the engine
and core never share mutable byte arrays.

\texttt{VintageStoryPlayerCapsuleStore} uses
\texttt{IServerPlayer.}\allowbreak{}\texttt{SetModData\textless{}}\allowbreak{}\texttt{byte[]\textgreater{}}\allowbreak{}
and \passthrough{\lstinline!GetModData<byte[]>!} under the key
\texttt{firstgearbank:}\allowbreak{}\texttt{physical-}\allowbreak{}\texttt{settlements},
but reports capability \passthrough{\lstinline!UnprovedCoSerialization!}
until the controlled crash matrix has repeatably proved a stronger
result. It never exposes this storage as automatic proof merely because
bytes exist.

Keep the public constructors API-backed. Add internal delegate-backed
constructors, exposed only to
\texttt{FirstGearBank.}\allowbreak{}\texttt{Adapter.}\allowbreak{}\texttt{Tests}
through an SDK \passthrough{\lstinline!InternalsVisibleTo!} item in the
adapter project file, so the storage contract tests can record exact
keys and buffer cloning without implementing the large Vintage Story
interfaces. Those delegates are call spies, not substitutes for engine
durability tests. Store a versioned, bounded, settlement-ID-sorted
capsule collection per player; malformed, duplicate, oversized, or
unknown-version bytes return \texttt{InvalidSettlementEvidence} and are
never overwritten during load.

\begin{itemize}
\tightlist
\item[$\square$]
  \textbf{Step 10: Replace the empty ModSystem with wiring only}
\end{itemize}

\begin{lstlisting}
public sealed class FirstGearBankModSystem : ModSystem
{
    private ICoreServerAPI? _serverApi;
    private FirstGearBankServerHost? _serverHost;

    public override void StartServerSide(ICoreServerAPI api);
    public override void Dispose();
}
\end{lstlisting}

\passthrough{\lstinline!StartServerSide!} validates assembly version
1.22.7.0, creates data/\allowbreak{}config/\allowbreak{}log paths with
verified public APIs, constructs the adapters/host, and subscribes host
methods to \passthrough{\lstinline!SaveGameCreated!},
\passthrough{\lstinline!SaveGameLoaded!}, and
\passthrough{\lstinline!GameWorldSave!}.
\passthrough{\lstinline!Dispose!} unsubscribes the exact delegates
before disposing the host/logger and is safe if startup stopped partway.
Do not initialize client authority or register packet DTOs in this
phase.

\begin{itemize}
\tightlist
\item[$\square$]
  \textbf{Step 11: Compile the direct engine boundary and run adapter
  tests}
\end{itemize}

\begin{lstlisting}
dotnet build "First Gear Bank/First Gear Bank.csproj" -c Release
dotnet test tests/FirstGearBank.Adapter.Tests/FirstGearBank.Adapter.Tests.csproj -c Release
\end{lstlisting}

Expected: zero errors; test output may contain the test-only API DLL,
while release package verification still excludes it.

\begin{itemize}
\tightlist
\item[$\square$]
  \textbf{Step 12: Package and perform the controlled 1.22.7 smoke run}
\end{itemize}

\begin{lstlisting}
dotnet run --project CakeBuild/CakeBuild.csproj -- --target=VerifyPackage
\end{lstlisting}

Deploy explicitly to the disposable StoryForge test installation, create
a backed-up disposable world, and execute the matrix in
\texttt{docs/}\allowbreak{}\texttt{testing/}\allowbreak{}\texttt{foundation-}\allowbreak{}\texttt{engine-}\allowbreak{}\texttt{smoke.}\allowbreak{}\texttt{md}.
At minimum verify: mod loads; exact game/API version is logged; fresh
state stages on the first world save; restart restores the same revision
and world ID; copied corruption quarantines; no client-side financial
service starts; a diagnostic capsule roundtrips through an ordinary
player save/restart without any inventory claim; and the file log
contains one bounded startup/restore result. Mark crash ordering and
every synthesized physical-settlement phase as deferred to Phase 3
rather than manufacturing evidence in a foundation-only mod.

Record the exact ZIP SHA-256, world backup, observed revisions, and
pass/fail in the smoke document. A normal capsule roundtrip proves
persistence only; it does not promote
\passthrough{\lstinline!UnprovedCoSerialization!} or claim inventory
ordering.

Regenerate the smoke document\textquotesingle s
\passthrough{\lstinline!.tex!} and \passthrough{\lstinline!.pdf!} from
Markdown and visually verify the rendered evidence before commit.

\begin{itemize}
\tightlist
\item[$\square$]
  \textbf{Step 13: Commit the adapter and observed smoke result}
\end{itemize}

\begin{lstlisting}
git diff --check
git add "First Gear Bank/FirstGearBankModSystem.cs" "First Gear Bank/Runtime" `
  "First Gear Bank/Persistence" "First Gear Bank/Configuration" "First Gear Bank/Settlement" `
  "First Gear Bank/Diagnostics" "First Gear Bank/First Gear Bank.csproj" `
  tests/FirstGearBank.Adapter.Tests README.md docs/testing/foundation-engine-smoke.md `
  docs/testing/foundation-engine-smoke.tex
git add -u "First Gear Bank/First Gear BankModSystem.cs"
git commit -m "feat: wire Vintage Story foundation lifecycle"
\end{lstlisting}

\subsubsection{Task 13: Prove the Foundation as One Recoverable Packaged
Vertical
Slice}\label{task-13-prove-the-foundation-as-one-recoverable-packaged-vertical-slice}

\textbf{Files:}

\begin{itemize}
\tightlist
\item
  Create:
  \texttt{tests/}\allowbreak{}\texttt{FirstGearBank.}\allowbreak{}\texttt{Core.}\allowbreak{}\texttt{Tests/}\allowbreak{}\texttt{Integration/}\allowbreak{}\texttt{FoundationRecoveryTests.}\allowbreak{}\texttt{cs}
\item
  Create:
  \texttt{tests/}\allowbreak{}\texttt{FirstGearBank.}\allowbreak{}\texttt{Core.}\allowbreak{}\texttt{Tests/}\allowbreak{}\texttt{Architecture/}\allowbreak{}\texttt{CoreDependencyTests.}\allowbreak{}\texttt{cs}
\item
  Create:
  \texttt{tests/}\allowbreak{}\texttt{FirstGearBank.}\allowbreak{}\texttt{Adapter.}\allowbreak{}\texttt{Tests/}\allowbreak{}\texttt{Integration/}\allowbreak{}\texttt{FoundationHostIntegrationTests.}\allowbreak{}\texttt{cs}
\item
  Create:
  \texttt{tests/}\allowbreak{}\texttt{CakeBuild.}\allowbreak{}\texttt{Tests/}\allowbreak{}\texttt{SourceStandardsTests.}\allowbreak{}\texttt{cs}
\item
  Create:
  \texttt{tests/}\allowbreak{}\texttt{FirstGearBank.}\allowbreak{}\texttt{Core.}\allowbreak{}\texttt{Tests/}\allowbreak{}\texttt{TestSupport/}\allowbreak{}\texttt{FoundationRecoveryRig.}\allowbreak{}\texttt{cs}
\item
  Create:
  \texttt{docs/}\allowbreak{}\texttt{testing/}\allowbreak{}\texttt{foundation-}\allowbreak{}\texttt{acceptance.}\allowbreak{}\texttt{md}
\item
  Create:
  \texttt{docs/}\allowbreak{}\texttt{testing/}\allowbreak{}\texttt{foundation-}\allowbreak{}\texttt{acceptance.}\allowbreak{}\texttt{tex}
  through regeneration
\item
  Generate but do not commit:
  \texttt{docs/}\allowbreak{}\texttt{testing/}\allowbreak{}\texttt{foundation-}\allowbreak{}\texttt{acceptance.}\allowbreak{}\texttt{pdf}
\item
  Modify: \passthrough{\lstinline!README.md:1!}
\item
  Modify:
  \texttt{docs/}\allowbreak{}\texttt{superpowers/}\allowbreak{}\texttt{plans/}\allowbreak{}\texttt{2026-}\allowbreak{}\texttt{09-}\allowbreak{}\texttt{08-}\allowbreak{}\texttt{first-}\allowbreak{}\texttt{gear-}\allowbreak{}\texttt{bank-}\allowbreak{}\texttt{v1-}\allowbreak{}\texttt{roadmap.}\allowbreak{}\texttt{md:}\allowbreak{}\texttt{1}
\item
  Modify:
  \texttt{docs/}\allowbreak{}\texttt{superpowers/}\allowbreak{}\texttt{plans/}\allowbreak{}\texttt{2026-}\allowbreak{}\texttt{09-}\allowbreak{}\texttt{08-}\allowbreak{}\texttt{first-}\allowbreak{}\texttt{gear-}\allowbreak{}\texttt{bank-}\allowbreak{}\texttt{v1-}\allowbreak{}\texttt{roadmap.}\allowbreak{}\texttt{tex:}\allowbreak{}\texttt{1}
  through regeneration
\item
  Generate but do not commit:
  \texttt{docs/}\allowbreak{}\texttt{superpowers/}\allowbreak{}\texttt{plans/}\allowbreak{}\texttt{2026-}\allowbreak{}\texttt{09-}\allowbreak{}\texttt{08-}\allowbreak{}\texttt{first-}\allowbreak{}\texttt{gear-}\allowbreak{}\texttt{bank-}\allowbreak{}\texttt{v1-}\allowbreak{}\texttt{roadmap.}\allowbreak{}\texttt{pdf}
\end{itemize}

\textbf{Interfaces:}

\begin{itemize}
\tightlist
\item
  Consumes: every Task 1-12 deliverable.
\item
  Produces: a green, packaged, source-linked Phase 1 baseline and
  recorded evidence for authoring the financial-core plan.
\end{itemize}

\begin{itemize}
\tightlist
\item[$\square$]
  \textbf{Step 1: Write the failing
  save/\allowbreak{}reload/\allowbreak{}replay integration test}
\end{itemize}

\begin{lstlisting}
[Fact]
public void Foundation_state_replays_in_process_but_rejects_the_scope_after_restart()
{
    var rig = FoundationRecoveryRig.CreateDeterministic();
    var scope = rig.OpenScope(PlayerIds.Ada);
    var request = rig.FinalizedDepositTestRequest(scope, sequence: 1, units: 1_000_000);

    var first = rig.Coordinator.Execute(request.Context, request.Mutation);
    var inProcessRetry = rig.Coordinator.Execute(request.Context, request.Mutation);

    Assert.True(first.IsSuccess);
    Assert.Equal(first.InvocationKey, inProcessRetry.InvocationKey);
    Assert.Equal(first.ResultingRevision, inProcessRetry.ResultingRevision);
    Assert.True(first.OperationIds.AsSpan().SequenceEqual(inProcessRetry.OperationIds.AsSpan()));

    var saved = rig.Persistence.EncodeCurrent();
    var restored = FoundationRecoveryRig.Restore(saved);
    var beforeStaleRetry = restored.State;
    var staleRetry = restored.Coordinator.Execute(request.Context, request.Mutation);

    Assert.False(staleRetry.IsSuccess);
    Assert.Equal(BankErrorCode.InactiveRequestScope, staleRetry.Failure?.Code);
    Assert.Equal(beforeStaleRetry.Revisions.Global, restored.State.Revisions.Global);
    Assert.Equal(0, restored.TimestampCaptureCount);
    Assert.Equal(0, restored.IdentifierIssueCount);
    var restoredFinance = restored.State.Finance.TryGetVerified().Value;
    Assert.Single(restoredFinance.Journal);
    Assert.Equal(
        PhysicalSettlementPhase.Finalized,
        Assert.Single(restored.State.Settlements.TryGetVerified().Value.Records).Phase);
    var replayed = restored.Projector.Replay(
        restoredFinance.Journal,
        ProjectionCheckpoint.Empty).Value;
    var projectionMatch = ProjectionValidator.MatchesReplay(restoredFinance.Projection, replayed);

    Assert.True(projectionMatch.IsSuccess);
    Assert.True(projectionMatch.Value);
}
\end{lstlisting}

This uses a matching finalized settlement fixture, not a claim that
playable inventory mutation exists. Add an admitted rejection retry
before staging, pruned response, fresh-scope success after restart,
restored automatic-key duplicate, marker repair, registry-section
corruption with finance available, journal corruption with finance
read-only, and exact raw-byte preservation.

\begin{itemize}
\tightlist
\item[$\square$]
  \textbf{Step 2: Run the integration tests and observe the first
  missing fixture/behavior failure}
\end{itemize}

\begin{lstlisting}
dotnet test tests/FirstGearBank.Core.Tests/FirstGearBank.Core.Tests.csproj -c Release
\end{lstlisting}

Expected: FAIL until the deterministic integration rig and any exposed
contract gaps are completed.

\begin{itemize}
\tightlist
\item[$\square$]
  \textbf{Step 3: Add only the deterministic fixture code needed by the
  integration tests}
\end{itemize}

Keep fake time, entropy, identifiers, and world storage under the test
project. Drive production coordinator, journal, projection, codecs,
migration registry, and loader without replacing them with mocks. Use
the same canonical request bytes before and after restore.

\begin{itemize}
\tightlist
\item[$\square$]
  \textbf{Step 4: Verify the full recovery test matrix passes}
\end{itemize}

\begin{lstlisting}
dotnet test tests/FirstGearBank.Core.Tests/FirstGearBank.Core.Tests.csproj -c Release
\end{lstlisting}

\begin{itemize}
\tightlist
\item[$\square$]
  \textbf{Step 5: Add architecture and source-standard regression tests}
\end{itemize}

\passthrough{\lstinline!CoreDependencyTests!} loads
\passthrough{\lstinline!FirstGearBank.Core.dll!} and rejects any
referenced assembly whose name starts with
\passthrough{\lstinline!Vintagestory!}, \passthrough{\lstinline!Cake!},
or the adapter assembly name.
\passthrough{\lstinline!SourceStandardsTests!} walks every
human-maintained \passthrough{\lstinline!.cs!} file in the repository
working tree, including files not yet committed, while excluding only
\passthrough{\lstinline!bin!}, \passthrough{\lstinline!obj!},
\passthrough{\lstinline!Artifacts!}, and generated code. It requires the
first token to be a
\texttt{/}\allowbreak{}\texttt{*\ }\allowbreak{}\texttt{.}\allowbreak{}\texttt{.}\allowbreak{}\texttt{.}\allowbreak{}\texttt{\ }\allowbreak{}\texttt{*/}\allowbreak{}
responsibility block and rejects documentation lines longer than 120
characters. Do not obtain the file list from
\passthrough{\lstinline!git ls-files!}, because that would skip the new
files at the very gate intended to check them. During the same review,
inspect every callable changed in Phase 1 for its immediately preceding
\texttt{/}\allowbreak{}\texttt{/}\allowbreak{}\texttt{/}\allowbreak{}\texttt{/}\allowbreak{}
block and exact three-blank-line spacing; do not use a brittle
declaration regex as false assurance.

\begin{lstlisting}
[Fact]
public void Core_has_no_engine_or_build_tool_dependency()
{
    var forbidden = typeof(CoreAssemblyMarker).Assembly
        .GetReferencedAssemblies()
        .Where(reference =>
            reference.Name?.StartsWith("Vintagestory", StringComparison.OrdinalIgnoreCase) == true ||
            reference.Name?.StartsWith("Cake", StringComparison.OrdinalIgnoreCase) == true ||
            reference.Name == "FirstGearBank")
        .ToArray();

    Assert.Empty(forbidden);
}
\end{lstlisting}

\begin{lstlisting}
dotnet test tests/CakeBuild.Tests/CakeBuild.Tests.csproj -c Release
dotnet test tests/FirstGearBank.Core.Tests/FirstGearBank.Core.Tests.csproj -c Release
\end{lstlisting}

Expected: FAIL on any remaining source-layout or dependency-boundary
violation; correct those violations before the final verification step.

\begin{itemize}
\tightlist
\item[$\square$]
  \textbf{Step 6: Run formatting, source, dependency, vulnerability, and
  full test checks}
\end{itemize}

\begin{lstlisting}
dotnet format "First Gear Bank.sln" --verify-no-changes --no-restore
dotnet test "First Gear Bank.sln" -c Release
dotnet list "First Gear Bank.sln" package --vulnerable --include-transitive
git diff --check
\end{lstlisting}

Expected: formatter reports no change; all tests pass; NuGet reports no
known vulnerable package; whitespace check passes. If the formatter
conflicts with the required three-blank-line standard, preserve the
mandated source layout and record the specific formatter limitation
instead of rewriting documented callables incorrectly.

\begin{itemize}
\tightlist
\item[$\square$]
  \textbf{Step 7: Build and inspect the final foundation archive twice}
\end{itemize}

\begin{lstlisting}
dotnet run --project CakeBuild/CakeBuild.csproj -- --target=VerifyPackage
$archive = Get-Item -LiteralPath "Releases/firstgearbank_1.0.0.zip"
$firstHash = (Get-FileHash -LiteralPath $archive.FullName -Algorithm SHA256).Hash
dotnet run --project CakeBuild/CakeBuild.csproj -- --target=VerifyPackage
$secondHash = (Get-FileHash -LiteralPath $archive.FullName -Algorithm SHA256).Hash
if ($firstHash -ne $secondHash) { throw "Verified package bytes changed between identical builds." }
\end{lstlisting}

Expected package root: \passthrough{\lstinline!modinfo.json!},
\passthrough{\lstinline!FirstGearBank.dll!},
\passthrough{\lstinline!FirstGearBank.Core.dll!}, and the configuration
template/resource or assets intentionally emitted by the adapter.
Expected absence: \passthrough{\lstinline!VintagestoryAPI.dll!},
\passthrough{\lstinline!VSSurvivalMod.dll!}, Cake DLLs, test assemblies,
PDBs, \passthrough{\lstinline!.deps.json!}, and
\passthrough{\lstinline!.runtimeconfig.json!}.

\begin{itemize}
\tightlist
\item[$\square$]
  \textbf{Step 8: Complete the acceptance record from command and engine
  evidence}
\end{itemize}

\texttt{docs/}\allowbreak{}\texttt{testing/}\allowbreak{}\texttt{foundation-}\allowbreak{}\texttt{acceptance.}\allowbreak{}\texttt{md}
must record exact command lines, exit codes, test counts, ZIP
hash/layout, installed API assembly version,
fresh/\allowbreak{}save/\allowbreak{}restart smoke results, quarantine
smoke results, log path/labels, unresolved engine gates, and the
accrued-liability-index gate result. Do not mark a row passing from
intended behavior or mock output alone.

Update \passthrough{\lstinline!README.md!} to state that the foundation
is installed but player-facing banking remains unavailable until the
later roadmap phases. Mark Phase 1 complete in the roadmap only when
every automated check and required 1.22.7 smoke row has observed
evidence. Inventory/capsule ordering and per-phase process-crash rows
are explicitly deferred Phase 3 evidence, not required Phase 1 rows.
Regenerate the acceptance record and roadmap
\passthrough{\lstinline!.tex!}/\passthrough{\lstinline!.pdf!} files from
Markdown and visually verify every page; stage only Markdown and LaTeX.

\begin{itemize}
\tightlist
\item[$\square$]
  \textbf{Step 9: Review the diff against the locked phase scope}
\end{itemize}

\begin{lstlisting}
git status --short
git diff --stat
git diff --check
rg -n "Vintagestory|ICore|Entity|BlockPos|ItemStack" FirstGearBank.Core
\end{lstlisting}

Expected: the final search returns only architecture-boundary test
strings or explanatory exclusions, never a core engine type/reference.
Confirm that no loan, overdraft, margin, repo, swap, stock, fund,
option, commodity, joint account, cross-world migration, currency
conversion, temporal CD, CD redemption, or transfer fee behavior entered
the diff.

\begin{itemize}
\tightlist
\item[$\square$]
  \textbf{Step 10: Commit the verified foundation gate}
\end{itemize}

\begin{lstlisting}
git add tests/FirstGearBank.Core.Tests/Integration tests/FirstGearBank.Core.Tests/Architecture `
  tests/FirstGearBank.Adapter.Tests/Integration tests/CakeBuild.Tests/SourceStandardsTests.cs `
  tests/FirstGearBank.Core.Tests/TestSupport/FoundationRecoveryRig.cs `
  docs/testing/foundation-acceptance.md docs/testing/foundation-acceptance.tex README.md `
  docs/superpowers/plans/2026-09-08-first-gear-bank-v1-roadmap.md `
  docs/superpowers/plans/2026-09-08-first-gear-bank-v1-roadmap.tex
git commit -m "test: verify the recoverable banking foundation"
git status --short
\end{lstlisting}

Expected: clean status except for any explicitly documented user-owned
pre-existing changes.

\subsection{Plan Self-Review}\label{plan-self-review}

\subsubsection{Foundation specification
coverage}\label{foundation-specification-coverage}

\begin{itemize}
\tightlist
\item
  Sections 1-3: Tasks 1-4 establish exact target/mod identity, pure-core
  boundary, server authority, and safe packaging.
\item
  Section 4: Task 5 implements unit scale, ties-to-even conversion,
  display independence, checked arithmetic, and cap.
\item
  Section 5 foundation: Tasks 5 and 8 implement hidden world/player
  identities and the exclusive lazy-opening rule.
\item
  Section 6: Tasks 7-9 implement immutable double entry, typed facts,
  replay, revisions, and request plus automatic-operation idempotency.
\item
  Section 18 foundation: Task 6 implements every static
  field/\allowbreak{}default/\allowbreak{}domain and JSONC ownership;
  derived activation is explicitly gated on the financial formulas that
  consume it.
\item
  Sections 20.2-20.4 foundation: Tasks 9-12 implement snapshots, staging
  semantics, integrity, migrations, replay, request and automatic-key
  persistence, explicit unavailable-domain slots, opaque
  quarantined-section pass-through, recovery evidence archival, and
  versioned hash domains.
\item
  Section 21: every task carries failure, privacy, trust,
  bounded-allocation, logging, and no-scan constraints.
\item
  Section 22: Tasks 1, 10, and 12 record verified APIs and isolate
  runtime-only evidence.
\item
  Section 23 foundation rows: Tasks 3-13 exercise package safety, money,
  journal, replay, request, schema, corruption, settlement evidence,
  adapter lifecycle, and multiplayer-independent recovery behavior.
\item
  Sections 7-17 and remaining 18-26 behavior are allocated to named
  executable plans in the roadmap and are not claimed by this foundation
  milestone.
\end{itemize}

\subsubsection{Type and boundary
consistency}\label{type-and-boundary-consistency}

\begin{itemize}
\tightlist
\item
  \passthrough{\lstinline!PlayerId!}, both invocation-key variants,
  \passthrough{\lstinline!CommandPayloadDigest!},
  \passthrough{\lstinline!FinancialTimestamp!},
  \passthrough{\lstinline!BankUnits!}, and the revision vector originate
  in Task 5 and are reused without aliases.
\item
  \passthrough{\lstinline!JournalRecord!} is the only monetary source
  accepted by projection and persistence.
\item
  \passthrough{\lstinline!BankStateRepository!} is the only
  published-state owner; the coordinator is the only mutation writer.
\item
  Mutations see only a domain-gated read view; unavailable authority is
  never represented by a fabricated empty value.
\item
  \passthrough{\lstinline!IBankWorldStateStore!},
  \texttt{IPlayerRecoveryCapsuleStore},
  \texttt{IFinancialTimestampSource}, and
  \passthrough{\lstinline!IBankLog!} point from core toward adapter
  implementations, never from core toward Vintage Story.
\item
  Persistence separates journal authority from projection caches and
  separates registry, delivery, settlement, branch, configuration,
  clock, and market failure domains.
\item
  A known quarantined section remains opaque and byte-preserved across
  permitted healthy saves; an unknown section kind blocks whole-envelope
  mutation because container v1 has no authenticated criticality flag.
\item
  \passthrough{\lstinline!GameWorldSave!} is a pre-write staging
  boundary. No code or log calls its return a durable disk
  acknowledgment.
\item
  The complete authenticated physical-recovery plan is reconstructive,
  while its cross-inventory co-serialization remains an engine-evidence
  gate; the accrued-liability result remains a separate mathematical
  evidence gate. Neither is represented as stronger than the API or
  proof supports.
\end{itemize}

\end{document}
